% Copyright (c) 2026 Scott Armstrong.
% Released under Apache 2.0 license.
\documentclass[11pt,letterpaper]{article}
\usepackage[T1]{fontenc}
\usepackage{amsmath,amsthm,amssymb,mathtools}
\usepackage[margin=1.05in]{geometry}
\usepackage{microtype}
\setlength{\emergencystretch}{1em}
\usepackage{xcolor}
\definecolor{paperlink}{RGB}{0,0,114}
\usepackage{hyperref}
\hypersetup{colorlinks=true,linktoc=all,
  linkcolor=paperlink,citecolor=paperlink,urlcolor=paperlink,pdfstartview=Fit,
  pdftitle={The Escauriaza--Seregin--Sverak theorem: version for the Lean formalization},
  pdfauthor={}}
\usepackage{enumitem}
\setlist[itemize]{leftmargin=2em,itemsep=0.25em,topsep=0.35em}
\setlist[enumerate]{leftmargin=2.25em,itemsep=0.25em,topsep=0.35em}
\setlength{\parindent}{1.25em}
\setlength{\parskip}{0pt}
\usepackage{fancyhdr}
\pagestyle{fancy}
\fancyhf{}
\fancyhead[L]{\small The Escauriaza--Seregin--\v{S}ver\'ak theorem}
\fancyhead[R]{\small\thepage}
\setlength{\headheight}{14pt}
\numberwithin{equation}{section}
\setcounter{tocdepth}{1}
\date{}

\theoremstyle{plain}
\newtheorem{theorem}{Theorem}[section]
\newtheorem{proposition}[theorem]{Proposition}
\newtheorem{lemma}[theorem]{Lemma}
\newtheorem{corollary}[theorem]{Corollary}
\theoremstyle{definition}
\newtheorem{definition}[theorem]{Definition}
\newtheorem{deviation}[theorem]{Deviation}
\theoremstyle{remark}
\newtheorem{remark}[theorem]{Remark}

% Copyright (c) 2026 Scott Armstrong.
% Released under Apache 2.0 license.
% Shared notation for paper/ess.tex. Add a macro here only when it is used
% at least three times; never redefine a macro with a different meaning.

% Number systems
\newcommand{\R}{\mathbb{R}}
\newcommand{\N}{\mathbb{N}}
\newcommand{\Z}{\mathbb{Z}}

% Norms and absolute values (Euclidean on vectors and matrices)
\newcommand{\abs}[1]{\left\lvert #1\right\rvert}
\newcommand{\norm}[1]{\left\lVert #1\right\rVert}
\newcommand{\nf}[2]{#1/#2}

% Differential operators
\DeclareMathOperator{\dv}{div}
\DeclareMathOperator{\curl}{curl}
\DeclareMathOperator{\supp}{supp}
\DeclareMathOperator*{\esssup}{ess\,sup}
\newcommand{\dt}{\partial_t}
\newcommand{\Lop}{L}                       % backward heat operator dt + Delta (Sections on Carleman, UC, BU)

% Sets
\newcommand{\Ball}[2]{B_{#1}(#2)}          % \Ball{r}{x}: open Euclidean ball of radius r centred at x
\newcommand{\Hplus}{\R^3_+}                % {x_3 > 0}
\newcommand{\Qp}{Q_+}                      % \Hplus \times (0,1)
\newcommand{\Cyl}[2]{Q(#1,#2)}             % \Cyl{R}{T} = B_R(0) \times (0,T)
\newcommand{\Pcyl}[2]{Q_{#1}(#2)}          % \Pcyl{r}{z}: past parabolic cylinder B_r(x) \times (t-r^2, t] of CKN

% Function classes
\newcommand{\Cc}{C_c^\infty}
\newcommand{\Ctwoone}{C^{2,1}}
\newcommand{\Lloc}[1]{L^{#1}_{\mathrm{loc}}}

% Navier--Stokes objects
\newcommand{\Leray}{\mathbb{P}}           % Leray projection
\newcommand{\Riesz}[1]{\mathcal{R}_{#1}}   % Riesz transform R_j
\newcommand{\J}{J}                         % the space J of ESS (divergence-free L^2 fields)
\newcommand{\Jeps}{J_\varepsilon}          % spatial mollification at scale epsilon


\title{The Escauriaza--Seregin--\v{S}ver\'ak theorem:\\ version for the Lean formalization}
\author{}

\begin{document}
\maketitle

\begin{abstract}
We state a suitable form of Leray's global existence theorem for the three-dimensional Navier--Stokes equations on $\R^3$, which is proved in the companion Caffarelli--Kohn--Nirenberg paper, and we prove the endpoint regularity theorem of Escauriaza, Seregin, and \v{S}ver\'ak. Every divergence-free initial datum in $L^2(\R^3)$ has a global Leray--Hopf solution that is suitable and has a canonical pressure in $L^{5/3}(\R^3\times(0,T))$ on each finite time interval. The formalized Caffarelli--Kohn--Nirenberg theorem then gives zero one-dimensional parabolic Hausdorff measure for its singular set. Every Leray--Hopf solution with finite $L^\infty(0,T;L^3(\R^3))$ norm is regular at every point of $\R^3\times(0,T)$, belongs to $L^5(\R^3\times(0,T))$, and is unique among Leray--Hopf solutions with the same datum. The endpoint argument combines an open-top smallness criterion, vorticity regularity, and explicit Carleman estimates for unique continuation and backward uniqueness. We also prove the Ladyzhenskaya--Prodi--Serrin theorem for the full range $3<s\le\infty$, smoothness under the critical $L^\infty_tL^3_x$ bound, and the resulting regularity criterion for $3\le s\le\infty$. These results, together with the unique continuation, backward uniqueness and Carleman theorems, are formally verified in Lean~4.
\end{abstract}

{\small
\tableofcontents
}
\clearpage

%======================================================================
\section{Introduction}\label{sec:intro}
%======================================================================

The three-dimensional Navier--Stokes equations pose two basic questions for finite-energy data: whether a global weak solution exists, and whether such a solution can develop singularities. Leray's construction gives global weak solutions on $\R^3$ for divergence-free initial velocities in $L^2(\R^3)$~\cite{Leray1934}. The Caffarelli--Kohn--Nirenberg theorem gives a geometric answer to the second question for \emph{suitable} weak solutions: their singular set has zero one-dimensional parabolic Hausdorff measure~\cite{CKN1982}. The corresponding theorem is available in the formalized CKN library~\cite{CKNLean}. Thus suitability is the condition that lets Leray's existence result carry the CKN partial-regularity conclusion.

Escauriaza, Seregin, and \v{S}ver\'ak proved a sharper endpoint statement: a Leray--Hopf solution whose velocity belongs to $L^\infty(0,T;L^3(\R^3))$ has no singular points in the open time interval~\cite[Theorems~1.3--1.4]{ESS2003}. They also proved that this bound forces the velocity into space-time $L^5$ and gives uniqueness among Leray--Hopf solutions with the same datum. This norm is invariant under Navier--Stokes scaling, so the result settles regularity at a critical integrability threshold. The proof here gives the endpoint $L^5$ and uniqueness conclusions by constructing a small-time $L^5$ solution and proving a weak--strong comparison; it also fixes the associated pressure and makes the endpoint regularity steps explicit.

\paragraph{Main results.}
Section~\ref{sec:leray-hopf} fixes the solution notions and contains the precise definitions used below. Statements (A)--(E) and (G) repeat results stated in Section~\ref{sec:leray-hopf}; (A)--(C) and (G) are proved in the Caffarelli--Kohn--Nirenberg paper~\cite{CKNLean} (Section~\ref{sec:existence-in-ckn}); statement (F) is Theorem~\ref{thm:ess-l5-unique}, stated and proved in Part~V. Statements (H)--(J) are Theorem~\ref{thm:lps} and Corollaries~\ref{cor:ess-smooth} and~\ref{cor:serrin-criterion}, stated and proved in Part~VI.

\begin{enumerate}[label=\textbf{(\Alph*)},leftmargin=2.8em]
\item \textbf{Leray existence.} For every $a\in\J$ there are $u$, $Du$ and $p$ such that $(u,Du)$ is a global Leray--Hopf solution with datum $a$, and $(u,Du,p,0)$ is a suitable weak solution on $\R^3\times(0,\infty)$ for every $q>\nf52$ (Theorem~\ref{thm:leray}).

\item \textbf{The CKN headline.} For every $a\in\J$ there is a global Leray--Hopf solution $(u,Du)$ with datum $a$ whose singular set in $\R^3\times(0,\infty)$ has zero one-dimensional parabolic Hausdorff measure (Corollary~\ref{cor:ckn-headline}).

\item \textbf{Associated pressure.} Let $(u,Du)$ be a Leray--Hopf solution on $\R^3\times[0,T]$ with datum $a\in\J$. There is $p\in L^{5/3}(\R^3\times(0,T))$ such that (S3) of Definition~\ref{def:sws} holds on $\R^3\times(0,T)$ with $f=0$. If in addition $\esssup_{t\in(0,T)}\norm{u(\cdot,t)}_{L^3(\R^3)}<\infty$, then this $p$ also belongs to $L^\infty((0,T);L^{3/2}(\R^3))$ (Theorem~\ref{thm:assoc-pressure}).

\item \textbf{Local endpoint regularity.} Let $u$, $Du$, $p$ be measurable on $Q=B_1\times(-1,0)$ with
\begin{gather*}
\esssup_{t\in(-1,0)}\int_{B_1}\abs{u(x,t)}^2\,dx<\infty,\qquad
\iint_Q\abs{u}^2+\abs{Du}^2<\infty,\qquad
p\in L^{3/2}(Q),\\
\esssup_{t\in(-1,0)}\int_{B_1}\abs{u(x,t)}^3\,dx<\infty,
\end{gather*}
such that $Du(\cdot,t)$ is the weak gradient of $u(\cdot,t)$ on $B_1$ for a.e. $t\in(-1,0)$, and (S2) and (S3) of Definition~\ref{def:sws} hold on $Q$ with $f=0$. Then there is $\gamma\in(0,1]$ and a function $\tilde u$ equal to $u$ almost everywhere on $B_{1/2}\times(-\nf14,0)$ that is parabolically H\"older continuous with exponent $\gamma$ on the closure of $B_{1/2}\times(-\nf14,0)$ (Theorem~\ref{thm:ess-local}).

\item \textbf{Global endpoint regularity.} Let $(u,Du)$ be a Leray--Hopf solution on $\R^3\times[0,T]$ with datum $a\in\J$, and suppose $\esssup_{t\in(0,T)}\int_{\R^3}\abs{u(x,t)}^3\,dx<\infty$. Then every point of $\R^3\times(0,T)$ is a regular point of $u$ (Theorem~\ref{thm:ess-global}).

\item \textbf{$L^5$ and uniqueness.} Under the same hypotheses, $u\in L^5(\R^3\times(0,T))$, and every other Leray--Hopf solution $(v,Dv)$ on $[0,T]$ with datum $a$ satisfies $u=v$ almost everywhere on $\R^3\times(0,T)$ (Theorem~\ref{thm:ess-l5-unique}).

\item \textbf{Forced Leray existence.} For data and forces in the class of
Theorem~\ref{thm:leray-forced}, there is one global forced Leray--Hopf
solution that is suitable for every $q>5/2$ for which $f$ is locally in
$L^q$; its singular set has zero one-dimensional parabolic Hausdorff
measure, and its associated pressure has the local split stated in
Theorems~\ref{thm:leray-forced-singularSet} and
\ref{thm:assoc-pressure-forced}.
\item \textbf{Ladyzhenskaya--Prodi--Serrin.} Every Leray--Hopf solution in $L^\ell(0,T;L^s(\R^3))$ with $3/s+2/\ell=1$ and $3<s\le\infty$ is unique among Leray--Hopf solutions with its datum and has a representative smooth on $\R^3\times(0,T]$, one-sided at $T$ (Theorem~\ref{thm:lps}).

\item \textbf{Smooth critical endpoint.} Every Leray--Hopf solution in $L^\infty(0,T;L^3(\R^3))$ has the same smoothness and uniqueness, as well as a finite space-time $L^5$ norm (Corollary~\ref{cor:ess-smooth}).

\item \textbf{Regularity criterion.} Combining (H) and (I): every Leray--Hopf solution in $L^\infty(0,T;L^3(\R^3))$, in $L^\ell(0,T;L^s(\R^3))$ with $3<s<\infty$ and $\ell=2s/(s-3)$, or in $L^2(0,T;L^\infty(\R^3))$ is unique among Leray--Hopf solutions with its datum and has a representative smooth on $\R^3\times(0,T]$, one-sided at $T$ (Corollary~\ref{cor:serrin-criterion}).
\end{enumerate}

\paragraph{Smoothness.}
Part~VI proves the Ladyzhenskaya--Prodi--Serrin theorem in the full range $3<s\le\infty$, including smoothness through the top time. At $(s,\ell)=(5,5)$ it turns the $L^5$ conclusion of Theorem~\ref{thm:ess-l5-unique} into Corollary~\ref{cor:ess-smooth}. Theorem~\ref{thm:ess-l5-unique} itself states only integrability and uniqueness; smoothness is stated separately, in Theorem~\ref{thm:lps} and Corollaries~\ref{cor:ess-smooth} and~\ref{cor:serrin-criterion}.

\paragraph{Part I: linear continuation.}
The ESS argument ultimately applies backward uniqueness to vorticity, but its linear inputs must match the regularity actually obtained from a suitable weak solution. We therefore formulate the linear fields in the space-time weak-derivative class $W^{2,1}_2$ and require continuity only where a trace is evaluated. Propositions~\ref{prop:carleman-gauss} and~\ref{prop:carleman-halfspace} first establish the Gaussian and anisotropic Carleman estimates for smooth compactly supported fields, then Lemma~\ref{lem:carleman-sobolev} extends them by density to the fields used in the continuation arguments. The first estimate propagates infinite-order vanishing across a spatial boundary through an $L^2$ Gaussian bound on parabolic boxes; this replaces the pointwise local regularity estimate used in ESS Lemma~4.3. The half-space estimate gives decay away from the boundary plane. Its initial-time cutoff is controlled by a trace inequality, and its noncompact spatial cutoff is justified by collar estimates, so no zero extension across the initial slice is needed. The final large-growth case of backward uniqueness is handled on finitely many rescaled time slabs. The constants are explicit: the Gaussian estimate uses $c_0=e^{4/3}(9+2\sqrt6)$, and the half-space estimate uses $a_0(\alpha)=2$ and $c_*(\alpha)=5+7/(2(2\alpha-1))$. These details make the linear results applicable in the weak vorticity class without adding smoothness assumptions.

\paragraph{Parts II and III: existence and the associated pressure.}
The Leray--Hopf weak equation is initially tested only against divergence-free vector fields, whereas suitability requires the momentum identity for arbitrary compactly supported vector tests. The proofs that every $a\in\J$ has a global Leray--Hopf solution that is suitable, that every Leray--Hopf solution has an associated pressure defined by the whole-space double Riesz transform of $u\otimes u$, and their versions with a force, do not use Parts~I and IV--VI. They were moved to the last part of the Caffarelli--Kohn--Nirenberg paper~\cite{CKNLean}, where the statements of Section~\ref{sec:leray-hopf} keep their labels; Section~\ref{sec:existence-in-ckn} records this and the whole-space pressure used in Parts~IV--VI. The construction there uses spatially regularized solutions given by an $L^2$ Fourier mild formulation, in place of the Oseen-kernel construction of O{\.z}a{\'n}ski--Pooley~\cite[Theorem~4.2]{OzanskiPooley2018}, and passes to the limit with global tail control by a compactness argument that does not rely on a Helly selection theorem, Rellich compactness, or Aubin--Lions. Since the regularized pairs do not solve the unregularized equations, the limit is taken directly in their momentum and local-energy identities. The Caffarelli--Kohn--Nirenberg theorem then gives the stated measure bound for the singular set.

\paragraph{Part IV: the $L^\infty L^3$ endpoint.}
The local endpoint theorem assumes the weak momentum equation and a uniform $L^3$ bound but does not assume the local energy inequality. A local $L^4$ estimate and space-time mollification recover the local energy equality; a momentum-pairing modulus supplies a weakly continuous $L^3$ representative up to the top time, and a fixed whole-space pressure split controls the harmonic remainder under rescaling. We then measure smallness on open past cylinders by
\[
r^{-2}\int_{t_0-r^2}^{t_0}\int_{B_r(x_0)}\bigl(\abs{u}^3+\abs{p}^{3/2}\bigr)\,dx\,dt\,.
\]
If this scale-invariant quantity is below the CKN threshold, Theorem~A gives interior H\"older regularity. To reach the terminal face, shifted centers and overlapping representatives give a one-sided H\"older extension; the top point is not treated as an interior CKN regular point. The good-point lemma makes this criterion open and yields a lower bound at every bad point. Thus failure of H\"older regularity on the target compact cylinder produces the quantitative lower bound needed for a blow-up contradiction.

After rescaling a bad point, compactness gives an ancient suitable solution with zero terminal velocity trace and a pressure equal to the whole-space Riesz pressure of its velocity tensor. Far from the origin, the $L^3$ and pressure tails are small; CKN epsilon regularity makes the velocity bounded there. We then bootstrap the weak vorticity equation only as far as needed: a localized vorticity energy estimate, local heat gains, product bounds, and div--curl estimates give $W^{2,1}_2$ control and continuity of vorticity. This finite weak-to-regular argument replaces the all-orders Stokes/maximal-regularity bounds and local strong-solvability input used in the exterior part of ESS's proof~\cite[Theorem~1.4, §§3--5]{ESS2003}. Half-space backward uniqueness forces the vorticity to vanish outside a ball. At almost every time the CKN singular-set theorem supplies a time interval of regular points; unique continuation spreads the vanishing across the ball, and an $L^3$ Liouville argument makes the entire blow-up velocity zero. That contradicts the bad-point lower bound. For the global theorem, parabolic rescaling applies the local result at a later center, placing each interior point strictly inside the resulting H\"older cylinder.

\paragraph{Part V: $L^5$ and uniqueness.}
The critical bound gives only an almost-every-time $L^3$ estimate, so the first step is to identify the initial datum in $L^3$ using good times and the strong $L^2$ trace. A heat-flow approximation then makes the linear solution small in $L^5\cap L^4$ on a short interval. The forced Stokes estimates control the quadratic term in those two norms, producing a local Leray--Hopf solution; the $(5,5)$ energy comparison identifies it with the given weak solution. Away from zero, the pressure tail and the local regularity estimate give a uniform velocity bound over the whole space, and the multiplicative inequality converts that bound and the finite dissipation into $L^5$. The final comparison applies to every Leray--Hopf solution with the same datum.

\paragraph{Lean formalization.}
The main results of this manuscript, together with the linear continuation results used in their proofs, are formally verified in the proof assistant Lean~4 with its mathematical library Mathlib, building on the Lean formalization of the Caffarelli--Kohn--Nirenberg theorem~\cite{CKNLean}. The formalization accompanies this manuscript~\cite{ESSLean}. Ten of the results in the table below are stated there as Lean theorems, each in its own file of the directory \texttt{ESS/Statements}; the theorem \texttt{ESS.essGlobal}, for example, is stated in \texttt{ESS/Statements/EssGlobal.lean}, and the other files are named in the same way. The six results in the first part of the table, on Leray existence and the associated pressure, are proved in the Caffarelli--Kohn--Nirenberg paper~\cite{CKNLean} (Section~\ref{sec:existence-in-ckn}) and are formalized in the CKN library, each as a Lean theorem in its own file of the directory \texttt{CKN/Statements}; the ESS library imports them.
\begin{center}
\small
\begin{tabular}{lll}
Result & Content & Lean theorem\\
\hline
Theorem~\ref{thm:leray} & Leray existence & \texttt{CKN.leray\_existence}\\
Corollary~\ref{cor:ckn-headline} & Leray existence, singular set & \texttt{CKN.leray\_existence\_singularSet}\\
Theorem~\ref{thm:assoc-pressure} & associated pressure & \texttt{CKN.associatedPressure}\\
Theorem~\ref{thm:leray-forced} & forced Leray existence & \texttt{CKN.lerayExistenceForced}\\
Theorem~\ref{thm:leray-forced-singularSet} & forced existence, singular set & \texttt{CKN.lerayExistenceForcedSingularSet}\\
Theorem~\ref{thm:assoc-pressure-forced} & forced associated pressure & \texttt{CKN.associatedPressureForced}\\
\hline
Theorem~\ref{thm:ess-local} & local endpoint regularity & \texttt{ESS.essLocal}\\
Theorem~\ref{thm:ess-global} & global endpoint regularity & \texttt{ESS.essGlobal}\\
Proposition~\ref{prop:carleman-gauss} & Gaussian Carleman inequality & \texttt{ESS.carlemanGaussian}\\
Proposition~\ref{prop:carleman-halfspace} & half-space Carleman inequality & \texttt{ESS.carlemanHalfSpace}\\
Theorem~\ref{thm:uc} & unique continuation & \texttt{ESS.uniqueContinuation}\\
Theorem~\ref{thm:bu} & backward uniqueness & \texttt{ESS.backwardUniqueness}\\
Theorem~\ref{thm:ess-l5-unique} & $L^5$ and uniqueness & \texttt{ESS.essL5Unique}\\
Theorem~\ref{thm:lps} & mixed-norm smoothness and uniqueness & \texttt{ESS.ladyzhenskayaProdiSerrin}\\
Corollary~\ref{cor:ess-smooth} & smooth critical endpoint & \texttt{ESS.essSmooth}\\
Corollary~\ref{cor:serrin-criterion} & criterion for $3\le s\le\infty$ & \texttt{ESS.serrinCriterion}\\
\end{tabular}
\end{center}
The definitions of $\J$, finite-time and global Leray--Hopf solutions and their forced versions, the force classes, the forced quadratic tensor, and the space-time weak-derivative class are stated in \texttt{CKN/Statements}. The ten ESS theorems are stated in \texttt{ESS/Statements}; they use these CKN definitions, and suitable weak solutions and regular points are also those of~\cite{CKNLean}. The six CKN theorems and the ten ESS theorems are proved in full, without \texttt{sorry}, and their proofs use only the standard axioms \texttt{propext}, \texttt{Classical.choice} and \texttt{Quot.sound}. The formal proofs follow the arguments of this manuscript and of the CKN paper, although they are not a line-by-line transcription of them. As an independent check of the statements, the directory \texttt{comparators} of the ESS library contains two pairs of files, \texttt{Linear} and \texttt{Regularity}. In each pair, the file \texttt{Challenge.lean} imports only Mathlib and restates its theorems with every definition written directly in terms of Mathlib, and the file \texttt{Solution.lean} derives each restatement from the corresponding theorem of the formalization. Together they restate all ten ESS theorems. The comparator pair for the Leray results is part of the CKN library; it restates five of the six theorems above, and Theorem~\ref{thm:assoc-pressure-forced} is not restated, because its pressure constructions are too long for a self-contained restatement.

\paragraph{Outline.}
Section~\ref{sec:notation} sets the notation, and Section~\ref{sec:leray-hopf} gives the solution classes and main statements. Part~I (Sections~\ref{sec:carleman}--\ref{sec:bu}) proves the linear continuation results; Parts~II and III (Section~\ref{sec:existence-in-ckn}) record the existence and associated-pressure results, which are proved in the CKN paper~\cite{CKNLean}; Part~IV (Sections~\ref{sec:local-regularity}--\ref{sec:thm-global}) proves the regularity endpoint; Part~V (Sections~\ref{sec:pv-trace}--\ref{sec:pv-assembly}) proves the $L^5$ and uniqueness conclusions; and Part~VI (Sections~\ref{sec:lps-statements}--\ref{sec:lps-smooth}) proves the mixed-norm theorem, the smooth endpoint corollary and the regularity criterion for $3\le s\le\infty$. The appendix records corrections and proof details in the comparison sources.

%======================================================================
\section{Notation and conventions}\label{sec:notation}
%======================================================================

\paragraph{Space, time and norms.}
Space is $\R^3$ and time is $\R$; a space-time point is $z=(x,t)$. On
vectors and matrices, $\abs{\cdot}$ is the Euclidean (Frobenius) norm. For a
field $w\colon\R^3\times\R\to\R^3$ with components $w_i$ we write
$\partial_j w_i$ for spatial and $\dt w_i$ for time derivatives,
$\abs{\nabla w}^2=\sum_{i,j}(\partial_j w_i)^2$ and
$\abs{\nabla^2 w}^2=\sum_{i,j,k}(\partial_j\partial_k w_i)^2$.
The ball $\Ball{r}{x}$ is open; $B_r=\Ball{r}{0}$.

\paragraph{Open past cylinders.}
For $z_0=(x_0,t_0)$ and $r>0$, write
\[
\begin{aligned}
 Q_r^-(z_0) &= B_r(x_0)\times(t_0-r^2,t_0),\\
 \overline{Q_r^-(z_0)} &= \overline{B_r(x_0)}\times[t_0-r^2,t_0].
\end{aligned}
\]

\paragraph{Space-time weak derivatives.}
Let $U\subset\R^3\times\R$ be open and $w\in\Lloc{1}(U;\R^3)$. Fields
$Dw=(Dw_{ij})$, $D^2w=(D^2w_{ijk})$ and $D_tw=(D_tw_i)$ in $\Lloc{1}(U)$ are
the \emph{space-time weak derivatives} $\partial_j w_i$,
$\partial_j\partial_k w_i$ and $\dt w_i$ of $w$ on $U$ if, for every
$\varphi\in\Cc(U)$ and all indices,
\[
\int_U w_i\,\partial_j\varphi=-\int_U Dw_{ij}\,\varphi,\qquad
\int_U Dw_{ij}\,\partial_k\varphi=-\int_U D^2w_{ijk}\,\varphi,\qquad
\int_U w_i\,\dt\varphi=-\int_U D_tw_i\,\varphi .
\]
Such data are unique almost everywhere. We then write $\nabla w$,
$\nabla^2 w$, $\dt w$ for them and $\Lop w=\dt w+\Delta w$ for the field
$D_tw_i+\sum_j D^2w_{ijj}$; inequalities involving these are understood almost
everywhere. The class $W^{2,1}_2(U)$ consists of those $w$ with
$\abs{w}+\abs{\nabla w}+\abs{\nabla^2w}+\abs{\dt w}\in L^2(U)$.

\section{Solution classes and main results}\label{sec:leray-hopf}

This section fixes the solution notions and states the main results. The
suitable weak solutions are exactly those of the formalized
Caffarelli--Kohn--Nirenberg theorem~\cite{CKNLean}; the Leray--Hopf solutions
transcribe (1.3)--(1.7) of~\cite{ESS2003}.

\paragraph{Test functions and suitable weak solutions.}
For open $\Omega\subset\R^3$ and an open interval $I\subset\R$ let
$\Cc(\Omega\times I)$ denote the smooth functions with compact support in
$\Omega\times I$. A velocity is a map $u\colon\R^3\times\R\to\R^3$, and its
weak spatial gradient is carried as separate data
$Du\colon\R^3\times\R\to\R^{3\times3}$ with $Du_{ij}=\partial_j u_i$.

\begin{definition}[Suitable weak solution; {\cite[Definition~1.5]{CKNLean}}]\label{def:sws}
Let $q>\nf52$, $\Omega\subset\R^3$ open, $I\subset\R$ an open interval, and
$f\in\Lloc{q}(\Omega\times I;\R^3)$. The quadruple $(u,Du,p,f)$ is a
\emph{suitable weak solution on $\Omega\times I$} if the following hold.
\begin{enumerate}[label=(S\arabic*)]
\item For every open $\Omega'$ with compact closure
  $\overline{\Omega'}\subset\Omega$ and every interval $J$ with compact closure
  $\overline J\subset I$: $u$, $Du$, $p$, $f$ are measurable on
  $\Omega'\times J$;
  $\esssup_{t\in J}\int_{\Omega'}\abs{u(x,t)}^2\,dx<\infty$;
  $\iint_{\Omega'\times J}\abs{u}^2+\abs{Du}^2<\infty$;
  $p\in L^{3/2}(\Omega'\times J)$; $f\in L^q(\Omega'\times J)$; and for
  a.e.\ $t\in J$, $Du(\cdot,t)$ is the weak gradient of $u(\cdot,t)$ on
  $\Omega'$.
\item $\iint u\cdot\nabla\psi=0$ for every $\psi\in\Cc(\Omega\times I)$.
\item For every $\varphi\in\Cc(\Omega\times I;\R^3)$,
  \[
  \iint\Bigl[-u\cdot\dt\varphi-u_iu_j\,\partial_j\varphi_i+Du_{ij}\,\partial_j\varphi_i
    -p\,\dv\varphi-f\cdot\varphi\Bigr]=0 .
  \]
\item For every $\psi\in\Cc(\Omega\times I)$ with $\psi\ge0$,
  \[
  2\iint\abs{Du}^2\psi\le\iint\Bigl[\abs{u}^2(\dt\psi+\Delta\psi)
    +(\abs{u}^2+2p)\,u\cdot\nabla\psi+2(f\cdot u)\,\psi\Bigr].
  \]
\end{enumerate}
\end{definition}

A point $z_0\in\Omega\times I$ is \emph{regular} if $u$ agrees almost
everywhere on some open neighborhood $N\subset\Omega\times I$ of $z_0$ with a
parabolically H\"older continuous function~\cite[Definition~2.2]{CKNLean}; the
\emph{singular set} is the set of points of $\Omega\times I$ that are not
regular.

\paragraph{Divergence-free data.}
Following~\cite{ESS2003}, $\J$ is the closure in $L^2(\R^3;\R^3)$ of the
smooth compactly supported divergence-free fields. We use the equivalent
description by the weak divergence.

\begin{lemma}\label{lem:J-weak-div}
A field $a\in L^2(\R^3;\R^3)$ belongs to $\J$ if and only if
$\int_{\R^3}a\cdot\nabla\psi=0$ for every $\psi\in\Cc(\R^3)$.
\end{lemma}
\begin{proof}[Proof of Lemma~\ref{lem:J-weak-div}]
If $a\in\J$, choose $a_k\in C_c^\infty(\R^3;\R^3)$ with
$\dv a_k=0$ and $a_k\to a$ in $L^2$. For $\psi\in C_c^\infty(\R^3)$,
integration by parts gives $\int a_k\cdot\nabla\psi=0$, and
$\abs{\int(a_k-a)\cdot\nabla\psi}\le
\norm{a_k-a}_2
\norm{\nabla\psi}_2\to0$. This proves the weak-divergence
identity.

Conversely, suppose that $a\in L^2$ satisfies that identity. With the
Fourier convention $\widehat f(\xi)=\int e^{-2\pi i x\cdot\xi}f(x)\,dx$,
for $\delta>0$ define a vector potential $A_\delta$ and its curl $b_\delta$
by
\begin{equation}\label{eq:J-vector-potential}
 \widehat{A_\delta}(\xi)=
   \frac{i\,\xi\times\widehat a(\xi)}{2\pi(|\xi|^2+\delta^2)},
 \qquad
 \widehat{b_\delta}(\xi)=
   \frac{|\xi|^2}{|\xi|^2+\delta^2}\widehat a(\xi).
\end{equation}
The weak divergence identity gives $\xi\cdot\widehat a(\xi)=0$ almost
everywhere, so $b_\delta=\curl A_\delta$. The multipliers show that
$A_\delta$ and its first weak derivatives belong to $L^2$. Spatially
mollifying this $H^1$ potential gives smooth $L^2$ potentials whose curls
converge to $b_\delta$ in $L^2$; cutting off each smooth potential on
expanding balls shows that its curl belongs to $\J$. Thus $b_\delta\in\J$.

For $\delta_n=1/(n+1)$ the curl multiplier is bounded by one and converges
to one almost everywhere. Plancherel and dominated convergence therefore give
\begin{equation}\label{eq:J-cutoff-convergence}
 \|b_{\delta_n}-a\|_2\longrightarrow0.
\end{equation}
The closedness of $\J$ in $L^2$ now gives $a\in\J$. This route regularizes
the Fourier potential and passes to its curl in $L^2$; it does not require an
$L^6$ potential estimate or a prescribed cutoff-curl sequence for $a$.
\end{proof}

\paragraph{Leray--Hopf solutions.}
The velocity of a Leray--Hopf solution is a function defined at every time,
and it is this function, not an equivalence class, that the definition
constrains at every time. The same function is used as the velocity of a
suitable weak solution.

\begin{definition}[Leray--Hopf solution]\label{def:leray-hopf}
Let $0<T<\infty$ and $a\in\J$. The pair $(u,Du)$ is a \emph{Leray--Hopf
solution on $\R^3\times[0,T]$ with initial datum $a$} if the following hold.
\begin{enumerate}[label=(LH\arabic*)]
\item\label{LH:class} \emph{Energy class, \cite[(1.3)]{ESS2003}.} $u$ and $Du$
  are measurable on $\R^3\times(0,T)$;
  $\esssup_{t\in(0,T)}\int\abs{u(x,t)}^2dx<\infty$;
  $\iint_{\R^3\times(0,T)}\abs{u}^2+\abs{Du}^2<\infty$; and for a.e.\
  $t\in(0,T)$, $Du(\cdot,t)$ is the weak gradient of $u(\cdot,t)$ on $\R^3$
  and $\int u(x,t)\cdot\nabla\psi(x)\,dx=0$ for every $\psi\in\Cc(\R^3)$.
\item\label{LH:weakcont} \emph{Weak continuity, \cite[(1.4)]{ESS2003}.} For
  every $w\in L^2(\R^3;\R^3)$ the map $t\mapsto\int u(x,t)\cdot w(x)\,dx$ is
  continuous on $[0,T]$.
\item\label{LH:eq} \emph{Weak equation, \cite[(1.5)]{ESS2003}.} For every
  $\varphi\in\Cc(\R^3\times(0,T);\R^3)$ with $\dv\varphi=0$,
  \[
  \iint\Bigl[-u\cdot\dt\varphi-u_iu_j\,\partial_j\varphi_i+Du_{ij}\,\partial_j\varphi_i\Bigr]=0 .
  \]
\item\label{LH:energy} \emph{Energy inequality, \cite[(1.6)]{ESS2003}.} For
  every $t_0\in[0,T]$,
  \[
  \frac12\int\abs{u(x,t_0)}^2dx+\int_0^{t_0}\!\!\int\abs{Du}^2\,dx\,dt\le\frac12\int\abs{a}^2dx .
  \]
\item\label{LH:initial} \emph{Initial datum, \cite[(1.7)]{ESS2003}.}
  $\int\abs{u(x,t)-a(x)}^2dx\to0$ as $t\downarrow0$.
\end{enumerate}
A \emph{global Leray--Hopf solution} with datum $a$ is a pair $(u,Du)$ that is
a Leray--Hopf solution on $\R^3\times[0,T]$ for every $T>0$.
\end{definition}

\begin{remark}[The global notion]\label{rem:global-LH}
\cite{ESS2003} define the global notion by replacing $[0,T]$ with $[0,\infty)$
in (1.3)--(1.7). Read literally, (1.3) would then require
$u\in L^2(0,\infty;W^{1,2})$, that is, $\int_0^\infty\!\int\abs{u}^2<\infty$,
which Leray's solutions need not satisfy; the energy inequality controls
only $\sup_t\int\abs{u}^2$ and $\int_0^\infty\!\int\abs{\nabla u}^2$. We therefore
define global solutions through every finite interval, which is the reading
under which~\cite[Theorem~1.1]{ESS2003} holds.
\end{remark}

\paragraph{Main results.}
The first theorem is Leray's existence theorem, with the constructed
solution suitable in the sense of Definition~\ref{def:sws}.

\begin{theorem}[Existence]\label{thm:leray}
For every $a\in\J$ there are $u$, $Du$ and $p$ such that $(u,Du)$ is a global
Leray--Hopf solution with datum $a$, and $(u,Du,p,0)$ is a suitable weak
solution on $\R^3\times(0,\infty)$ for every $q>\nf52$.
\end{theorem}

Combined with the Caffarelli--Kohn--Nirenberg theorem~\cite[Theorem~2.10]{CKNLean},
this gives the headline of~\cite{CKN1982}.

\begin{corollary}\label{cor:ckn-headline}
For every $a\in\J$ there is a global Leray--Hopf solution $(u,Du)$ with datum
$a$ whose singular set in $\R^3\times(0,\infty)$ has zero one-dimensional
parabolic Hausdorff measure.
\end{corollary}

Every Leray--Hopf solution has a pressure.

\begin{theorem}[Associated pressure]\label{thm:assoc-pressure}
Let $(u,Du)$ be a Leray--Hopf solution on $\R^3\times[0,T]$ with datum
$a\in\J$. There is $p\in L^{5/3}(\R^3\times(0,T))$ such that (S3) of
Definition~\ref{def:sws} holds on $\R^3\times(0,T)$ with $f=0$. If in
addition $\esssup_{t\in(0,T)}\norm{u(\cdot,t)}_{L^3(\R^3)}<\infty$, then
this $p$ also belongs to $L^\infty((0,T);L^{3/2}(\R^3))$.
\end{theorem}

\paragraph{Forced Leray--Hopf solutions.}
The same all-time velocity representative is used in the forced equation,
the work term, and the initial trace.

\begin{definition}[Forced Leray--Hopf solution]\label{def:forced-leray-hopf}
Let $0<T<\infty$, $a\in\J$, and
$f\in L^2(\R^3\times(0,T);\R^3)$. The pair $(u,Du)$ is a
\emph{forced Leray--Hopf solution on $\R^3\times[0,T]$ with datum $a$ and
force $f$} if it satisfies (LH1), (LH2), and (LH5) of
Definition~\ref{def:leray-hopf}, and also the following two clauses:
\begin{enumerate}[label=(FLH\arabic*)]
\item For every $\varphi\in C_c^\infty(\R^3\times(0,T);\R^3)$ with
  $\dv\varphi=0$,
  \[
  \iint\Bigl[-u\cdot\partial_t\varphi-u_i u_j\,\partial_j\varphi_i
       +Du_{ij}\,\partial_j\varphi_i-f\cdot\varphi\Bigr]\,dx\,dt=0\,.
  \]
\item For every $t\in[0,T]$,
  \[
  \frac12\int_{\R^3}|u(x,t)|^2\,dx
  +\int_0^t\!\int_{\R^3}|Du|^2\,dx\,ds
  \le\frac12\int_{\R^3}|a|^2\,dx
  +\int_0^t\!\int_{\R^3}f\cdot u\,dx\,ds\,.
  \]
\end{enumerate}
A \emph{global forced Leray--Hopf solution} is one pair $(u,Du)$ that is
a forced Leray--Hopf solution on every finite interval, with the same
representative and force.
\end{definition}

The force work is an ordinary Lebesgue integral: on every finite slab,
$u\in L^\infty_tL^2_x$ and $f\in L^2_{t,x}$, so $f\cdot u\in L^1_{t,x}$.

\begin{theorem}[Forced suitable existence]\label{thm:leray-forced}
Let $a\in\J$ and let $f\colon\R^3\times(0,\infty)\to\R^3$ satisfy
$f\in L^2(\R^3\times(0,T))$ for every finite $T$ and
$f\in L^{q_0}_{\mathrm{loc}}(\R^3\times(0,\infty))$ for some
$q_0>5/2$. There are one velocity $u$, weak gradient $Du$, and pressure
$p$ such that $(u,Du)$ is a global forced Leray--Hopf solution with datum
$a$. For every $q>5/2$ for which
$f\in L^q_{\mathrm{loc}}(\R^3\times(0,\infty))$, the same quadruple
$(u,Du,p,f)$ is a suitable weak solution on
$\R^3\times(0,\infty)$ in the sense of
Definition~\ref{def:sws}. In particular, its local energy inequality
contains the term $2(f\cdot u)\psi$.
\end{theorem}

\begin{theorem}[Forced singular-set conclusion]
\label{thm:leray-forced-singularSet}
Under the hypotheses of Theorem~\ref{thm:leray-forced}, for every
$q>5/2$ for which $f\in L^q_{\mathrm{loc}}$, the witness there satisfies
\[
 \mathcal H^1_{\mathrm{par}}\bigl(\operatorname{Sing}(u)\bigr)=0\,.
\]
\end{theorem}

\begin{theorem}[Forced associated pressure]\label{thm:assoc-pressure-forced}
Let $(u,Du)$ be a forced Leray--Hopf solution on
$\R^3\times[0,T]$ with datum $a\in\J$ and
$f\in L^2(\R^3\times(0,T);\R^3)$. Let
$p_N=P[u\otimes u]$ and let $p_f$ be the force pressure of
\cite[Lemma~13.1]{CKNLean}. Then
\[
 p_N\in L^{5/3}(\R^3\times(0,T)),\qquad
 p_f\in L^2((0,T);L^6(\R^3)),\qquad
 p=p_N+p_f\in L^{3/2}_{\mathrm{loc}}(\R^3\times(0,T)).
\]
For every $\varphi\in C_c^\infty(\R^3\times(0,T);\R^3)$,
\[
 \iint\Bigl[-u\cdot\partial_t\varphi-u_i u_j\,\partial_j\varphi_i
       +Du_{ij}\,\partial_j\varphi_i-p\,\dv\varphi
       -f\cdot\varphi\Bigr]\,dx\,dt=0\,.
\]
When $f=0$, this pressure is the pressure of Theorem~\ref{thm:assoc-pressure}
and has all of its stated global and conditional integrability properties.
\end{theorem}


The next theorem is the local result~\cite[Theorem~1.4]{ESS2003}. Its
hypotheses are those of a suitable weak solution without the local energy
inequality, imposed on the whole cylinder $B_1\times(-1,0)$ rather than on
compactly contained subcylinders, together with the $L^\infty L^3$ bound.

\begin{theorem}[Local regularity]\label{thm:ess-local}
Let $u$, $Du$, $p$ be measurable on $Q=B_1\times(-1,0)$ with
\begin{gather*}
\esssup_{t\in(-1,0)}\int_{B_1}\abs{u(x,t)}^2dx<\infty,\qquad
\iint_Q\abs{u}^2+\abs{Du}^2<\infty,\qquad
p\in L^{3/2}(Q),\\
\esssup_{t\in(-1,0)}\int_{B_1}\abs{u(x,t)}^3dx<\infty ,
\end{gather*}
such that $Du(\cdot,t)$ is the weak gradient of $u(\cdot,t)$ on $B_1$ for a.e.\
$t\in(-1,0)$, and (S2) and (S3) of Definition~\ref{def:sws} hold on $Q$ with
$f=0$. Then there is $\gamma\in(0,1]$ and a function $\tilde u$ equal to $u$
almost everywhere on $B_{1/2}\times(-\nf14,0)$ that is parabolically
H\"older continuous with exponent $\gamma$ on the closure of
$B_{1/2}\times(-\nf14,0)$.
\end{theorem}

The global result is~\cite[Theorem~1.3]{ESS2003} in the form of its step
(3.5): the solution has no singular points.

\begin{theorem}[Regularity of $L^\infty L^3$ Leray--Hopf solutions]\label{thm:ess-global}
Let $(u,Du)$ be a Leray--Hopf solution on $\R^3\times[0,T]$ with datum
$a\in\J$, and suppose
$\esssup_{t\in(0,T)}\int_{\R^3}\abs{u(x,t)}^3dx<\infty$. Then every point of
$\R^3\times(0,T)$ is a regular point of $u$.
\end{theorem}

\begin{remark}[Smoothness on positive-time intervals]\label{rem:not-claimed}
Theorem~\ref{thm:ess-l5-unique} records the $L^5$ and uniqueness conclusions of~\cite[Theorem~1.3]{ESS2003}. The Prodi--Serrin theorem~\cite[Theorem~8.17, printed p.~173]{RobinsonRodrigoSadowski2016} further implies that its $L^5$ solution is smooth on every closed interval $[\delta,T]$ with $0<\delta<T$; smoothness is a consequence, not a separate clause of the theorem.
\end{remark}

%======================================================================
\part*{Part I. Carleman inequalities, unique continuation and backward uniqueness}
%======================================================================

\section{Carleman inequalities}\label{sec:carleman}

The weighted estimates below convert control of the backward heat operator into control of a field and its gradient. The first weight amplifies small times while remaining Gaussian in space; the second also grows rapidly in the normal variable $x_3$ on the translated half-space.

Throughout this section, $\Lop=\dt+\Delta$. All vector and matrix norms are Euclidean, and derivatives and integrations are componentwise. The calculations are the compact-support versions of (6.2)--(6.11) and (6.13)--(6.26) in~\cite{ESS2003}.

\begin{proposition}\label{prop:carleman-gauss}
There is an absolute constant
\[
 c_0=e^{4/3}(9+2\sqrt{6})\,.
\]
For every $a>0$ and every $w\in \Cc(\R^3\times(0,2);\R^3)$, with
$h(t)=t e^{(1-t)/3}$, one has
\begin{equation}\label{eq:carleman-gauss}
\begin{aligned}
 \int_{\R^3\times(0,2)} h^{-2a}(t)e^{-|x|^2/(4t)}
 \left(\frac{a}{t}|w|^2+|\nabla w|^2\right)\,dx\,dt\\
 \le c_0\int_{\R^3\times(0,2)}h^{-2a}(t)e^{-|x|^2/(4t)}|Lw|^2\,dx\,dt\,.
\end{aligned}
\end{equation}
\end{proposition}

\begin{proof}
Set $q=a+1$, $\phi(x,t)=-|x|^2/(8t)-q\log h(t)$, and $v=e^\phi w$. The field $v$ is smooth and compactly supported in $\R^3\times(0,2)$. Define the conjugated operator
\[
 \mathcal L_\phi v:=e^\phi Lw
 =\dt v-\dv(v\otimes\nabla\phi)-\nabla v\nabla\phi+\Delta v
   +( |\nabla\phi|^2-\dt\phi)v\,.
\]
Here $\dv(v\otimes\nabla\phi)_i=\sum_j\partial_j(v_i\partial_j\phi)$ and $(\nabla v\nabla\phi)_i=\sum_j(\partial_jv_i)(\partial_j\phi)$. In particular,
\[
 \mathcal L_\phi v=(\dt+\Delta)v-2\nabla\phi\cdot\nabla v
                    +( |\nabla\phi|^2-\dt\phi-\Delta\phi)v\,.
\]

For this compactly supported field, write $t\mathcal L_\phi=S+A$, where
\begin{align}
 Sv&=t\Delta v+t( |\nabla\phi|^2-\dt\phi)v-\tfrac12v,\label{eq:carleman-S}\\
 Av&=\tfrac12\big(\dt(tv)+t\dt v\big)
       -t\big(\dv(v\otimes\nabla\phi)+\nabla v\nabla\phi\big)\,.\label{eq:carleman-A}
\end{align}
For compactly supported $v$, the product rule and the flux identity below compute the integrated cross term directly. The time integrations have no endpoint terms because $v$ is supported in a compact subset of $(0,2)$, and the spatial integrations have no boundary terms because $v$ has compact spatial support.

All integrals in this proof are over $\R^3\times(0,2)$ unless specified otherwise. Write $\mathcal D_\phi(v)$ for the density in the following direct product-rule computation, with $i,j\in\{1,2,3\}$ summed:
\begin{align}\label{eq:carleman-commutator}
 I:=\int\mathcal D_\phi(v)\,dx\,dt
 &=4\int t^2\left(\phi_{,ij}\,\partial_i v\cdot\partial_jv
             +\phi_{,ij}\phi_{,i}\phi_{,j}|v|^2\right)\\
 &\quad+\int t^2|v|^2
       \left(\dt^2\phi-2\dt|\nabla\phi|^2-\Delta^2\phi\right)
       +\int t|\nabla v|^2
       -\int t|v|^2\left(|\nabla\phi|^2-\dt\phi\right)\,.\nonumber
\end{align}
Consequently, expanding the square and using the integrated flux identity below gives
\begin{equation}\label{eq:carleman-commutator-energy}
 P:=\int t^2|\mathcal L_\phi v|^2\,dx\,dt
 =\int|Sv|^2\,dx\,dt+\int|Av|^2\,dx\,dt+I.
\end{equation}
For vector fields, sum the scalar fluxes over components:
\begin{align*}
 F_0={}&t^2\bigl(|\nabla\phi|^2-\dt\phi\bigr)|v|^2
          -t^2|\nabla v|^2-\tfrac12t|v|^2,\\
 F_i={}&2t^2\partial_i v\cdot\dt v+t\,v\cdot\partial_i v
       -4t^2\partial_i v\cdot\left(\sum_{j=1}^3\phi_{,j}\partial_jv\right)
       +2t^2\phi_{,i}|\nabla v|^2\\
 &\quad-2t^2(\Delta\phi)\,v\cdot\partial_i v
       +t^2\partial_i(\Delta\phi)|v|^2
       -2t^2\bigl(|\nabla\phi|^2-\dt\phi\bigr)\phi_{,i}|v|^2
       +t\phi_{,i}|v|^2,\qquad i=1,2,3.
\end{align*}
Expanding $2(Sv)\cdot(Av)$ and differentiating the displayed fluxes by the product rule, with equality of mixed space-time derivatives, gives the pointwise identity
\begin{equation*}
 2(Sv)\cdot(Av)=\mathcal D_\phi(v)+\dt F_0+\sum_{i=1}^3\partial_iF_i\,.
\end{equation*}
The fluxes and their derivatives are supported in $\supp v$, a compact subset of the open cylinder, so their time and spatial integrals vanish. Thus $I=2\int Sv\cdot Av$. Expanding $|t\mathcal L_\phi v|^2=|Sv+Av|^2$ pointwise gives the displayed identity for $P$; nonnegativity of the two square integrals yields $I\leq P$. The explicit value below gives $I\geq0$. This integrated estimate requires no separate space-time pairing identities for $S$ and $A$.

For the present $\phi$, let $r=h'/h=1/t-1/3$. Then
\[
 \nabla\phi=-\frac{x}{4t},\qquad \phi_{,ij}=-\frac{\delta_{ij}}{4t},
 \qquad |\nabla\phi|^2=\frac{|x|^2}{16t^2},
 \qquad |\nabla\phi|^2-\dt\phi=-|\nabla\phi|^2+q r,
\]
while $\Delta^2\phi=0$ and
\[
 \dt^2\phi-2\dt|\nabla\phi|^2=\frac{q}{t^2}.
\]
Substituting these identities into~\eqref{eq:carleman-commutator}, the $\nabla v$ terms cancel, as do the two terms containing $t|\nabla\phi|^2|v|^2$. The remainder is
\begin{equation}\label{eq:carleman-I}
 I=q\int(1-tr)|v|^2\,dx\,dt
   =\frac{q}{3}\int t|v|^2\,dx\,dt\,.
\end{equation}
In particular,~\eqref{eq:carleman-commutator-energy} gives $0\le I\le P$.

The pointwise identity
\begin{equation}\label{eq:carleman-gradient-pointwise}
 |\nabla v|^2=\tfrac12(\dt+\Delta)|v|^2-v\cdot(\dt v+\Delta v)
\end{equation}
holds componentwise. To integrate it, use the expression for $\mathcal L_\phi v$ above. Spatial integration by parts yields
$\int v\cdot(2\nabla\phi\cdot\nabla v)=-\int(\Delta\phi)|v|^2$; this has no boundary term because $v$ is compactly supported. The $\Delta\phi$ terms then cancel, giving
\begin{equation}\label{eq:carleman-gradient-integrated}
 \int t^2|\nabla v|^2
 =-\int t|v|^2-\int t^2v\cdot\mathcal L_\phi v
   +\int t^2|v|^2( |\nabla\phi|^2-\dt\phi)\,.
\end{equation}
Here the time integration by parts has no endpoint terms because the support is compact in $(0,2)$, and the spatial Laplacian integration by parts has no boundary term because the support is compact in $\R^3$.

Let
\[
 E=\int t^2\left(|\nabla v|^2+|v|^2|\nabla\phi|^2\right)\,.
\]
Combining~\eqref{eq:carleman-gradient-integrated} with
$|\nabla\phi|^2-\dt\phi=-|\nabla\phi|^2+q h'/h$ gives
\[
 E=-\int t|v|^2-\int t^2v\cdot\mathcal L_\phi v
       +q\int t^2\frac{h'}h|v|^2
 \le 3I+\left|\int t^2v\cdot\mathcal L_\phi v\right|\,.
\]
Indeed, the coefficient of $|v|^2$ before the last cross term is
$-t+qt^2h'/h=at-qt^2/3\le qt$, and $q\int t|v|^2=3I$ by~\eqref{eq:carleman-I}. Since $0<t<2$,~\eqref{eq:carleman-I} and Cauchy--Schwarz imply
\[
 \left|\int t^2v\cdot\mathcal L_\phi v\right|
 \le \left(\int t^2|v|^2\right)^{1/2}P^{1/2}
 \le (6I)^{1/2}P^{1/2}\le\sqrt6\,P\,.
\]
It follows that
\begin{equation}\label{eq:carleman-E}
 E\le(3+\sqrt6)P\,.
\end{equation}
Finally, $e^\phi|\nabla w|\le|\nabla v|+|v||\nabla\phi|$, so
$e^{2\phi}|\nabla w|^2\le2(|\nabla v|^2+|v|^2|\nabla\phi|^2)$. Since
$e^{2\phi}=h^{-2(a+1)}e^{-|x|^2/(4t)}$ and
$(th^{-1})^2=e^{-2(1-t)/3}$ satisfies $e^{-2/3}\le(th^{-1})^2\le e^{2/3}$ on $(0,2)$,~\eqref{eq:carleman-E} yields
\begin{align*}
 \int h^{-2a}e^{-|x|^2/(4t)}
       \left((a+1)\frac{|w|^2}{t}+|\nabla w|^2\right)
 &\le e^{2/3}(9+2\sqrt6)P\\
 &\le e^{4/3}(9+2\sqrt6)
       \int h^{-2a}e^{-|x|^2/(4t)}|Lw|^2\,.
\end{align*}
The first inequality uses $(a+1)\ge a$ and the lower bound on $(th^{-1})^2$; the second uses its upper bound. This proves~\eqref{eq:carleman-gauss} with the stated $c_0$.
\end{proof}

\begin{proposition}\label{prop:carleman-halfspace}
For every $\alpha\in(1/2,1)$ there are constants $a_0(\alpha)>0$ and $c_*(\alpha)>0$ such that, for every $a>a_0(\alpha)$ and every
$w\in \Cc(\{x_3>1\}\times(0,1);\R^3)$, the following estimate holds with
$x'=(x_1,x_2)$ and
\[
 \phi(x,t)=-\frac{|x'|^2}{8t}+a(1-t)x_3^{2\alpha}t^{-\alpha}\,:
\]
\begin{equation}\label{eq:carleman-halfspace}
\begin{aligned}
 \int_{\{x_3>1\}\times(0,1)}t^2e^{2\phi}
       \left(a\frac{|w|^2}{t^2}+\frac{|\nabla w|^2}{t}\right)\,dx\,dt\\
 \le c_*(\alpha)\int_{\{x_3>1\}\times(0,1)}t^2e^{2\phi}|Lw|^2\,dx\,dt\,.
\end{aligned}
\end{equation}
One may take $a_0(\alpha)=2$ and
\[
 c_*(\alpha)=5+\frac{7}{2(2\alpha-1)}\,.
\]
\end{proposition}

\begin{proof}
Fix $\alpha\in(1/2,1)$, $a\ge2$, and a test field as in the statement. Set $Q=\{x_3>1\}\times(0,1)$, $v=e^\phi w$, and $\mathcal L_\phi v=e^\phi Lw$. The function $x_3^{2\alpha}$ is smooth on $x_3>1$, and $v$ has compact support in $Q$. Apply the conjugation and commutator identities above on $Q$, with all their integrations by parts justified by this compact support: there are no contributions from spatial infinity, the plane $x_3=1$, or the time endpoints $t=0,1$.

The gradients and Hessians of the two parts of $\phi$ are
\begin{align}
 \nabla\phi^{(1)}&=-\frac{x'}{4t},&
 \nabla\phi^{(2)}&=2\alpha a(1-t)x_3^{2\alpha-1}t^{-\alpha}e_3,
 \qquad \phi^{(1)}=-\frac{|x'|^2}{8t},\quad
 \phi^{(2)}=a(1-t)x_3^{2\alpha}t^{-\alpha}.
 \label{eq:carleman-half-gradients}
\end{align}
Their gradients are orthogonal, so
\begin{equation}\label{eq:carleman-half-orthogonal}
 \nabla\phi^{(1)}\cdot\nabla\phi^{(2)}=0,
 \qquad |\nabla\phi|^2=|\nabla\phi^{(1)}|^2+|\nabla\phi^{(2)}|^2\,.
\end{equation}
The only nonzero entries of the Hessians are
\begin{equation}\label{eq:carleman-half-hessians}
 \phi^{(1)}_{,ij}=-\frac{\delta_{ij}}{4t}\quad(1\le i,j\le2),
 \qquad
 \phi^{(2)}_{,33}=2\alpha(2\alpha-1)a(1-t)x_3^{2\alpha-2}t^{-\alpha}.
\end{equation}
It follows that
\begin{equation}\label{eq:carleman-half-hess-grad}
 \phi_{,ij}\phi_{,i}\phi_{,j}
 =-\frac1{4t}|\nabla\phi^{(1)}|^2
  +\frac{2\alpha(2\alpha-1)a(1-t)}{t^\alpha}
       x_3^{2\alpha-2}|\nabla\phi^{(2)}|^2
 \ge-\frac{|x'|^2}{64t^3}\,.
\end{equation}

To split the commutator~\eqref{eq:carleman-commutator}, define, for $s=1,2$,
\begin{align*}
 I_s={}&4\int_Q t^2\left(\phi^{(s)}_{,ij}\partial_i v\cdot\partial_jv
       +\phi^{(s)}_{,ij}\phi^{(s)}_{,i}\phi^{(s)}_{,j}|v|^2\right)\,dx\,dt\\
 &+\int_Q t^2|v|^2\left(\dt^2\phi^{(s)}-2\dt|\nabla\phi^{(s)}|^2-\Delta^2\phi^{(s)}
       -\frac1t|\nabla\phi^{(s)}|^2+\frac1t\dt\phi^{(s)}\right)\,dx\,dt\,.
\end{align*}
The Hessians and gradients in~\eqref{eq:carleman-half-gradients}--\eqref{eq:carleman-half-hessians} have disjoint tangential and normal coordinates. Thus the mixed terms vanish, and~\eqref{eq:carleman-commutator} becomes
\begin{equation}\label{eq:carleman-half-I-split}
 I=I_1+I_2+\int_Q t|\nabla v|^2\,dx\,dt\,.
\end{equation}
For $\phi^{(1)}$, direct differentiation gives
\begin{align*}
 4t^2\phi^{(1)}_{,ij}\partial_i v\cdot\partial_jv
 &=-t|\nabla'v|^2,\\
 4t^2\phi^{(1)}_{,ij}\phi^{(1)}_{,i}\phi^{(1)}_{,j}
 &=-t|\nabla\phi^{(1)}|^2,\\
 t^2\left(\dt^2\phi^{(1)}-2\dt|\nabla\phi^{(1)}|^2-\Delta^2\phi^{(1)}
       -t^{-1}|\nabla\phi^{(1)}|^2+t^{-1}\dt\phi^{(1)}\right)
 &=t|\nabla\phi^{(1)}|^2\,.
\end{align*}
Therefore $I_1=-\int_Qt|\nabla'v|^2$ and
\begin{equation}\label{eq:carleman-half-I1}
 I=\int_Q t|\partial_3v|^2\,dx\,dt+I_2\,.
\end{equation}

For the second part of the weight, $\phi^{(2)}_{,33}>0$ on $Q$, and
$\phi^{(2)}_{,33}(\partial_3\phi^{(2)})^2\ge0$. Dropping these nonnegative terms from $I_2$ gives
\begin{equation}\label{eq:carleman-half-A}
 I_2\ge\int_Q t^2|v|^2(A_1+A_2+A_3)\,dx\,dt,
\end{equation}
where
\[
 A_1=-\dt|\nabla\phi^{(2)}|^2,\qquad
 A_2=A_1-\Delta^2\phi^{(2)}-\frac1t|\nabla\phi^{(2)}|^2,\qquad
 A_3=\dt^2\phi^{(2)}+\frac1t\dt\phi^{(2)}\,.
\]
Differentiating $x_3^{2\alpha}$ four times yields
\[
 \Delta^2\phi^{(2)}=2\alpha(2\alpha-1)(2\alpha-2)(2\alpha-3)
              a(1-t)x_3^{2\alpha-4}t^{-\alpha}.
\]
Using the exact derivative of $|\nabla\phi^{(2)}|^2$, and
$2t+(2\alpha-1)(1-t)\ge2\alpha-1$, we obtain
\begin{equation}\label{eq:carleman-half-A2}
 A_2\ge\frac{1-t}{t^\alpha}x_3^{2\alpha-4}a(2\alpha-1)
 \left[\frac{4\alpha^2a x_3^{2\alpha+2}}{t^{\alpha+1}}
              -2\alpha(2\alpha-2)(2\alpha-3)\right]\,.
\end{equation}
On $Q$, $x_3^{2\alpha+2}t^{-(\alpha+1)}\ge1$. The bracket in~\eqref{eq:carleman-half-A2} is therefore at least
\[
 8\alpha^2-2\alpha(2\alpha-2)(2\alpha-3)
 =-4\alpha(2\alpha-1)(\alpha-3)>0,
\]
where the last inequality uses $1/2<\alpha<1$. Hence $A_2>0$ on $Q$ for every $a\ge2$.

A direct differentiation also gives
\begin{equation}\label{eq:carleman-half-A3}
 A_3=a x_3^{2\alpha}t^{-\alpha-2}
       \left(\alpha^2-(1-\alpha)^2t\right)
 \ge a(2\alpha-1)\frac{x_3^{2\alpha}}{t^{\alpha+2}}\,.
\end{equation}
For $A_1$, direct calculation gives
\begin{align*}
 A_1-\frac1t|\nabla\phi^{(2)}|^2
 &=4\alpha^2a^2x_3^{4\alpha-2}(1-t)t^{-2\alpha-1}
       \left(2t+(2\alpha-1)(1-t)\right)\\
 &\ge(2\alpha-1)4\alpha^2a^2(1-t)x_3^{4\alpha-2}t^{-2\alpha-1}\ge0.
\end{align*}
Thus $A_1\ge t^{-1}|\nabla\phi^{(2)}|^2$. From~\eqref{eq:carleman-half-I1}--\eqref{eq:carleman-half-A3}, the commutator energy identity, and $x_3^{2\alpha}t^{-\alpha}\ge1$, we obtain
\begin{equation}\label{eq:carleman-half-mass}
\begin{aligned}
 P&:=\int_Q t^2|\mathcal L_\phi v|^2\,dx\,dt
       \ge I\ge a(2\alpha-1)U+G,\\
 U&:=\int_Q\frac{x_3^{2\alpha}}{t^\alpha}|v|^2\,dx\,dt,
 \qquad G:=\int_Qt|v|^2|\nabla\phi^{(2)}|^2\,dx\,dt\,.
\end{aligned}
\end{equation}
In particular, for $\theta=2\alpha-1>0$ and $V=\int_Q|v|^2\,dx\,dt$, we have
$V\le U$ and $P\ge a\theta U+G$.

Integrating~\eqref{eq:carleman-gradient-pointwise} with the factor $t$ gives
\begin{equation}\label{eq:carleman-half-gradient-id}
 \int_Qt|\nabla v|^2\,dx\,dt
 =-\frac12\int_Q|v|^2\,dx\,dt-\int_Qtv\cdot\mathcal L_\phi v\,dx\,dt
   +\int_Qt|v|^2(|\nabla\phi|^2-\dt\phi)\,dx\,dt\,.
\end{equation}
The time integration by parts has no endpoint terms because $v$ is compactly supported in $(0,1)$; the spatial integrations by parts have no boundary terms because $v$ is compactly supported in $\{x_3>1\}$. Since
$|\nabla\phi^{(1)}|^2-\dt\phi^{(1)}=-|\nabla\phi^{(1)}|^2$, rearranging~\eqref{eq:carleman-half-gradient-id} yields
\begin{align}\label{eq:carleman-half-gradient-rewrite}
 \int_Q t\left(|\nabla v|^2+|v|^2|\nabla\phi|^2\right)\,dx\,dt
 ={}&-\frac12\int_Q|v|^2-\int_Qtv\cdot\mathcal L_\phi v\\
 &+2\int_Qt|v|^2|\nabla\phi^{(2)}|^2
   -\int_Qt|v|^2\dt\phi^{(2)}\,dx\,dt\,.
\end{align}
Moreover,
\[
 -t\dt\phi^{(2)}
 =a\frac{x_3^{2\alpha}}{t^\alpha}\big(\alpha+(1-\alpha)t\big)
 \le a\frac{x_3^{2\alpha}}{t^\alpha},
\]
since $0<t<1$ and $0<\alpha<1$. The identity
$e^\phi\nabla w=\nabla v-v\nabla\phi$ implies
$\tfrac12e^{2\phi}|\nabla w|^2\le|\nabla v|^2+|v|^2|\nabla\phi|^2$. Therefore~\eqref{eq:carleman-half-gradient-rewrite} gives
\begin{equation}\label{eq:carleman-half-gradient-estimate}
 D:=\frac12\int_Qt e^{2\phi}|\nabla w|^2\,dx\,dt
 \le\left|\int_Qtv\cdot\mathcal L_\phi v\,dx\,dt\right|+2G+aU\,.
\end{equation}

It remains to bound the cross term with constants uniform in $a$. Cauchy--Schwarz, followed by Young's inequality with equal weights, gives the explicit absorption
\begin{align*}
 \left|\int_Qtv\cdot\mathcal L_\phi v\right|
 &\le (VP)^{1/2}\le(UP)^{1/2}
   =\left(\frac{a\theta U\,P}{a\theta}\right)^{1/2}\\
 &\le\frac12a\theta U+\frac{1}{2a\theta}P
 \le\left(\frac12+\frac{1}{2a\theta}\right)P
 \le\left(\frac12+\frac{1}{4\theta}\right)P,
\end{align*}
where the penultimate inequality uses $a\theta U\le P$ from~\eqref{eq:carleman-half-mass}, and the last uses $a\ge2$. The same bound gives $aV\le P/\theta$, while~\eqref{eq:carleman-half-mass} gives $G\le P$ and $aU\le P/\theta$. Consequently,
\[
 a\int_Qe^{2\phi}|w|^2\,dx\,dt+\int_Qte^{2\phi}|\nabla w|^2\,dx\,dt
 =aV+2D
 \le\left(5+\frac{7}{2\theta}\right)P.
\]
Since $P=\int_Qt^2e^{2\phi}|Lw|^2\,dx\,dt$ and $\theta=2\alpha-1$, this is~\eqref{eq:carleman-halfspace} with the stated constants.
\end{proof}

\begin{remark}[Departure from \cite{ESS2003}]
The printed proof of Proposition~6.2~\cite[pp.~240--243]{ESS2003} ends by saying that Cauchy--Schwarz and (6.23), (6.26) imply the claim, without displaying the uniform constant or the absorption. The estimates above supply that omitted step: $a\ge2$ makes $A_2$ positive, and Young's inequality bounds the cross term uniformly, giving $a_0(\alpha)=2$ and $c_*(\alpha)=5+7/(2(2\alpha-1))$. The proposition's statement is unchanged.
\end{remark}

\paragraph{Extension to Sobolev fields and a Caccioppoli inequality.}
The uniqueness theorems apply the two inequalities to cut-off fields of class $W^{2,1}_2$, and they need two elementary estimates for the differential inequality $\abs{\Lop w}\le c_1(\abs{w}+\abs{\nabla w})$.

\begin{lemma}\label{lem:carleman-sobolev}
Set $U_G=\R^3\times(0,2)$ and $U_H=\{x_3>1\}\times(0,1)$. The following extensions retain the constants of the stated propositions:
\begin{enumerate}
\item For every $a>0$, the inequality in Proposition~\ref{prop:carleman-gauss} holds for all $w\in W^{2,1}_2(U_G;\R^3)$ compactly supported in $U_G$, with the same absolute constant $c_0$.
\item For every $\alpha\in(1/2,1)$ and every $a>a_0(\alpha)$, the inequality in Proposition~\ref{prop:carleman-halfspace} holds for all $w\in W^{2,1}_2(U_H;\R^3)$ compactly supported in $U_H$, with the same constant $c_*(\alpha)$.
\end{enumerate}
\end{lemma}

\begin{proof}
Fix one of the two open sets $U$ in the statement and let $K=\supp w\Subset U$. Extend $w$ by zero to $\R^3\times\R$. Because $w$ vanishes in a neighborhood of $\partial U$, this extension has as its weak derivatives the zero extensions of $Dw$, $D^2w$, and $D_tw$: testing against a smooth function on $\R^3\times\R$ reduces, after inserting a cutoff equal to one near $K$, to the weak-derivative identities on $U$. Thus, for a standard space-time mollifier $\rho_\delta$,
\[
\partial_j(w*\rho_\delta)=(Dw)_{\cdot j}*\rho_\delta,\qquad
\partial_k\partial_j(w*\rho_\delta)=(D^2w)_{\cdot jk}*\rho_\delta,\qquad
\dt(w*\rho_\delta)=(D_tw)*\rho_\delta\,.
\]
These identities use only the stated spatial weak derivatives through order two and the time weak derivative; no mixed derivative is needed. If $\delta$ is smaller than the distance from $K$ to $\partial U$, then $w*\rho_\delta\in \Cc(U;\R^3)$. All these mollifications are supported in a fixed compact set $K'\Subset U$, on which every scalar weight and coefficient in the two inequalities is bounded. The approximate-identity theorem gives convergence of $w*\rho_\delta$, its first and second spatial derivatives, and its time derivative to the corresponding fields in $L^2(K')$. Hence the weighted $L^2$ norms of the fields and of $\Lop(w*\rho_\delta)$ converge to those of $w$ and $\Lop w$. Apply the appropriate proposition to $w*\rho_\delta$ and pass to the limit. The constants do not depend on $\delta$.
\end{proof}

\begin{remark}[Departure from \cite{ESS2003}]
Propositions 6.1 and 6.2 of~\cite{ESS2003}, printed pp.~238--243, are stated for smooth compactly supported fields. The argument above supplies the density passage needed to use those estimates for compactly supported $W^{2,1}_2$ fields.
\end{remark}

\begin{lemma}\label{lem:caccioppoli}
Let $0<r\le1$, $c_1\ge0$, $x\in\R^3$, and $t\in\R$. Set $Q=\Ball{2r}{x}\times(t,t+4r^2)$. Suppose that $w\in W^{2,1}_2(Q;\R^3)$ and
\[
\abs{\Lop w}\le c_1\bigl(\abs{w}+\abs{\nabla w}\bigr)\quad\text{a.e.\ on }Q.
\]
Then, for the absolute constant $K=256$,
\[
\iint_{\Ball{r}{x}\times(t,t+r^2)}\abs{\nabla w}^2
\le K\bigl(1+c_1^2+r^{-2}\bigr)
\iint_{\Ball{2r}{x}\times(t,t+4r^2)}\abs{w}^2\,.
\]
The same estimate holds when $t$ is the initial time of the domain, provided $w$ has a continuous extension to $\overline{\Ball{2r}{x}}\times\{t\}$; no derivative at that slice is assumed.
\end{lemma}

\begin{proof}
Choose $0<4\varepsilon<r^2$. Take the smooth spatial weight $\eta$
supported in $\Ball{2r}{x}$, equal to one on $\Ball{r}{x}$, with
$|\nabla\eta|\le400/(99r)$. Let $\theta$ be the product of a lower
time transition from zero at $t+\varepsilon$ to one at $t+2\varepsilon$
and an upper transition from one through $t+(101/100)r^2$ to zero at
$t+3r^2$. It is supported in $[t+\varepsilon,t+3r^2]$ and equals one
on $[t+2\varepsilon,t+r^2]$. Its lower transition is nondecreasing and
finishes before $t+r^2$; after $t+r^2$, $\theta$ is nonincreasing and
$|\theta'|\le800/(199r^2)$. Set $\psi(y,s)=\eta(y)^2\theta(s)^2$.
Its product derivatives satisfy
\[
\sum_j|\partial_j\psi|^2\le
  \frac{4(400/99)^2}{r^2}\psi,\qquad
\partial_s\psi\ge-\frac{1600}{199r^2}.
\]
The weak derivatives give the integrated identity
\[
\iint_Q\Lop w\cdot w\psi
=-\frac12\iint_Q|w|^2\partial_s\psi
 -\iint_Q|\nabla w|^2\psi
 -\iint_Q\sum_{i,j}w_i\,\partial_jw_i\,\partial_j\psi,
\]
where $Q=\Ball{2r}{x}\times(t,t+4r^2)$. The time cutoff is supported inside
the open interval, so this identity uses no trace or derivative at $t$.
Applying the differential inequality and Young's inequality gives
\[
\iint_Q|\nabla w|^2\psi
\le2\left(c_1+c_1^2+
  \frac{4(400/99)^2+800/199}{r^2}\right)\iint_Q|w|^2.
\]
The coefficient is at most $256(1+c_1^2+r^{-2})$. Since $\psi=1$ on
$\Ball{r}{x}\times(t+2\varepsilon,t+r^2)$, let $\varepsilon\downarrow0$
and use monotone convergence. This proves the estimate, including at an
initial time, without using a derivative there.
\end{proof}

\begin{lemma}\label{lem:ftc-small-time}
Let $B\subset\R^3$ be open and $\tau>0$. Suppose that $w:B\times(0,\tau)\to\R^3$ has space-time weak time derivative $D_tw\in L^2(B\times(0,\tau))$, extends continuously to $B\times[0,\tau)$, and satisfies $w(x,0)=0$ for every $x\in B$. Then, for every $0<t<\tau$,
\[
\int_B\abs{w(x,t)}^2\,dx
\le t\int_0^t\int_B\abs{D_tw(x,s)}^2\,dx\,ds\,.
\]
\end{lemma}

\begin{proof}
Fix $\Omega\Subset B$ and $T<\tau$. Continuity makes $w$ square-integrable on $\Omega\times(0,T)$, and $D_tw$ is square-integrable there by hypothesis. Testing the weak time-derivative identity with products $\varphi(x)\zeta(s)$ from countable dense families and applying Fubini shows that, for almost every $x\in\Omega$, the section $s\mapsto w(x,s)$ has weak derivative $D_tw(x,s)$ on $(0,T)$. Both sections are in $L^2((0,T))$ for almost every such $x$, so this section has an absolutely continuous representative. That representative agrees almost everywhere in $s$ with the given continuous section; hence they agree for every $s\in(0,T)$, and continuity at zero identifies its value there as zero. Therefore, for almost every $x\in\Omega$ and every $t<T$,
\[
w(x,t)=\int_0^t D_tw(x,s)\,ds,\qquad
\abs{w(x,t)}^2\le t\int_0^t\abs{D_tw(x,s)}^2\,ds
\]
by Cauchy--Schwarz. Integrate over $\Omega$, then exhaust $B$ by relatively compact open subsets and use monotone convergence. Taking $T$ with $t<T<\tau$ proves the claim.
\end{proof}

\begin{remark}[Departure from \cite{ESS2003}]
In §5, Lemma 5.2 of~\cite{ESS2003}, printed pp.~234--236, the zero extension across the initial slice is used without specifying a trace topology; Equation (5.7) in that proof and (5.19) in the proof of Lemma 5.3, printed pp.~234 and 236, are pointwise gradient bounds invoked from parabolic regularity. The estimates above instead provide a spacetime $L^2$ gradient bound in the weak-derivative class and the small-time trace inequality under the stated continuous representative convention.
\end{remark}

\section{Unique continuation across spatial boundaries}\label{sec:uc}

The integral form of Lemma 4.3 replaces the maximum in (4.4) by the
space-time $L^2$ norms already present in (4.1). Its proof keeps the Gaussian
weight in Proposition 6.1 and estimates the cutoff errors directly in $L^2$;
in particular, it does not use the local $L^2$-to-pointwise estimate (4.18)
of~\cite{ESS2003}.

\begin{lemma}[Gaussian estimate from integral vanishing]\label{lem:uc-gaussian}
Let $R,T,c_1>0$ with $T\leq1$, let $x_0\in\R^3$, and let
$w\colon B_R(x_0)\times[0,T)\to\R^3$ be continuous. Suppose that $w$ has
space-time weak derivatives $\nabla w$, $\nabla^2w$, and $\dt w$ on
$B_R(x_0)\times(0,T)$ and
\[
 \abs{w}+\abs{\nabla w}+\abs{\nabla^2w}+\abs{\dt w}
       \in L^2(B_R(x_0)\times(0,T)),\qquad
 \abs{\Lop w}\leq c_1(\abs{w}+\abs{\nabla w})
 \quad\text{a.e.}
\]
Assume also that for every $m\in\N$ there are $C_m,r_m>0$ such that
\begin{equation}\label{eq:uc-integral-vanishing}
 \int_0^{r^2}\int_{B_r(x_0)}\abs{w(y,s)}^2\,dy\,ds
       \leq C_m r^m\qquad(0<r<\min\{r_m,R,\sqrt{T}\}).
\end{equation}
Put
\[
 \mathcal N_{x_0}(R,T;w)=
 \int_0^{3T/4}\int_{B_{3R/4}(x_0)}
       \bigl(\abs{w}^2+T\abs{\nabla w}^2\bigr)\,dy\,ds.
\]
There are absolute constants
\[
 \beta=\frac1{100},\qquad
 \beta_1=\frac38\sqrt{2\beta},\qquad \beta_2=16\beta,
\]
and a number $\gamma=\gamma(c_1)\in(0,3/16)$ such that, whenever
$0<t\leq\gamma T$, $\abs{x-x_0}\leq\beta_1R$, and
$\beta_2t\leq\abs{x-x_0}^2$, one has
\begin{equation}\label{eq:uc-gaussian-box}
 \int_t^{2t}\int_{B_{\sqrt{2t}}(x)}\abs{w(y,s)}^2\,dy\,ds
 \leq C(c_1)\mathcal N_{x_0}(R,T;w)
       \exp\!\left(-\frac{\abs{x-x_0}^2}{2t}\right),
\end{equation}
where $C(c_1)$ depends only on $c_1$ and the absolute constants in
Proposition~\ref{prop:carleman-gauss}. Equivalently, the average of $|w|^2$
over this box is bounded by the right-hand side of~\eqref{eq:uc-gaussian-box}
divided by $2^{3/2}|B_1|t^{5/2}$.
\end{lemma}

\begin{proof}
Translation allows us to take $x_0=0$. Fix an output point $x$ satisfying the
conclusion's spatial conditions. The spatial condition and the definitions of
$\beta_1,\beta_2$ imply
\[
 \abs x+\sqrt{2t}
 \leq\left(\beta_1+\sqrt{\frac{2}{\beta_2}}\,\beta_1\right)R
 =\left(\beta_1+\frac{3}{16}\right)R<R,
 \qquad 2t\leq2\gamma T<\frac{3T}{8}<T.
\]
Thus the box in~\eqref{eq:uc-gaussian-box} lies inside the domain of $w$.
Set $\mu=\sqrt{2\beta}$ and $X=x/\mu$, and put
\[
 \lambda=\sqrt{2t},\qquad \rho=\frac{2\abs{X}}{\lambda},\qquad
 v(y,s)=w(\lambda y,\lambda^2s).
\]
The affine change of variables in the definition of weak derivatives gives
$D_yv=\lambda Dw$, $D_y^2v=\lambda^2D^2w$, and
$D_sv=\lambda^2D_tw$, evaluated at $(\lambda y,\lambda^2s)$. Consequently,
on $Q(\rho,2)=B_\rho\times(0,2)$,
\[
 \abs{\Lop v}\leq c_1\lambda(\abs v+\abs{\nabla v})\quad\text{a.e.}
\]
Here we used $\lambda<1$, which follows from $t\leq\gamma T$,
$T\leq1$, and $\gamma<3/16$. The output point conditions give
\[
 \abs X\leq\frac{\beta_1R}{\mu}=\frac{3R}{8},\qquad
 \abs X^2\geq\frac{\beta_2t}{\mu^2}=8t.
\]
Consequently,
\[
 \rho\geq4,\qquad \lambda\rho=2\abs X\leq\frac{3R}{4},\qquad
 2\lambda^2=4t\leq\frac{3T}{4}.
\]
Thus $Q(\rho,2)$ is mapped into
$B_{3R/4}(0)\times(0,3T/4)$. The integral vanishing also transfers to $v$:
for every $m$ and all sufficiently small $r$,
\[
 \int_0^{r^2}\int_{B_r}\abs v^2\leq
       \lambda^{-5}C_m(\lambda r)^m.
\]

Let $h(s)=s e^{(1-s)/3}$ and
\[
 a=\frac{\beta\rho^2}{2\log h(3/2)},\qquad
 \omega_a(y,s)=h(s)^{-2a}e^{-\abs y^2/(4s)}.
\]
Since $h(3/2)>1$, $a>0$. Choose smooth cutoffs $\theta_\rho(y)$ and
$\eta(s)$ with $\theta_\rho=1$ on $B_{\rho-1}$ and supported in $B_\rho$,
and $\eta=1$ on $(0,3/2)$ and $\eta=0$ for $s\geq7/4$. Their derivatives
through order two are bounded independently of $\rho\geq4$. Choose
$\chi_\varepsilon\in C^\infty(\R)$ equal to zero for $s\leq\varepsilon$,
equal to one for $s\geq2\varepsilon$, and satisfying
$\abs{\chi_\varepsilon'}\leq C\varepsilon^{-1}$. For sufficiently small
$\varepsilon>0$, set
$z_\varepsilon=\chi_\varepsilon\theta_\rho\eta v$, extending it by zero
outside $Q(\rho,2)$. The weak product rule, obtained by applying the weak
derivative identities to the product of a test function and a smooth cutoff,
shows that $z_\varepsilon$ and its first spatial, second spatial, and time
weak derivatives are in $L^2$ and that it has compact support in
$\R^3\times(0,2)$. The density extension
Lemma~\ref{lem:carleman-sobolev} therefore permits applying
Proposition~\ref{prop:carleman-gauss} to $z_\varepsilon$.

Let
\[
 K=\bigl(B_\rho\times(0,2)\bigr)
       \setminus\bigl(B_{\rho-1}\times(0,3/2)\bigr).
\]
Using the weak product rule twice and the inequality for $Lv$, one obtains
\[
 \abs{\Lop z_\varepsilon}^2\leq
 Cc_1^2\lambda^2(\abs{z_\varepsilon}^2+
                         \abs{\nabla z_\varepsilon}^2)
 +C(c_1)\mathbf 1_K(\abs v^2+\abs{\nabla v}^2)
 +C\varepsilon^{-2}\mathbf 1_{(0,2\varepsilon)}\abs v^2.
\]
Indeed, the derivatives of $\theta_\rho\eta$ are supported in $K$;
in the leading term, replace
$\chi_\varepsilon\theta_\rho\eta\nabla v$ by
$\nabla z_\varepsilon$ minus the cutoff-gradient term, which is again
supported in $K$. The last term contains the derivative of
$\chi_\varepsilon$.

For $\rho\geq4$, $a\geq a_*:=8\beta/\log h(3/2)>0$. The left side of the
Carleman inequality controls
\[
 \kappa\int\omega_a(\abs{z_\varepsilon}^2+
                         \abs{\nabla z_\varepsilon}^2),
 \qquad \kappa=\min\{1,a_*/2\},
\]
because $s<2$ on the support. Let $C$ be the absolute coefficient of
$c_1^2\lambda^2$ in the preceding bound and let $c_0$ be the constant in
Proposition~\ref{prop:carleman-gauss}. Since $\lambda^2=2t\leq2\gamma$,
choose
\[
 0<\gamma(c_1)<\min\left\{\frac{3}{16},
                  \frac{\kappa}{4Cc_0c_1^2}\right\}.
\]
Then the first term on the right is absorbed by the Carleman left side. The
result is
\begin{align}
 \int\omega_a\left(\frac{a}{s}\abs{z_\varepsilon}^2+
                              \abs{\nabla z_\varepsilon}^2\right)
 &\leq C(c_1)\int_K\omega_a(\abs v^2+\abs{\nabla v}^2)\notag\\
 &\quad+C\varepsilon^{-2}
       \int_0^{2\varepsilon}\int_{B_\rho}\omega_a\abs v^2.
 \label{eq:uc-cutoff-carleman}
\end{align}

We verify that the last term tends to zero. Fix an integer $m>4a+4$ and a
radius $r_*>0$ small enough that the scaled form of
\eqref{eq:uc-integral-vanishing} with this exponent $m$ holds whenever
$0<r<r_*$. For $0<s<1$, $h(s)\geq c s$ with an absolute $c>0$. Put
$\tau_j=2^{1-j}\varepsilon$ and
$I_j=(\tau_j/2,\tau_j)$, for $j=0,1,\ldots$; these intervals cover
$(0,2\varepsilon)$ up to endpoints. Partition space into
$S_{j,0}=B_{\sqrt{\tau_j}}$ and, for $k\geq1$,
\[
 S_{j,k}=\{y:2^{k-1}\sqrt{\tau_j}\leq\abs y<
                                  2^k\sqrt{\tau_j}\}.
\]
If $2^k\sqrt{\tau_j}< r_*$, the set
$S_{j,k}\times I_j$ is contained in
$B_{2^k\sqrt{\tau_j}}\times]0,(2^k\sqrt{\tau_j})^2[$. On this shell the
Gaussian is at most $e^{-4^k/16}$ for $k\geq1$ (and at most one for $k=0$),
so its weighted integral is bounded by
\[
 C_m\tau_j^{m/2-2a}2^{km}e^{-4^k/16},
\]
with the exponential factor omitted at $k=0$. Summing these shells gives at
most $C_{m,a}\tau_j^{m/2-2a}$. On the remaining shells,
$\abs y\geq r_*/2$; their total contribution is at most
\[
 C\norm{v}_{L^2(Q(\rho,2))}^2\tau_j^{-2a}
          e^{-r_*^2/(16\tau_j)}.
\]
Consequently the final term of~\eqref{eq:uc-cutoff-carleman} is bounded by
\[
 C_{m,a}\varepsilon^{-2}\sum_{j\geq0}\tau_j^{m/2-2a}
 +C\varepsilon^{-2}\norm{v}_{L^2(Q(\rho,2))}^2
       \sum_{j\geq0}\tau_j^{-2a}e^{-r_*^2/(16\tau_j)}.
\]
The first sum is
$O(\varepsilon^{m/2-2a-2})=o(1)$. For every $q>0$ the exponential
dominates $\tau^{-q}$ as $\tau\downarrow0$; choose $q>2a+2$. The second
sum after multiplication by $\varepsilon^{-2}$ is then
$O(\varepsilon^{q-2a-2})=o(1)$. This proves the required limit of the
initial-time error.

For the remaining cutoff term, we claim
\begin{equation}\label{eq:uc-collar-weight}
 \sup_K\omega_a\leq e^{-\beta\rho^2}.
\end{equation}
On the part where $s\geq3/2$, monotonicity of $h$ on $[3/2,2]$ gives
$h(s)^{-2a}\leq h(3/2)^{-2a}=e^{-\beta\rho^2}$. On the spatial collar
$\abs y\geq\rho-1$, $s<3/2$, and
$(\rho-1)^2\geq9\rho^2/16$. The function
\[
 g(s)=h(s)^{-2a}e^{-\rho^2/(16s)}
\]
is increasing on $(0,2)$: differentiating its logarithm and substituting
$a=\beta\rho^2/(2\log h(3/2))$ reduces the sign to
$1/16+(\beta/\log h(3/2))(-s+s^2/3)\geq0$. Since $s-s^2/3\leq3/4$ for
all $s$, with equality at $s=3/2$, this follows from
$\beta\leq\log h(3/2)/12$; here $\log h(3/2)=\log(3/2)-1/6>0.23$. Moreover,
$g(2)\leq h(2)^{-2a}\leq h(3/2)^{-2a}=e^{-\beta\rho^2}$, since $h(2)>h(3/2)$.
On the spatial collar, $\abs y^2\geq9\rho^2/16$ gives
$e^{-\abs y^2/(4s)}\leq e^{-9\rho^2/(64s)}\leq e^{-\rho^2/(16s)}$, so
$\omega_a(y,s)\leq g(s)\leq g(2)\leq e^{-\beta\rho^2}$. This
proves~\eqref{eq:uc-collar-weight}. Thus the weighted collar integral is at
most
\[
 e^{-\beta\rho^2}\int_{Q(\rho,2)}(\abs v^2+\abs{\nabla v}^2).
\]
Fatou's lemma applied to the nonnegative Carleman integrand in
\eqref{eq:uc-cutoff-carleman}, together with the vanishing initial-time term,
now yields
\begin{equation}\label{eq:uc-interior-weight}
 D:=\int_0^{3/2}\int_{B_{\rho-1}}\omega_a
                  (\abs v^2+\abs{\nabla v}^2)
 \leq C(c_1)e^{-\beta\rho^2}
       \int_{Q(\rho,2)}(\abs v^2+\abs{\nabla v}^2).
\end{equation}

The change of variables gives
\[
 \int_{Q(\rho,2)}(\abs v^2+\abs{\nabla v}^2)
 =\lambda^{-5}\int_{\lambda Q(\rho,2)}\abs w^2
  +\lambda^{-3}\int_{\lambda Q(\rho,2)}\abs{\nabla w}^2
 \leq 2\lambda^{-5}\mathcal N_0(R,T;w),
\]
where $\lambda Q(\rho,2)$ denotes
$B_{\lambda\rho}\times(0,2\lambda^2)$ and we used
$\lambda^2\leq2T$. Put $\mu=\sqrt{2\beta}$ and
$\Omega=B_{\mu X/\lambda}\times(1/2,1)$. The condition $\rho\geq4$
implies $B_{\mu X/\lambda}(1)\subset B_{\rho-1}$, since
$\mu\leq1$ and $\abs{X}/\lambda=\rho/2$. On $\Omega$,
$h(s)\leq1$ and
$\abs y^2\leq2\mu^2\abs X^2/\lambda^2+2$, so
\[
 \int_\Omega\abs v^2\leq
 e^{\mu^2\abs X^2/\lambda^2+1}D.
\]
Rescaling this box back to the original variables, using
$\rho=2\abs X/\lambda$ and $\lambda^2=2t$, gives
\[
 \int_t^{2t}\int_{B_\lambda(\mu X)}\abs{w(y,s)}^2\,dy\,ds
 \leq C(c_1)\mathcal N_0(R,T;w)
      e^{-\beta\abs X^2/t}.
\]
By construction $x=\mu X$, so the target ball is $B_\lambda(x)$ and
$\beta\abs X^2/t=\abs x^2/(2t)$. This is
\eqref{eq:uc-gaussian-box}. The average statement follows because the box
has measure $2^{3/2}|B_1|t^{5/2}$.
\end{proof}

\begin{theorem}[Unique continuation across a spatial boundary]\label{thm:uc}
Let $R,T,c_1>0$ and let $w\colon B_R\times[0,T)\to\R^3$ be continuous.
Suppose that $w$ has space-time weak derivatives $\nabla w$, $\nabla^2w$,
and $\dt w$ on $Q(R,T)=B_R\times(0,T)$ and
\begin{equation}\label{eq:uc-w212}
 \abs{w}+\abs{\nabla w}+\abs{\nabla^2w}+\abs{\dt w}\in L^2(Q(R,T)),
 \qquad \abs{\Lop w}\leq c_1(\abs w+\abs{\nabla w})\quad\text{a.e.}
\end{equation}
Assume that for every $k\in\N$ there is $C_k>0$ such that
\begin{equation}\label{eq:uc-pointwise-vanishing}
 \abs{w(x,t)}\leq C_k(\abs x+\sqrt t)^k
       \qquad ((x,t)\in Q(R,T)).
\end{equation}
Then $w(x,0)=0$ for every $x\in B_R$.
\end{theorem}

\begin{proof}
Set $T_0=\min\{T,1\}$ and restrict $w$ to $B_R\times[0,T_0)$. For
$0<r<\min\{1,R,\sqrt{T_0}\}$, (\ref{eq:uc-pointwise-vanishing}) gives
$|w(y,s)|^2\leq C_k^2(2r)^{2k}$ on
$B_r\times(0,r^2)$. This cylinder has measure $|B_1|r^5$, so
\[
 \int_0^{r^2}\int_{B_r}|w|^2\leq |B_1|C_k^2 2^{2k}r^{2k+5}.
\]
Given any integer $m$, choose $k$ with $2k+5\geq m$; since $r<1$, this
proves the integral vanishing hypothesis of Lemma~\ref{lem:uc-gaussian} at
the origin.

Apply that lemma on $B_R\times(0,T_0)$. If $x\ne0$ and $|x|<\beta_1R$,
then for all sufficiently small $t>0$ its conditions hold. The average of
$|w|^2$ over $B_{\sqrt{2t}}(x)\times(t,2t)$ is bounded by
\[
 C(c_1)\mathcal N_0(R,T_0;w)t^{-5/2}
       e^{-|x|^2/(2t)},
\]
which tends to zero. By continuity of $w$ at $(x,0)$, these averages tend to
$|w(x,0)|^2$; hence $w(x,0)=0$. At $x=0$, (\ref{eq:uc-pointwise-vanishing})
and continuity give $w(0,0)=0$.

We record why the same Gaussian estimate supplies the integral vanishing
needed to move the center. Consider any application of
Lemma~\ref{lem:uc-gaussian} centered at $x_0$ on a ball of radius $R_0$,
and let $y_0$ satisfy $0<d:=|y_0-x_0|<\beta_1R_0$. Choose $r>0$ small
enough that $B(y_0,3r)$ lies in $B(x_0,\beta_1R_0)$,
$|y-x_0|\geq d/2$ for $y\in B(y_0,3r)$,
$\beta_2r^2<d^2/4$, and $r^2<\gamma T_0$ for the time length $T_0$ of
this application. For $j\geq0$, set $\tau_j=2^{-j-1}r^2$ and
$q_j=\tfrac12\sqrt{2\tau_j}$. Since $q_j\leq r/2$, choose a maximal family
of pairwise disjoint balls of radius $q_j/2$ with centers
$x_{j\ell}\in B(y_0,r)$. These balls lie in $B(y_0,r+q_j/2)$, so comparison
of volumes bounds their number by
$C(r/q_j)^3=C r^3\tau_j^{-3/2}$. Maximality implies that the balls
$B_{q_j}(x_{j\ell})$ cover $B(y_0,r)$, and their centers lie in
$B(y_0,2r)$.
The corresponding boxes
$B_{\sqrt{2\tau_j}}(x_{j\ell})\times(\tau_j,2\tau_j)$ cover
$B(y_0,r)\times(0,r^2)$ as $j$ varies. Every such box is eligible for
the Gaussian estimate, and its integral is at most
$C\mathcal N_{x_0}(R_0,T_0;w)e^{-c d^2/\tau_j}$, with an absolute $c>0$.
Therefore
\[
 \int_0^{r^2}\int_{B_r(y_0)}|w|^2
 \leq C\mathcal N_{x_0}(R_0,T_0;w)r^3
       \sum_{j\geq0}\tau_j^{-3/2}e^{-c d^2/\tau_j}.
\]
For every $p>0$, $e^{-cd^2/\tau}\leq C(p,d)\tau^p$ for sufficiently
small $\tau$. Taking $p$ arbitrarily large shows that the last expression
is $O(r^m)$ for every $m$. Thus $w$ has integral vanishing of infinite
order at $(y_0,0)$. If $y_0=x_0$, that property is the hypothesis of the
application itself. This argument uses the $L^2$ box estimate only and
does not assume pointwise spatial regularity at positive times.

It remains to reach every point of $B_R$. Define
\[
 r_0=\beta_1R,\qquad r_{j+1}=r_j+\beta_1(R-r_j).
\]
We have proved both $w(y,0)=0$ and integral vanishing at every
$|y|<r_0$. Suppose these two conclusions hold for every $|y|<r_j$.
Given $y$ with $r_j\leq|y|<r_{j+1}$, choose $s<r_j$ sufficiently close
to $r_j$ that $|y|-s<\beta_1(R-s)$, and put
$x_0=s y/|y|$. Then $|x_0|=s<r_j$ and
$|y-x_0|<\beta_1(R-|x_0|)$. The restriction of $w$ to
$B(x_0,R-|x_0|)\times[0,T_0)$ is continuous there. Its weak derivative
fields restrict to this smaller cylinder, all four $L^2$ norms remain finite,
and the differential inequality holds a.e. by restriction. Thus the smaller
open cylinder satisfies the regularity and inequality hypotheses of
Lemma~\ref{lem:uc-gaussian}; the ball lies in $B_R$, and integral vanishing
at $x_0$ is the induction hypothesis.
Applying the lemma with this center gives the Gaussian estimate at $y$ for
all sufficiently small $t$, since $y\ne x_0$. Dividing its box integral by
the box measure gives averages tending to zero, and continuity yields
$w(y,0)=0$. The covering argument above gives integral vanishing at $y$,
which supplies the induction hypothesis for the next step.

Finally $r_j=R-(1-\beta_1)^j(R-r_0)$ increases to $R$. For each
$y\in B_R$, some finite $j$ has $|y|<r_j$, proving the conclusion.
\end{proof}

\begin{remark}[Departure from \cite{ESS2003}]
The class in Theorem~\ref{thm:uc} uses the space-time weak-derivative
$W^{2,1}_2$ convention and a continuous representative up to $t=0$, both fixed
in the statement of the theorem; this makes the pointwise growth assumption and
the value $w(x,0)$ unambiguous. ESS's Lemma~4.3~\cite[pp.~230--233]{ESS2003} estimates a pointwise value
through (4.18). Here the Carleman estimate instead yields an $L^2$ bound on
parabolic boxes, with the cutoff collar controlled by the $L^2$ norms in
(4.1). The initial-time cutoff limit is justified by the dyadic shell sum
above, and the final passage from the inner ball to $B_R$ is the explicit
radius iteration in the proof.
\end{remark}

\section{Backward uniqueness on a half-space}\label{sec:bu}

Backward uniqueness uses two different decays. The first Carleman estimate turns
small Gaussian growth into decay away from the boundary plane. The anisotropic
estimate then makes the remaining cutoff errors vanish as its parameter tends
to infinity. We work throughout with the space-time weak derivatives defined
in Section~\ref{sec:notation}; in particular, all differential inequalities
below hold almost everywhere.

\pagebreak[4]

\begin{lemma}[Gaussian averages]\label{lem:bu-gaussian}
Set $A_0=10^{-12}$ and $\beta=10^{-6}$. For every $c_1>0$ there is
$\gamma=\gamma(c_1)\in(0,1/12)$ with the following property. Let
$A\in[0,A_0]$, and suppose that $w$ is continuous on
$\Hplus\times[0,1)$, has space-time weak derivatives
$Dw,D^2w,D_tw\in L^2$ on every bounded subset of $\Qp$, and satisfies
\begin{align*}
  |D_tw+\Delta w|&\le c_1(|Dw|+|w|) &&\text{a.e. in }\Qp,\\
  w(x,0)&=0 &&\text{for }x\in\Hplus,\\
  |w(x,t)|&\le e^{A|x|^2} &&\text{for }(x,t)\in\Qp.
\end{align*}
Writing $\bar A=\max\{A,A_0/2\}$, there is
$C=C(c_1,A)>0$ such that, for every $x_3>2$ and $0<t<\gamma$,
\begin{equation}\label{eq:bu-gaussian-average}
  \frac{1}{t^{5/2}}
  \int_t^{5t/2}\int_{\Ball{\sqrt{3t}}{x}} |w(y,s)|^2\,dy\,ds
  \le C e^{8\bar A|x|^2}e^{-\beta x_3^2/(12t)}\,.
\end{equation}
\end{lemma}

\begin{proof}
Fix $x_3>2$ and $0<t<\gamma$, and write
$h(s)=s e^{(1-s)/3}$. Put $\lambda=\sqrt{3t}$,
$\rho=(x_3-1)/\lambda$, and $\sigma=1/6$. The rescaled
field
\[
  v(y,s)=w(x+\lambda y,\lambda^2(s-\sigma))
\]
is defined for $y\in B_\rho$ and $\sigma<s<3$. Its weak derivatives obey
\[
  |D_sv+\Delta_yv|\le c_1\lambda(|D_yv|+|v|),\qquad
  |v(y,s)|\le e^{2A|x|^2+2A\lambda^2|y|^2}\,.
\]
The first inequality follows from the affine change of variables and
$\lambda<1$; the second follows from
$|x+\lambda y|^2\le2|x|^2+2\lambda^2|y|^2$.

The field $v$ already uses the Carleman time coordinate: its source time is
$\lambda^2(s-\sigma)$, where $\sigma=1/6$. Multiply it by the mollified
spatial cutoff, equal to one on $B_{\rho/2}$ and supported in
$B_{3\rho/4}$, the final-time cutoff evaluated at $s-\sigma$ (one through
$3/2$ and zero after $7/4$), and the initial-time cutoff evaluated at
$s-\sigma$ (zero through $\varepsilon$ and one from $2\varepsilon$).
Extend their product $V_\varepsilon$ by zero. Its spatial transition is
contained in $13\rho/20\leq|y|\leq3\rho/4$; the late transition is
$\sigma+3/2\leq s\leq\sigma+7/4$; and the initial transition is
$\sigma+\varepsilon\leq s\leq\sigma+2\varepsilon$. The weak product rule
gives its first spatial, second spatial, and time weak derivatives in
$L^2$, with compact support in $\mathbb R^3\times(0,2)$. Thus
Lemma~\ref{lem:carleman-sobolev} permits applying
Proposition~\ref{prop:carleman-gauss} to $V_\varepsilon$. Applying the
Gaussian Carleman estimate, set
\[
 I_\varepsilon=\int_0^2\int_{\R^3}h^{-2a}(s)e^{-|y|^2/(4s)}
 \left(\frac{a}{s}|V_\varepsilon|^2+|D_yV_\varepsilon|^2\right)dy\,ds.
\]
Let $S_\rho=\{13\rho/20\leq|y|\leq3\rho/4\}$ and
$L_\rho=B_{3\rho/4}$. The localized product-rule estimate has a
spatial error on $S_\rho\times(\sigma,\sigma+7/4)$ involving
$|v|+|D_yv|$, a late-time error on
$L_\rho\times[\sigma+3/2,\sigma+7/4]$ involving $|v|$, and an
initial-strip error on
$B_\rho\times[\sigma+\varepsilon,\sigma+2\varepsilon]$ involving
$\varepsilon^{-1}|v|$. Applying the Gaussian Carleman estimate to
$V_\varepsilon$ and squaring this bound gives
\begin{align}
 I_\varepsilon
 &\leq C_Gc_1^2\lambda^2 I_\varepsilon\notag\\
 &\quad+C(c_1)\int_{S_\rho\times(\sigma,\sigma+7/4)}
       \omega_a(|v|^2+|D_yv|^2)
   +C(c_1)\int_{L_\rho\times(\sigma+3/2,\sigma+7/4)}\omega_a|v|^2\notag\\
 &\quad+C\varepsilon^{-2}
       \int_{\sigma+\varepsilon}^{\sigma+2\varepsilon}
             \int_{B_\rho}\omega_a|v|^2.
\end{align}
Here $C_G$ is absolute and $\omega_a=h(s)^{-2a}e^{-|y|^2/(4s)}$.


The initial-strip error is controlled without extending $w$ by zero. Since
$v(y,s)=w(x+\lambda y,\lambda^2(s-\sigma))$,
Lemma~\ref{lem:ftc-small-time} gives
\[
 \varepsilon^{-2}\int_{\sigma+\varepsilon}^{\sigma+2\varepsilon}
       \int_{B_\rho}|v(y,s)|^2\,dy\,ds
 \leq 2\lambda^{-1}\int_0^{2\lambda^2\varepsilon}
       \int_{B(x,\lambda\rho)}|D_tw(z,r)|^2\,dz\,dr\longrightarrow0.
\]
For fixed $a$, $\omega_a$ is bounded on
$B_\rho\times[\sigma,\sigma+1/6]$, so the weighted initial error also
tends to zero. We use this only to choose one sufficiently small
$\varepsilon$; the proof below does not pass the weighted left side to an
$\varepsilon\downarrow0$ limit.

Choose $\gamma(c_1)$ so that
\[
 \gamma\le\frac1{24},\qquad
 \gamma\le\frac{\beta}{3H},\qquad
 3C_Gc_1^2\gamma\le\frac12,
 \qquad H=\log h(3/2)>\frac16.
\]
For $0<t<\gamma$, set $a=\beta\rho^2/H$. Then
$\rho^2>1/(3t)>H/\beta$, so $a>1$ and $\rho>4$; also
$\lambda^2=3t\le3\gamma$ absorbs the first term in the estimate for
$I_\varepsilon$.

For the chosen $a=\beta\rho^2/H$, put
$T_\rho=e^{8\bar A|x|^2}e^{-2\beta\rho^2}$. The spatial-shell errors lie
on $S_\rho=\{13\rho/20\leq|y|\leq3\rho/4\}$; the late-time mass
error lies in $B_{3\rho/4}$ over Carleman times
$[\sigma+3/2,\sigma+7/4]$, where the unshifted rescaled time
$s-\sigma$ lies in $[3/2,7/4]$. Take $r=(16\sqrt a)^{-1}$. A finite
Besicovitch cover of $S_\rho$ by balls of radius $r$ has spatial overlap
at most $N$, and its doubled balls lie in $B_\rho$. For Caccioppoli in
the rescaled source time, set $a_j=\sigma+jr^2/2$, use inner intervals
$(a_j,a_j+r^2)$, and enlarge them to $J_j=(a_j,a_j+4r^2)$. These cover
$(\sigma,\sigma+7/4)$ up to a null set, remain below $2$, and have
overlap at most $8$; the resulting outer cylinders therefore have
overlap at most $8N$. The Gaussian weight varies by a bounded factor on
each outer cylinder because $\rho r=\sqrt{H/\beta}/16$ is fixed.
Caccioppoli on each inner cylinder, including its initial-time version
at $s=\sigma$, controls the shell gradient by outer-cylinder radial
mass. The radial mass tail bound, the late-time bound, and the cutoff
derivative bounds give shell and late errors bounded by
$C(c_1,A)(1+a)(1+\rho)^3T_\rho$. Choose $0<\varepsilon<1/12$ small
enough that the weighted initial error in the cutoff estimate above
is at most $T_\rho$; the initial-strip limit above permits this choice.
Absorbing the term proportional to $c_1^2\lambda^2$ then yields the
direct one-cutoff estimate
\begin{equation}\label{eq:bu-gaussian-collar}
 I_\varepsilon\leq C(c_1,A)(1+a)(1+\rho)^3
 e^{8\bar A|x|^2}e^{-2\beta\rho^2}.
\end{equation}
The spatial cutoff equals one on $B_1$, and for this choice of
$\varepsilon$ the time cutoff equals one on the averaging core
$B_1\times(1/2,1)$. Since $\omega_a\geq e^{-3/2}$ there, the core mass
is directly controlled:
\[
 D:=\int_{1/2}^1\int_{B_1}|v(y,s)|^2\,dy\,ds\leq e^{3/2}I_\varepsilon.
\]
Finally, $a=\beta\rho^2/H$ implies that the polynomial factor is absorbed by the Gaussian tail:
$ (1+a)(1+\rho)^3e^{-2\beta\rho^2}\leq C e^{-\beta\rho^2}$. Thus $D\leq C(c_1,A)e^{8\bar A|x|^2}e^{-\beta\rho^2}.$

Changing variables back on $B_1\times(1/2,1)$ gives the left side of
\eqref{eq:bu-gaussian-average} up to the absolute factor $3^{5/2}$.
Finally, $\rho^2=(x_3-1)^2/(3t)\ge x_3^2/(12t)$ for $x_3>2$.
\end{proof}

\begin{remark}[Departure from \cite{ESS2003}]
The initial time is not extended by zero. A cutoff above $s=1/6+\varepsilon$
and Lemma~\ref{lem:ftc-small-time} replace that extension; its differentiated
error tends to zero by the displayed $L^2$ estimate. The proof uses
Lemma~\ref{lem:caccioppoli} for the collar gradients and proves an average
estimate, so it does not use the pointwise parabolic regularity bounds or
estimate (5.17) of ESS~\cite[§5, pp.~234--236]{ESS2003}.
The explicit choice $A_0=10^{-12}$, $\beta=10^{-6}$ also supplies the
uniform positive Gaussian exponent needed when the input value $A$ is zero;
it satisfies $A_0<1/32$, $8A_0<1/256$, $16A_0/24<1/48$,
$16A_0/24<\beta/48$, $\beta<1/256$, and $\beta<H/96$, where
$H=\log(3/2)-1/6>1/6$.
\end{remark}

\begin{lemma}[Vanishing on a short time interval]\label{lem:bu-small-time}
Fix $c_1>0$ and let $A\in[0,A_0]$. Suppose $w$ satisfies the continuity,
weak-derivative, differential-inequality and zero-trace hypotheses of
Lemma~\ref{lem:bu-gaussian}, as well as $|w(x,t)|\le e^{A|x|^2}$ on $\Qp$.
There is $\gamma_1=\gamma_1(c_1)\in(0,\gamma(c_1)/2]$ such that
\[
  w(x,t)=0\qquad\text{for every }x\in\Hplus\text{ and }0<t<\gamma_1\,.
\]
\end{lemma}

\begin{proof}
Fix $\alpha=3/4$, and let $c_*(\alpha)$ and $a_0(\alpha)$ be the constants
in Proposition~\ref{prop:carleman-halfspace}. Put
$\lambda=\sqrt{2\gamma_1}$ and, for $y_3>0$ and $1/2<s<1$, set
\[
 v(y,s)=w(\lambda y,\lambda^2(s-1/2)),\qquad
 F(y_3,s)=\frac{(1-s)y_3^{2\alpha}}{s^\alpha}.
\]
Then $|D_sv+\Delta_yv|\le c_1\lambda(|D_yv|+|v|)$ a.e. and
$v(\cdot,s)\to0$ locally in $L^2$ as $s\downarrow1/2$, by the assumed
continuous zero trace. Choose
$y_-=3/\lambda+1$, $y_+=3/\lambda+3/2$, and
$y_0=3/\lambda+2$. Define
$B=2F(y_0,1/2)$ and $\Phi(y_3,s)=F(y_3,s)-B$. Let $\psi_1$ be smooth,
zero on $(-\infty,y_-]$ and one on $[y_+,\infty)$. Let $\psi_2$ be smooth,
zero on $(-\infty,-3/4]$ and one on $[-2/3,\infty)$. Set
$\eta(y_3,s)=\psi_1(y_3)\psi_2(\Phi(y_3,s)/B)$.

The derivatives of $\eta$ are supported where $\Phi\le-D/2$, with
$D=-2\Phi(y_+,1/2)>0$. Indeed, $F$ increases in $y_3$ and decreases in
$s$, so the support of $\psi_1'$ has $\Phi\le\Phi(y_+,1/2)=-D/2$.
On the support of $\psi_2'$, $\Phi\le-2B/3$. Since $\lambda<1/2$,
\[
 \frac{D}{2B}=1-\frac12\left(\frac{y_+}{y_0}\right)^{3/2}
 <1-\frac{45}{128}<\frac23,
\]
which gives the same bound there. Also $\Phi>0$ implies $y_3>y_+$, so
$\eta=1$ on that region.

For $R>1$, choose a smooth spatial cutoff $\chi_R$ equal to one on $B_R$,
zero outside $B_{2R}$, and satisfying
$|D\chi_R|\le C/R$, $|D^2\chi_R|\le C/R^2$. For $\varepsilon>0$, choose
$\vartheta_\varepsilon$ equal to zero for $s\le1/2+\varepsilon$, equal to
one for $s\ge1/2+2\varepsilon$, and with
$|\vartheta_\varepsilon'|\le C/\varepsilon$. The field
$V_{R,\varepsilon}=\chi_R\eta\vartheta_\varepsilon v$ has compact support
in $(\R^3_++e_3)\times(0,1)$: its spatial support stays above $y_3>1$,
and on the bounded set $B_{2R}$ the condition $\psi_2(\Phi/B)\ne0$
keeps $s$ a positive distance below $1$. It belongs to the Sobolev class of
Lemma~\ref{lem:carleman-sobolev} by the weak product rule.

Apply Proposition~\ref{prop:carleman-halfspace} to $V_{R,\varepsilon}$
with parameter $a>\max\{a_0(3/4),1\}$. The terms containing
$D_sv+\Delta_yv$ are bounded by
$C c_1^2\lambda^2$ times the Carleman left side, plus cutoff errors.
Choose $\gamma_1$ so that
\[
  2C c_*(3/4)c_1^2\gamma_1\le\frac14,
  \qquad \gamma_1\le\frac{\gamma(c_1)}2,
\]
where $C$ is the absolute product-rule constant. These terms are then
absorbed. The errors from $D\eta$ are supported in
\[
\omega=\{y_3>y_-,\ 1/2<s<1,\ \Phi(y_3,s)\le-D/2\}.
\]
Indeed, the $\psi_1'$ term is supported in $(y_-,y_+)$, while each term from
$\psi_2'$ is multiplied by $\psi_1$ and vanishes for $y_3\le y_-$. The
preceding cutoff estimates give $\Phi\le-D/2$ on both supports. After
absorption, these errors give
\begin{align}
 I_a&:=\int_{1/2}^1\int_{\R^3_+}e^{-|y'|^2/(4s)}e^{2a\Phi}
       \big(a|\eta v|^2+s|D_y(\eta v)|^2\big)\,dy\,ds \notag\\
 &\le C(c_1)\int_\omega (1+y_3)^4 e^{-|y'|^2/(4s)}e^{2a\Phi}
       (|v|^2+|D_yv|^2)\,dy\,ds+E_{R,a}+E_{\varepsilon,a}\,.
 \label{eq:bu-halfspace-carleman}
\end{align}
Here $E_{R,a}$ is the error from $D\chi_R,D^2\chi_R$, and
$E_{\varepsilon,a}$ is the error from $\vartheta_\varepsilon'$. The factor
$(1+y_3)^4$ bounds the squares of the first and second derivatives of the
fixed cutoff $\eta$ on $y_3>1$.

The lower-time error vanishes as $\varepsilon\downarrow0$. On the bounded
spatial support $B_{2R}$ the Carleman weight is bounded for fixed $a,R$, and
Lemma~\ref{lem:ftc-small-time} gives
\[
 E_{\varepsilon,a}\le C(a,R,\lambda)
 \varepsilon^{-2}\int_{1/2+\varepsilon}^{1/2+2\varepsilon}
       \int_{B_{2R}}|v|^2\,dy\,ds
 \le C(a,R,\lambda)\int_0^{2\lambda^2\varepsilon}
       \int_{B_{2\lambda R}}|D_tw|^2\,dx\,dt\longrightarrow0.
\]

We next control the noncompact collar by shifted dyadic cells. For $d=2^{-k}$, each layer $d/2<\delta=s-1/2<d$ is partitioned into spatial cubes of side $\sqrt d/128$ and time intervals of length $d/4096$, with 2048 time labels. At a cell center $Y$ with time $\delta_Y\in[d/2,d]$, the Gaussian average estimate gives an average-box mass bounded by
\[
 C(c_1,A,\lambda)d^2
   e^{8A_0\lambda^2|Y|^2}e^{-\beta Y_3^2/(12d)}.
\]
Caccioppoli with $r=\sqrt{\delta_Y}/8$ contributes a factor bounded by $C(c_1)/d$; hence the cell energy has the bound
\[
 \int_{C_{k,m,\ell}}(|v|^2+|D_yv|^2)
 \leq C(c_1,A,\lambda)d
       \exp\!\left(16A_0\lambda^2|Y|^2
                    -\frac{\beta Y_3^2}{48d}\right).
\]
After including $(1+y_3)^4$ and the tangential Gaussian, Gaussian absorption bounds the cell sum by a fixed spatial lattice profile times a dyadic factor $C2^{N(k+1)}e^{-c\beta 2^{k+1}}$ for a fixed integer $N$. The spatial profile is summable, there are only 2048 time labels on each layer, and the dyadic Gaussian series converges. Consequently, after omitting $e^{2a\Phi}$,
\[
 J:=\int_\omega (1+y_3)^4e^{-|y'|^2/(4s)}
       (|v|^2+|D_yv|^2)\,dy\,ds<\infty,
\]
with a bound independent of $a$. The cell-energy power here is $d$: it comes from the $d^2$ Gaussian average and the $d^{-1}$ Caccioppoli factor.

For fixed $a$, the spatial cutoff error tends to zero as $R\to\infty$: its derivatives are supported in $\{|y|>R\}$, and the weighted integral is finite because the negative quadratic Gaussian dominates the positive normal weight $e^{2a y_3^{2\alpha}}$, with $2\alpha<2$. For every fixed $R$, the lower-time error tends to zero as $\varepsilon\downarrow0$ by Lemma~\ref{lem:ftc-small-time}. Thus, for fixed $a$, the cutoff estimate has the form $I_a\leq C(c_1)J e^{-Da}+S_R+E_R(\varepsilon)$, where $S_R\to0$ as $R\to\infty$ and $E_R(\varepsilon)\to0$ for each fixed $R$. Applying the two-parameter limit in that order, first $\varepsilon\downarrow0$ and then $R\to\infty$, gives
\[
 I_a\leq C(c_1,A,\lambda)e^{-Da},
\]
with a constant independent of $a$.

For any compact $K\subset\{y_3>0,\ 1/2<s<1,\ \Phi>0\}$,
$\eta=1$ and $\Phi$ has a positive minimum on $K$. Hence
$a e^{2a\min_K\Phi}\int_K|v|^2\le I_a\to0$ as $a\to\infty$.
Continuity gives $v=0$ on this open curved region.

It remains to fill the half-space at each fixed time $s_0\in(1/2,1)$.
Write
\[
 g(s)=\left(\frac{B s^\alpha}{1-s}\right)^{1/(2\alpha)},
\]
so that $\Phi(y_3,s)>0$ exactly when $y_3>g(s)$. Given any target
$y=(y',y_3)\in\Hplus$, choose $H>\max\{2g(s_0),y_3\}$ and let
$c=(y',H)$, $R=H-y_3/2$. Then $|y-c|=H-y_3<R<H$, so
$y\in\Ball{R}{c}\subset\Hplus$. By continuity of $g$, there is $T>0$ such
that $s_0+T<1$ and $g(s)<H-R/2$ for $s\in[s_0,s_0+T]$. Thus $v$ is zero
on a space-time neighborhood of $(c,s_0)$ lying above the graph.
The translated field $u(z,t)=v(c+z,s_0+t)$ satisfies all hypotheses of
Theorem~\ref{thm:uc} on $\Ball{R}{0}\times(0,T)$: its weak derivatives and
differential inequality are inherited; the cylinder has compact closure in
the half-space; and its all-orders condition holds because $u=0$ when
$|z|+\sqrt{t}<r_0$ for some $r_0>0$. On the rest of the compact cylinder
$u$ is bounded, so the condition holds with $C_k=r_0^{-k}\sup|u|$. Theorem
\ref{thm:uc} therefore gives $v(\cdot,s_0)=0$ on $B_R(c)$, in particular
$v(y,s_0)=0$. As $y$ and $s_0$ were arbitrary, these balls cover the
connected half-space at every $s_0\in(1/2,1)$. No time reversal is used.
Returning to $(x,t)$ proves the lemma.
\end{proof}

\begin{remark}[Departure from \cite{ESS2003}]
The application of Proposition~\ref{prop:carleman-halfspace} to the
noncompact cutoff is justified by first inserting $\chi_R$; its shell terms
vanish by the Gaussian-average estimate and Caccioppoli. A cutoff above
$s=1/2+\varepsilon$, controlled by Lemma~\ref{lem:ftc-small-time}, avoids
extending the field by zero at its initial trace. The source invokes
Theorem~\ref{thm:uc} to fill the curved zero region in one sentence~\cite[§5, pp.~236--238]{ESS2003}; here
the center $c=(y',H)$ and radius $R=H-y_3/2$ are chosen for each target point,
and Theorem~\ref{thm:uc} is applied forward from the time $s_0$. The
modified transition interval for $\psi_2$ ensures every cutoff derivative
lies in $\Phi\le-D/2$.
\end{remark}

\begin{lemma}[Iteration in time]\label{lem:bu-iterate}
Suppose $w$ satisfies the hypotheses of Lemma~\ref{lem:bu-small-time} with
some $A\in[0,A_0]$. Then $w\equiv0$ on $\Qp$.
\end{lemma}

\begin{proof}
Set $\gamma_0=0$ and
$\gamma_{k+1}=\gamma_k+(1-\gamma_k)\gamma_1$. If $w=0$ for
$0<t<\gamma_k$, define
\[
 w^{(k)}(y,s)=w\big(\sqrt{1-\gamma_k}\,y,
                   \gamma_k+(1-\gamma_k)s\big),
 \qquad (y,s)\in\Qp.
\]
The affine change of variables preserves the space-time weak derivative
class and continuity up to $s=0$. Its initial trace is zero: for $k=0$ this
is the given trace, and for $k\ge1$ it follows from continuity and
$w(\cdot,t)=0$ for $t<\gamma_k$. Moreover,
\[
 |D_sw^{(k)}+\Delta_yw^{(k)}|
 \le c_1\big(|D_yw^{(k)}|+|w^{(k)}|\big),\qquad
 |w^{(k)}(y,s)|\le e^{A|y|^2},
\]
because the time and Laplacian factors are both $1-\gamma_k$, while the
gradient factor is $\sqrt{1-\gamma_k}$. The local $L^2$ conditions are
preserved by this nonsingular change of variables. Lemma~\ref{lem:bu-small-time}
applied to $w^{(k)}$ shows $w=0$ for $0<t<\gamma_{k+1}$. Induction gives
$\gamma_k=1-(1-\gamma_1)^k\uparrow1$, and every $t<1$ lies below some
$\gamma_k$.
\end{proof}

\begin{remark}[Departure from \cite{ESS2003}]
The initial trace at each iteration is justified by continuity from the
already vanishing time interval; the rescaling factors and the recurrence
are displayed rather than left implicit in the proof of Lemma 5.4~\cite[§5, p.~238]{ESS2003}.
\end{remark}

\begin{theorem}[Backward uniqueness on a half-space]\label{thm:bu}
Let $c_1>0$, $M\in\R$, and $\Qp=\Hplus\times(0,1)$. Suppose that $w$ is
continuous on $\Hplus\times[0,1)$, has space-time weak derivatives
$Dw,D^2w,D_tw$ on $\Qp$, and satisfies
\begin{align*}
 w(x,0)&=0 &&(x\in\Hplus),\\
 |D_tw+\Delta w|&\le c_1(|Dw|+|w|) &&\text{a.e. on }\Qp,\\
 |w(x,t)|&\le e^{M|x|^2} &&((x,t)\in\Qp),
\end{align*}
with $w,Dw,D^2w,D_tw\in L^2$ on every bounded subset of $\Qp$. Then
$w\equiv0$ on $\Qp$.
\end{theorem}

\begin{proof}
If $M\le A_0$, replace $M$ by $A=\max\{M,0\}$; the growth bound then
implies $|w(x,t)|\le e^{A|x|^2}$ and $A\in[0,A_0]$. Lemma~\ref{lem:bu-iterate}
proves the conclusion. Suppose next that $M>A_0$, and set
$\lambda^2=A_0/(2M)<1/2$. Define $v_0(y,s)=w(\lambda y,\lambda^2s)$.
The weak derivatives transform by the affine chain rule, giving
\[
 |D_sv_0+\Delta_yv_0|
 \le c_1\big(\lambda|D_yv_0|+\lambda^2|v_0|\big)
 \le c_1(|D_yv_0|+|v_0|),\qquad
 |v_0(y,s)|\le e^{(A_0/2)|y|^2}.
\]
The zero trace, continuity, and local $L^2$ conditions are preserved, so
Lemma~\ref{lem:bu-iterate} gives $v_0\equiv0$. Hence
$w(x,t)=0$ for $0<t<\lambda^2$.

To cover the remaining time, set $t_k=k\lambda^2$ while $t_k<1$ and
$\delta_k=\min\{\lambda^2,1-t_k\}$. Once $w(\cdot,t_k)=0$, define
\[
 v_k(y,s)=w(\sqrt{\delta_k}\,y,t_k+\delta_k s),
 \qquad (y,s)\in\Qp.
\]
For $k=0$, the zero trace is the hypothesis; for $k\ge1$, continuity and
the preceding slab conclusion give $w(\cdot,t_k)=0$. The field $v_k$ is
continuous up to $s=0$ and has the required local weak derivatives by the
change of variables. Since $0<\delta_k\le\lambda^2<1/2$,
\[
 |D_sv_k+\Delta_yv_k|
 \le c_1\big(\sqrt{\delta_k}|D_yv_k|+\delta_k|v_k|\big)
 \le c_1(|D_yv_k|+|v_k|),\qquad
 |v_k(y,s)|\le e^{M\delta_k|y|^2}
 \le e^{(A_0/2)|y|^2}.
\]
Thus Lemma~\ref{lem:bu-iterate} applies on every slab and yields
$w=0$ for $t_k<t<t_k+\delta_k$. The construction takes at most
$\lceil1/\lambda^2\rceil$ slabs; the final one has length
$\delta_k=1-t_k<\lambda^2$ if the endpoint is not reached exactly. They
cover $0<t<1$, proving the theorem.
\end{proof}

\begin{remark}[Departure from \cite{ESS2003}]
The final large-growth iteration in the source asks for another application
of the small-growth lemma with exponent $M$, outside its stated range~\cite[§5, p.~238]{ESS2003}. Here
each time slab is spatially and parabolically rescaled by its own length
$\delta_k\le\lambda^2$; this keeps the growth exponent at most $A_0/2$ and
the differential-inequality constant at most $c_1$. The last partial slab
is included explicitly.
\end{remark}

%======================================================================
\part*{Parts II and III. Existence and the associated pressure (proved in the CKN paper)}
%======================================================================

\section{Leray existence and the associated pressure}\label{sec:existence-in-ckn}

The six results of Section~\ref{sec:leray-hopf} that concern the existence of
suitable Leray--Hopf solutions and their pressures are Theorem~\ref{thm:leray}
(existence), Corollary~\ref{cor:ckn-headline} (existence with singular set of
zero one-dimensional parabolic Hausdorff measure), Theorem~\ref{thm:assoc-pressure}
(associated pressure), and their forced versions
Theorem~\ref{thm:leray-forced}, Theorem~\ref{thm:leray-forced-singularSet} and
Theorem~\ref{thm:assoc-pressure-forced}. They are stated in
Section~\ref{sec:leray-hopf} because the theorems of Parts~IV--VI are about
Leray--Hopf solutions, but they do not depend on anything in Parts~I and
IV--VI. They are proved in the last part of the Caffarelli--Kohn--Nirenberg
paper~\cite{CKNLean}, where the statements carry the same labels and the
proofs occupy the sections on the whole-space pressure
(\texttt{sec:riesz-pressure}), the Leray-regularized equations
(\texttt{sec:regularised}), uniform bounds, compactness, the limit
(\texttt{sec:limit}) and forced Leray solutions (\texttt{sec:forced-leray}).
They are formalized in the CKN library as \texttt{CKN.leray\_existence},
\texttt{CKN.leray\_existence\_singularSet}, \texttt{CKN.associatedPressure},
\texttt{CKN.lerayExistenceForced},
\texttt{CKN.lerayExistenceForcedSingularSet} and
\texttt{CKN.associatedPressureForced}.

The proofs of Parts~IV--VI use, besides these statements, the following
notation and facts from the CKN paper. The whole-space pressure is defined in
the CKN paper as follows, and is used here with the same sign convention.
Further facts from the same part (\cite[Lemma~11.1]{CKNLean} on
compactness, \cite[Lemma~11.5]{CKNLean} on restriction, the
energy interpolation \cite[Lemma~8.4]{CKNLean}, the
Riesz duality \cite[Lemma~8.2]{CKNLean}, the regularized solutions
of \cite[Theorem~9.4]{CKNLean} and the limit of
\cite[Proposition~12.1]{CKNLean}) are cited where they are used.

\begin{definition}[Whole-space pressure]\label{def:riesz-pressure}
For $1<r<\infty$, on Schwartz functions let $\Riesz{j}$ denote the Fourier
multiplier $i\xi_j/|\xi|$. The product $\Riesz{i}\Riesz{j}$ has a bounded
extension to $L^r(\R^3)$. For a jointly measurable tensor field $F=(F_{ij})$
whose slices $F_{ij}(\cdot,t)$ belong to $L^r(\R^3)$ for almost every $t$,
define
\begin{equation}\label{eq:riesz-pressure}
P[F](\cdot,t)=\sum_{i,j=1}^3 \Riesz{i}\Riesz{j}F_{ij}(\cdot,t).
\end{equation}
The extensions for different exponents agree on their common domains, so this
defines one pressure operator. Indeed, bounded compactly supported inputs
belong to $L^r\cap L^s\cap L^2$ for any two finite exponents $r,s$; the CKN
extensions agree there with the same negative raw operator. Truncating an
input in $L^r\cap L^s$ by balls and bounded values gives convergence in both
norms, so the two outputs agree as distributions. With this convention
\begin{equation}\label{eq:riesz-pressure-sign}
-\Delta P[F]=\sum_{i,j=1}^3\partial_i\partial_jF_{ij}
\end{equation}
in the spatial distribution sense for almost every $t$. The CKN raw second
derivative operator has the opposite sign; the all-exponent bounded
extension in the CKN convention represents its negative, which is
$\Riesz{i}\Riesz{j}$ on the common $L^r\cap L^2$ domain.

There is a jointly measurable representative of $P[F]$, and for almost every
$t$,
\begin{equation}\label{eq:riesz-slice-bound}
\norm{P[F](\cdot,t)}_{L^r(\R^3)}
\le C_r\sum_{i,j=1}^3\norm{F_{ij}(\cdot,t)}_{L^r(\R^3)}.
\end{equation}
Here and below $C_r$ depends only on $r$ in dimension three. If
$F\in L^r(\R^3\times(0,T))$, then
\begin{equation}\label{eq:riesz-spacetime-bound}
\norm{P[F]}_{L^r(\R^3\times(0,T))}
\le C_r\sum_{i,j=1}^3\norm{F_{ij}}_{L^r(\R^3\times(0,T))}.
\end{equation}
\end{definition}

%======================================================================
\part*{Part IV. Regularity of $L^\infty L^3$ solutions}
%======================================================================

\section{Local regularity near small-energy points}\label{sec:local-regularity}

We use the CKN past cylinder
\[
  \mathcal Q_r(z_0)=B_r(x_0)\times(t_0-r^2,t_0],
  \qquad z_0=(x_0,t_0),
\]
and write $\overline{\mathcal Q_r(z_0)}=\overline B_r(x_0)\times
[t_0-r^2,t_0]$.  The parabolic distance is
\[
  d_{\rm par}((x,t),(y,s))=\max\{|x-y|,|t-s|^{1/2}\}.
\]
All vector norms below are Euclidean.
Throughout this section, ``suitable'' means suitable in the sense of
Definition~\ref{def:sws}, with force $f=0$; the spatial gradient is the
separate datum $Du$ fixed there.

\begin{lemma}[Theorem A on a cylinder of radius $r$]\label{lem:thmA-rescaled}
There are constants $\varepsilon_0,\gamma_0,C_4$ with
$0<\varepsilon_0$, $0<\gamma_0\le 2/3$, and $0\le C_4$, supplied by
CKN Theorem A at $q=3$, with the following property.  Let $D=\Omega\times I$ be a CKN
solution domain and let $(u,Du,p,0)$ be suitable there with exponent $q=3$.
Suppose $z_0=(x_0,t_0)\in D$, $r>0$,
$\overline{\mathcal Q_r(z_0)}\subset D$, and
\begin{equation}\label{eq:thmA-rescaled-smallness}
 r^{-2}\iint_{\mathcal Q_r(z_0)}
       \bigl(|u|^3+|p|^{3/2}\bigr)\,dx\,dt<\varepsilon_0.
\end{equation}
Then there is a representative $u_{z_0,r}$ of $u$ on
$\mathcal Q_{r/2}(z_0)$ which extends to its closure and satisfies
\begin{align*}
 \sup_{\overline{\mathcal Q_{r/2}(z_0)}}|u_{z_0,r}|&\le C_4r^{-1},\\
 |u_{z_0,r}(z)-u_{z_0,r}(z')|&\le
 C_4r^{-1-\gamma_0}d_{\rm par}(z,z')^{\gamma_0}
 \quad(z,z'\in\overline{\mathcal Q_{r/2}(z_0)}).
\end{align*}
\end{lemma}

\begin{proof}
Set $\Phi(y,s)=(x_0+ry,t_0+r^2s)$ and define
$u_r(y,s)=r u(\Phi(y,s))$, $Du_r(y,s)=r^2Du(\Phi(y,s))$, and
$p_r(y,s)=r^2p(\Phi(y,s))$.  The rescaled domain is
$D_r=\Phi^{-1}(D)$.  The CKN rescaling theorem, together with the
equivalence of its suitable and integrable-suitable predicates, makes
$(u_r,Du_r,p_r,0)$ suitable on $D_r$.  The affine map $\Phi$ sends the
closed unit past cylinder onto $\overline{\mathcal Q_r(z_0)}$, so the
assumed domain inclusion gives the closed-cylinder hypothesis of Theorem A.
The Jacobian of $\Phi$ is $r^5$, and hence
\[
 \iint_{\mathcal Q_1(0)}|u_r|^3
   =r^{-2}\iint_{\mathcal Q_r(z_0)}|u|^3,\qquad
 \iint_{\mathcal Q_1(0)}|p_r|^{3/2}
   =r^{-2}\iint_{\mathcal Q_r(z_0)}|p|^{3/2}.
\]
The strict real-valued bound in \eqref{eq:thmA-rescaled-smallness} implies
the non-strict extended-integral bound required by Theorem A; the force term
vanishes.  That theorem gives a representative $w$ on the closed half
unit-cylinder with supremum at most $C_4$ and parabolic Hölder seminorm at
most $C_4$.  Define
$u_{z_0,r}(z)=r^{-1}w(\Phi^{-1}(z))$.  Since $\Phi$ preserves null sets,
this is a representative of $u$ on the half-cylinder.  Moreover,
$d_{\rm par}(\Phi^{-1}z,\Phi^{-1}z')=r^{-1}d_{\rm par}(z,z')$.
The two stated bounds follow.
\end{proof}


\begin{lemma}[Top-boundary $\varepsilon$-regularity]\label{lem:thmA-top}
Let $\rho>r>0$, $z_0=(x_0,t_0)$, and suppose $(u,Du,p,0)$ is suitable
with exponent $q=3$ on
$B_\rho(x_0)\times(t_0-\rho^2,t_0)$.  If
\begin{equation}\label{eq:thmA-top-smallness}
 r^{-2}\iint_{B_r(x_0)\times(t_0-r^2,t_0)}
       (|u|^3+|p|^{3/2})\,dx\,dt<\varepsilon_0/8,
\end{equation}
then, with $V=B_{r/2}(x_0)\times(t_0-r^2/4,t_0)$, $u$ has a
representative on $V$ which extends to the closure and obeys
\begin{align*}
 \sup_{\overline V}|u_{\rm top}|&\le 2C_4/r,\\
 |u_{\rm top}(z)-u_{\rm top}(z')|
 &\le 4\cdot 8^{\gamma_0}C_4r^{-1-\gamma_0}
          d_{\rm par}(z,z')^{\gamma_0}
 \quad(z,z'\in\overline V).
\end{align*}
\end{lemma}

\begin{proof}
For $x\in B_{r/2}(x_0)$ and $0<s<r^2/4$, set
$z_{x,s}=(x,t_0-s)$ and apply Lemma~\ref{lem:thmA-rescaled} with radius
$r'=r/2$.  The closed full cylinder
$\overline{\mathcal Q_{r/2}(z_{x,s})}$ lies in
$B_r(x_0)\times(t_0-r^2,t_0)$: spatially,
$\overline B_{r/2}(x)\subset B_r(x_0)$, and its time interval is
$[t_0-s-r^2/4,t_0-s]$, strictly between $t_0-r^2$ and $t_0$.
Consequently
\[
 (r/2)^{-2}\iint_{\mathcal Q_{r/2}(z_{x,s})}
       (|u|^3+|p|^{3/2})
 \le 4r^{-2}\iint_{B_r(x_0)\times(t_0-r^2,t_0)}
       (|u|^3+|p|^{3/2})<\varepsilon_0/2.
\]
The smallness and domain hypotheses of that lemma hold uniformly in $x,s$.
It gives a representative on the closed half-cylinder
$\overline B_{r/4}(x)\times[t_0-s-r^2/16,t_0-s]$, with the displayed
bounds.

These open interiors, as $x$ and $s$ vary, cover
$V=B_{r/2}(x_0)\times(t_0-r^2/4,t_0)$.  Indeed, for $(y,t)\in V$, put
$a=t_0-t$ and choose
$s=a-\min\{a/2,r^2/32\}$.  Then $0<s<r^2/4$ and
$0<a-s<r^2/16$, so $(y,t)$ lies strictly inside the half-cylinder centered
at $(y,t_0-s)$.  The local representatives agree on every nonempty
overlap of their open interiors: they are continuous there and equal to
$u$ almost everywhere there, while a nonempty open overlap has positive
spacetime measure.  A countable subcover of $V$ shows that the glued
representative equals $u$ almost everywhere on $V$.

The local estimates have the same exponent and the same constants.  If
$z,z'\in V$ and $d_{\rm par}(z,z')<r/8$, choose the spatial center to be
that of the later point and choose a time center strictly between the later
point and $t_0$, sufficiently close to the later point.  More explicitly,
write $a_{\rm late}<a_{\rm early}$ for their time deficits from $t_0$ and
take
$c=a_{\rm late}-\min\{a_{\rm late}/2,r^2/64\}>0$.
Then both deficits lie strictly between $c$ and $c+r^2/16$, because
$a_{\rm early}-a_{\rm late}<r^2/64$, and both spatial points lie in the
radius-$r/4$ ball about the later point.  Thus both points lie in one of
the half-cylinder interiors, giving the stated local Hölder estimate.
For $d_{\rm par}(z,z')\ge r/8$, the uniform supremum bound gives
\[
 |u_{\rm top}(z)-u_{\rm top}(z')|
 \le 4C_4/r
 \le 4\cdot 8^{\gamma_0}C_4r^{-1-\gamma_0}
       d_{\rm par}(z,z')^{\gamma_0}.
\]
Hence the glued representative is Hölder on all of $V$, with a uniform
constant.  A Hölder function on $V$ extends uniquely with the same bound
to its closure, including the face $t=t_0$.
\end{proof}


Set
\[
 D=B_1(0)\times(-1,0),\qquad
 \widehat D=B_1(0)\times(-1,0],\qquad F_{\rm top}=B_1(0)\times\{0\}.
\]
Throughout, $(u,Du,p,0)$ is suitable with exponent $q=3$ on $D$.

\begin{definition}[Good point]\label{def:good-point}
A point $z_0=(x_0,t_0)\in\widehat D$ is \emph{good} if there is an $r>0$
such that
\[
 B_r(x_0)\times(t_0-r^2,t_0)\subset D,\qquad
 \overline B_r(x_0)\times[t_0-r^2,t_0)\subset\widehat D,
\]
and
\[
 r^{-2}\iint_{B_r(x_0)\times(t_0-r^2,t_0)}
       (|u|^3+|p|^{3/2})\,dx\,dt<\varepsilon_0/8.
\]
\end{definition}

\begin{lemma}[Good points and Hölder gluing]\label{lem:good-open-glue}
The good-point set is relatively open in $\widehat D$.  Let
$K\subset\overline{B_{1/2}(0)\times(-1/4,0)}$ be compact and suppose every
point of $K$ is good.  Then there are a relatively open neighborhood
$N\subset\widehat D$ of $K$ and a function $\widetilde u:N\to\mathbb R^3$
such that $\widetilde u=u$ almost everywhere on $N\cap D$ and
$\widetilde u|_K$ is parabolically Hölder with exponent $\gamma_0$.
In particular, this representative is Hölder on $K$ including its top
face $t=0$, and extends continuously there.

Moreover, if $z_0=(x_0,t_0)\in K$ is not good, then for every $R>0$ such
that $B_R(x_0)\times(t_0-R^2,t_0)\subset D$,
\begin{equation}\label{eq:bad-point-lower-bound}
 R^{-2}\iint_{B_R(x_0)\times(t_0-R^2,t_0)}
       (|u|^3+|p|^{3/2})\,dx\,dt\ge\varepsilon_0/8.
\end{equation}
\end{lemma}

\begin{proof}
Write $G=|u|^3+|p|^{3/2}$.  Fix a good point with witness radius $r$.
As $r'\uparrow r$, the cylinders
$B_{r'}(x_0)\times(t_0-r'^2,t_0)$ increase to the witness cylinder.
Monotone convergence and the strict smallness allow us to replace $r$ by
some $r'<r$ for which the normalized integral is still below
$\varepsilon_0/8$.  The resulting cylinder has positive spatial and time
margin inside the original witness cylinder.

First note that $G\in L^1_{\rm loc}(D)$.  On each local box
$B\times J\Subset D$, clause (S1) of Definition~\ref{def:sws} gives
$u\in L^\infty(J;L^2(B))$ and $u\in L^2(J;H^1(B))$.  The local Sobolev
inequality gives
\[
 \|u(\cdot,t)\|_{L^6(B)}
 \le C_B\bigl(\|Du(\cdot,t)\|_{L^2(B)}+\|u(\cdot,t)\|_{L^2(B)}\bigr).
\]
Interpolation between $L^2(B)$ and $L^6(B)$ then yields
\[
 \int_J\|u(\cdot,t)\|_{L^3(B)}^3\,dt
 \le \Bigl(\mathop{\rm ess\,sup}_{t\in J}\|u(\cdot,t)\|_{L^2(B)}\Bigr)^{3/2}
       |J|^{1/4}
       \Bigl(\int_J\|u(\cdot,t)\|_{L^6(B)}^2\,dt\Bigr)^{3/4}<\infty.
\]
The pressure clause of (S1) gives $p\in L^{3/2}(B\times J)$, so
$G\in L^1(B\times J)$.

For centers sufficiently close to $z_0$ in the relative topology of
$\widehat D$, the translated open-top cylinders of radius $r'$ remain
admissible in Definition~\ref{def:good-point}.  If $t_0=0$, these centers
satisfy $t'\le0$; the spatial and lower-time margins from $r'<r$ keep their
cylinders inside the original witness cylinder.  The function $G$ is
integrable there by the witness's strict smallness, so it dominates the
translated integrands.  If $t_0<0$, choose $\eta>0$ small enough that
\[
 K_0=\overline B_{r'+\eta}(x_0)\times
       [t_0-r'^2-\eta,t_0+\eta]\Subset D.
\]
Every translated cylinder with center sufficiently close to $z_0$ is
contained in $K_0$.  In both cases the cylinder indicators converge almost
everywhere as the center varies, since their boundaries have measure zero.
Dominated convergence, using $G$ on the witness cylinder at the top and
$\mathbf 1_{K_0}G$ at interior points, gives continuity of the translated
integral.  The strict inequality persists, proving relative openness.

We next construct local representatives around points of $K$.  Fix a good
point $z_0=(x_0,t_0)$ and a witness radius.  Shrink it as in the preceding
paragraph to a radius $r'<r$, retaining strict smallness.  Choose
$\rho$ with $r'<\rho<r$.  The suitable solution restricts to
$B_\rho(x_0)\times(t_0-\rho^2,t_0)$ by
\cite[Lemma~11.5]{CKNLean}, and this cylinder lies in $D$.  Lemma
\ref{lem:thmA-top} gives a representative Hölder on the closed past
half-cylinder centered at $z_0$.  If $t_0=0$, its restriction to a
relative neighborhood of $z_0$ in $\widehat D$ is the required one-sided
representative.

If $t_0<0$, use the interior CKN conclusion to obtain a full open
neighborhood.  At center $(x_0,t_0)$ and radius $r'/2$, the cylinder is
contained in the witness cylinder and its normalized integral is less
than $\varepsilon_0/2$.  Its integral varies continuously under a small time translation.
  By the local-integrability argument above, $G\in L^1_{\rm loc}(D)$.
  Choose a fixed compact set $K_1\Subset D$ containing the closures of all
  these cylinders for sufficiently small translations.  Their indicators
  converge almost everywhere, so dominated convergence applies with
  dominating function $\mathbf 1_{K_1}G$.  Therefore, for some sufficiently
  small $\sigma>0$, the CKN cylinder of radius $r'/2$ centered at
$(x_0,t_0+\sigma)$ has normalized integral less than $\varepsilon_0$ and
closed full cylinder in $D$.  Choose also $\sigma<r'^2/16$ and
$\sigma<-t_0$.  Lemma
\ref{lem:thmA-rescaled} then gives a Hölder representative on an open
neighborhood of $z_0$.

Thus each point of $K$ has a relative neighborhood in $\widehat D$ with
a local representative equal to $u$ almost everywhere on its intersection
with $D$, Hölder with exponent $\gamma_0$.  For any two such neighborhoods
with nonempty intersection, the intersection contains a nonempty open
subset of $D$: this is immediate for an interior point of the intersection;
if it meets only the top face, both neighborhoods are top neighborhoods
and their spatial and past-time interiors overlap below $t=0$.  On that
open subset the two representatives agree almost everywhere and are
continuous, hence agree pointwise; continuity up to the top gives agreement
on the full intersection.  Compactness of $K$ gives a finite subcover.
The local representatives consequently glue to one function on its
relatively open union $N$, and the finite cover shows it equals $u$ almost
everywhere on $N\cap D$.

The finite relative open cover of compact $K$ has a Lebesgue number
$\delta>0$ for the parabolic metric.  Let $H$ be the maximum of the local
Hölder constants and let $M$ bound the glued function on $K$; both are
finite.  Pairs in $K$ at distance less than $\delta$ lie in one patch and
satisfy its Hölder estimate.  For pairs at distance at least $\delta$,
the difference is at most $2M$, which is at most
$2M\delta^{-\gamma_0}d_{\rm par}(z,z')^{\gamma_0}$.  This proves the
claimed Hölder bound on all of $K$.

Finally, suppose the normalized integral in \eqref{eq:bad-point-lower-bound}
were strictly less than $\varepsilon_0/8$ for some such $R$.  As
$r'\uparrow R$ through smaller radii, the open-top cylinders increase to
the cylinder of radius $R$.  Monotone convergence and the strict margin
give an $r'<R$ with normalized integral still less than
$\varepsilon_0/8$.  Since $r'<R$, its closed spatial ball and its
half-open time interval lie inside
$B_R(x_0)\times(t_0-R^2,t_0)$, hence inside $\widehat D$ as required in
Definition~\ref{def:good-point}.  This makes $z_0$ good, a contradiction.
This proves \eqref{eq:bad-point-lower-bound}.
\end{proof}

\begin{remark}[Top-face regularity]
At the top face $t=0$, the regularity conclusion used here is a continuous
Hölder extension from the past, not CKN's open-neighborhood
$\mathrm{IsRegularPoint}$ predicate, since $t=0$ is outside $D$.  Theorem A
alone gives that extension only after the shifted-center gluing in
Lemma~\ref{lem:thmA-top}.  This is the one-sided conclusion needed for
Hölder continuity on the closure in Theorem~\ref{thm:ess-local}; it does
not assert that the boundary point is regular in CKN's sense.
\end{remark}

\begin{lemma}[Projection of the singular set to time]\label{lem:time-projection}
Let $I\subset\mathbb R$ be an open interval, let
$S\subset\mathbb R^3\times I$, and let $\pi_t(x,t)=t$.  If the
one-dimensional parabolic Hausdorff measure of $S$ is zero, then
$\mathcal L^1(\pi_tS)=0$.  In particular, if $S$ is the singular set of a
CKN-suitable solution on $\mathbb R^3\times I$ with $q=3$ and zero force,
then for almost every $t\in I$ and every $\rho>0$ there is a $\delta>0$
such that
\[
 \overline B_\rho(0)\times(t-\delta,t+\delta)
 \quad\text{consists entirely of regular points.}
\]
\end{lemma}

\begin{proof}
The CKN parabolic Hausdorff measure uses the diameter gauge for
$d_{\rm par}$.  If a set in a parabolic cover has diameter $d$, its time
projection has Euclidean diameter at most $d^2$, because
$|t-s|\le d_{\rm par}((x,t),(y,s))^2$.  Thus a parabolic cover of $S$ with
diameters $d_i$ induces a time cover whose one-half-dimensional diameter
cost is at most $\sum_i d_i$.  Taking covers with arbitrarily small total
cost shows that the Euclidean Hausdorff measure of dimension $1/2$ of
$\pi_tS$ is zero.  To see this implies Lebesgue nullity, enclose each time
cover set in an interval of length $\ell_i$; for covers with
$\ell_i\le\eta$, the total interval length is bounded by
$\sqrt{\eta}\sum_i\sqrt{\ell_i}$.  The latter sum can be made arbitrarily
small, proving $\mathcal L^1(\pi_tS)=0$.

For a CKN-suitable solution in the stated class, the CKN version of
Theorem C gives zero one-dimensional parabolic Hausdorff measure of its
singular set.  Also, its regular-point locus is open: if a point has an
open neighborhood with a Hölder representative, that same neighborhood
and representative witness regularity for each of its points.  Hence $S$
is relatively closed in $\mathbb R^3\times I$.  Fix
$t\in I\setminus\pi_tS$ and $\rho>0$.  The compact slice
$\overline B_\rho(0)\times\{t\}$ is disjoint from $S$, and so is covered by
finitely many product neighborhoods contained in the regular-point locus.
The minimum of their time radii is positive; decrease it if necessary so
that $(t-\delta,t+\delta)\subset I$.  These neighborhoods cover
$\overline B_\rho(0)\times(t-\delta,t+\delta)$, as required.
\end{proof}

\begin{remark}
Lemma~\ref{lem:time-projection} supplies neighborhoods consisting of
regular points.  By itself it does not supply the higher Sobolev bounds,
vorticity differential inequality, or infinite-order vanishing hypotheses
of ESS Theorem~4.1; those inputs must be established before that theorem is
applied (ESS 2003, printed p.~230, (4.1)--(4.3)).
\end{remark}

\begin{lemma}[Shifted-cylinder regularity]\label{lem:regular-point-shift}
Let $\Omega\subset\mathbb R^3$ be open and let $I\subset\mathbb R$ be an
open interval.  Set $D=\Omega\times I$, and let $z_0=(x_0,t_0)\in D$ and
$R>0$.  Set
$z_1=(x_0,t_0+R^2/8)$ and suppose
$\overline{\mathcal Q_{R/2}(z_1)}\subset D$.  If $u$ has a representative
Hölder on $\overline{\mathcal Q_{R/2}(z_1)}$, equal to $u$ almost everywhere
on its interior, then $z_0$ is a regular point in the CKN sense.
\end{lemma}

\begin{proof}
The closed half-cylinder is
\[
 \overline B_{R/2}(x_0)\times
 [t_0-R^2/8,t_0+R^2/8].
\]
Thus $z_0$ is strictly inside it in both space and time.  Choose
$0<\delta<R/\sqrt 8$.  The open parabolic ball
$N=\{z:d_{\rm par}(z,z_0)<\delta\}$ lies inside the open spatial ball
and strictly between the two time faces of the half-cylinder, hence
$N\subset D$.  Restrict the given representative to $N$.  It remains
Hölder there and agrees almost everywhere with $u$ there, which is exactly
the CKN open-neighborhood definition of a regular point.

In the global application $D=\mathbb R^3\times(0,T)$ and $0<t_0<T$.
Choose $R$ so that
$t_0-7R^2/8>0$ and $t_0+R^2/8<T$.  The full cylinder centered at $z_1$
has time interval $(t_0-7R^2/8,t_0+R^2/8]$, so its closure is in $D$;
in particular, the half-cylinder hypothesis above holds.  This verifies
the scale and time-domain conditions used with the shifted center.
\end{proof}

\subsection{Vorticity regularity}\label{subsec:vorticity-regularity}

Vorticity removes the pressure from the momentum equation. When the velocity is bounded, its equation is a heat equation with a bounded zeroth-order multiplier in the divergence-form flux. Local heat estimates and div--curl recovery then give the spatial derivatives needed for unique continuation.

For $z_0=(x_0,t_0)$, we use $Q_r^-(z_0)$ and its closure as fixed in
Section~\ref{sec:notation}. Weak derivatives of the velocity are
always carried by the separate field $Du$ of Definition~\ref{def:sws}.

\begin{lemma}[Weak vorticity equation]\label{lem:vorticity-weak-eq}
Let $U=\Omega\times I$ with $\Omega\subset\R^3$ open and $I\subset\R$ an open interval. If $(u,Du,p,0)$ is a suitable weak solution on $U$, define
\[
 \omega_i=\sum_{k,\ell=1}^3\epsilon_{ik\ell}Du_{\ell k},
 \qquad
 C_{ji}=u_j\omega_i-\omega_j u_i\,.
\]
Then, for every $\psi\in\Cc(U;\R^3)$,
\begin{equation}\label{eq:vorticity-weak-eq}
 \iint_U\left(-\omega_i\dt\psi_i-\omega_i\Delta\psi_i\right)
 =\iint_U C_{ji}\,\partial_j\psi_i\,.
\end{equation}
Here and below repeated indices are summed from $1$ to $3$. Equivalently, in distributions on $U$,
\[
 \dt\omega_i-\Delta\omega_i
 =-\sum_{j=1}^3\partial_j\big(u_j\omega_i-\omega_j u_i\big),
 \qquad i=1,2,3\,.
\]
The tensor convention in this formula is $C_{ji}=u_j\omega_i-\omega_j u_i$, with divergence $(\operatorname{div}C)_i=\sum_j\partial_j C_{ji}$.
\end{lemma}

\begin{proof}
Condition (S1) of Definition~\ref{def:sws} gives $u\in L^\infty_{\rm loc}(I;L^2_{\rm loc})\cap L^2_{\rm loc}(I;H^1_{\rm loc})$ and $\omega\in L^2_{\rm loc}$. The three-dimensional interpolation inequality gives $u\in L^{10/3}_{\rm loc}(U)$, so $u\omega\in L^{5/4}_{\rm loc}(U)$ and the right side of~\eqref{eq:vorticity-weak-eq} is integrable.

Use $\varphi=\curl\psi$ in (S3) of Definition~\ref{def:sws}. This is an admissible compactly supported smooth vector field and $\operatorname{div}\varphi=0$, so the pressure term vanishes. The distributional curl identity $\langle\curl u,\psi\rangle=\langle u,\curl\psi\rangle$ changes the time term into $-\iint\omega\cdot\dt\psi$. For the viscous term, the definition of $Du$ and spatial integration by parts give
\[
 \iint_U Du_{ij}\,\partial_j(\curl\psi)_i
 =-\iint_U u_i\Delta(\curl\psi)_i
 =-\iint_U\omega_i\Delta\psi_i\,.
\]
These integrations have no boundary terms because $\psi$ is compactly supported in $U$.

It remains to identify the nonlinear term. On every compact subcylinder, spatial mollification of $u$ preserves its zero divergence on a neighborhood of the test support and converges to $u$ in $L^2_tH^1_x$ and in $L^2_{t,x}$. The corresponding curls converge to $\omega$ in $L^2_{t,x}$. Thus the products $u^\varepsilon\otimes u^\varepsilon$ converge in $L^1_{t,x}$ and $u^\varepsilon\omega^\varepsilon$ converge in $L^1_tL^{3/2}_x$ on the test support. For a smooth divergence-free spatial field $v$, the product rule and $\operatorname{div}(\curl v)=0$ give
\[
 (\curl((v\cdot\nabla)v))_i
 =\partial_j\big(v_j(\curl v)_i-(\curl v)_jv_i\big)\,.
\]
Passing to the limit in the distributional pairing proves the same identity for $u$. Consequently the nonlinear term in (S3), after moving it to the other side, is $\iint C_{ji}\partial_j\psi_i$. This proves~\eqref{eq:vorticity-weak-eq}.
\end{proof}

\begin{lemma}[Localized vorticity energy estimate]\label{lem:localized-vorticity-energy}
Let $J=(a,\tau)$ be a bounded interval. Suppose $z\in L^2(J\times\R^3;\R^3)$ solves
\[
 \dt z_i-\Delta z_i
 =-\partial_j\big(v_jz_i-z_jv_i\big)+\partial_jF_{ji}+g_i
\]
in distributions, where $v\in L^\infty(J\times\R^3;\R^3)$, $F\in L^2(J\times\R^3;\R^{3\times3})$, and $g\in L^2(J;H^{-1}(\R^3;\R^3))$. If the initial trace is $z(a)=z_0\in L^2(\R^3)$, then
\[
 z\in C([a,\tau];L^2(\R^3))\cap L^2(J;H^1(\R^3))
\]
and
\begin{equation}\label{eq:localized-vorticity-energy}
\begin{split}
 \sup_{t\in[a,\tau]}\|z(t)\|_2^2+\int_a^\tau\|\nabla z(t)\|_2^2\,dt
 &\le C\exp\!\left(C(1+\|v\|_\infty^2)(\tau-a)\right)\\
 &\quad\cdot\left(\|z_0\|_2^2+\|F\|_{L^2(J\times\R^3)}^2
       +\|g\|_{L^2(J;H^{-1})}^2\right).
\end{split}
\end{equation}
where $C$ is an absolute constant.
\end{lemma}

\begin{proof}
First take smooth data and a smooth solution. Pair the equation with $z$ and integrate over $\R^3$. Integration by parts has no boundary terms, initially by spatial cutoff and then by passage to the limit in $H^1$. Since
$|v_jz_i-z_jv_i|\le 2|v||z|$, Young's inequality gives
\[
 \frac{d}{dt}\|z(t)\|_2^2+\|\nabla z(t)\|_2^2
 \le C\|v\|_\infty^2\|z(t)\|_2^2
       +C\|F(t)\|_2^2+C\|g(t)\|_{H^{-1}}^2\,.
\]
The $g$ pairing is estimated in $H^{-1},H^1$ duality; the $L^2$ part of the $H^1$ norm contributes $C\|g(t)\|_{H^{-1}}^2+C\|z(t)\|_2^2$, which is included in the Grönwall term. Grönwall's inequality proves~\eqref{eq:localized-vorticity-energy} for smooth solutions.

The trace convention for a distributional solution is the following. Since $z\in L^2(J;L^2)$ and
\[
 \dt z=\Delta z-\operatorname{div}(vz-zv)+\operatorname{div}F+g
       \in L^2(J;H^{-2}(\R^3)),
\]
$z\in W^{1,2}(J;H^{-2})$, so it has a representative in $C([a,\tau];H^{-2})$. The assertion $z(a)=z_0$ means that this representative has value $z_0$ at $a$, under the embedding $L^2\subset H^{-2}$.

We now prove the assertion for the given distributional solution $z$ directly, without constructing a second solution. By the Riesz representation in $H^1$, at almost every time $w=(1-\Delta)^{-1}g(t)\in H^1(\R^3;\R^3)$ gives $g=g_0+\operatorname{div}G$ in the sense $g_i=g_{0,i}+\partial_jG_{ji}$, with $g_0=w$ and $G_{ji}=-\partial_jw_i$. Here $g_0\in L^2(J\times\R^3;\R^3)$, $G\in L^2(J\times\R^3;\R^{3\times3})$ and $\|g_0\|_{L^2(J\times\R^3)}^2+\|G\|_{L^2(J\times\R^3)}^2=\|g\|_{L^2(J;H^{-1})}^2$; the argument below uses only such a representation. Since $z\in L^2(J\times\R^3)$ and $|v_jz_i-z_jv_i|\le\|v\|_\infty(|z_i|+|z_j|)$, the flux
\[
 H_{ji}=-(v_jz_i-z_jv_i)+F_{ji}+G_{ji}
\]
belongs to $L^2(J\times\R^3)$, with $\|H\|_2^2\le 36\|v\|_\infty^2\|z\|_2^2+3\|F\|_2^2+3\|G\|_2^2$ on every time slab, and each component solves the linear equation
\[
 \dt z_i-\Delta z_i=\partial_jH_{ji}+g_{0,i}
\]
in distributions on $J\times\R^3$, with all data in $L^2(J\times\R^3)$.

Let $k\in C_c^\infty(\R^3)$, and write $k\ast$ for convolution in the space variable. Fix $x\in\R^3$ and test the equation with $k(x-\cdot)\,\theta$, $\theta\in C_c^\infty(J)$. This shows that $t\mapsto(k\ast z_i(t))(x)$ has the distributional derivative
\[
 S_{k,i}(x,t)=\bigl((\Delta k)\ast z_i(t)+(\partial_jk)\ast H_{ji}(t)+k\ast g_{0,i}(t)\bigr)(x)
\]
on $J$, which belongs to $L^2(J)$. By the trace convention, applied to the smooth compactly supported test $k(x-\cdot)$, the continuous representative of this function has the value $(k\ast z_{0,i})(x)$ at $a$. Hence, for almost every $t\in J$,
\[
 (k\ast z_i(t))(x)=Z_{k,i}(x,t):=(k\ast z_{0,i})(x)+\int_a^tS_{k,i}(x,s)\,ds,
\]
and $t\mapsto Z_{k,i}(x,t)$ is absolutely continuous on $[a,\tau]$. Differentiate $Z_{k,i}(x,t)^2$ in $t$, integrate over $(a,t)$ and over $x\in\R^3$ (Fubini), and integrate by parts in $x$. The kernels are smooth with compact support and all fields are in $L^2$, so no boundary terms occur, and $(\partial_jk)\ast H_{ji}=\partial_j(k\ast H_{ji})$ and $(\Delta k)\ast z_i=\Delta Z_{k,i}$ for almost every $t$. Together with $2|\langle\nabla Z_k,k\ast H\rangle|\le\|\nabla Z_k\|_2^2+\|k\ast H\|_2^2$ and $2|\langle Z_k,k\ast g_0\rangle|\le\|Z_k\|_2^2+\|k\ast g_0\|_2^2$, this gives for every $t\in[a,\tau]$
\begin{equation*}
 \|Z_k(t)\|_2^2+\int_a^t\|\nabla Z_k\|_2^2
 \le\|k\ast z_0\|_2^2+\int_a^t\bigl(\|k\ast H\|_2^2+\|k\ast z\|_2^2+\|k\ast g_0\|_2^2\bigr).
\end{equation*}

Let $\rho_\varepsilon$ be a standard spatial mollifier and apply this inequality with $k=\rho_\varepsilon-\rho_\delta$. Since $\|\rho_\varepsilon\ast h-h\|_{L^2}\to0$ for every $h$ in $L^2(\R^3)$ and in $L^2(J\times\R^3)$, the right side tends to zero as $\varepsilon,\delta\to0$. Hence the curves $Z_{\rho_\varepsilon}$ are Cauchy in $C([a,\tau];L^2(\R^3))$, and $\nabla Z_{\rho_\varepsilon}$ is Cauchy in $L^2(J\times\R^3)$. The limit curve is continuous in $L^2$, equals $z(t)$ for almost every $t$, and has the value $z_0$ at $t=a$; the limit of the gradients is the weak gradient of $z$. This proves $z\in C([a,\tau];L^2(\R^3))\cap L^2(J;H^1(\R^3))$. Passing to the limit in the inequality with $k=\rho_\varepsilon$ gives it for $z$ without the kernel. With the bound on $\|H\|_2^2$, the quantities $e(t)=\|z(t)\|_2^2$ and $d(t)=\int_a^t\|\nabla z\|_2^2$ satisfy
\[
 e(t)+d(t)\le\|z_0\|_2^2+3\|F\|_2^2+3\|G\|_2^2+\|g_0\|_2^2+(36\|v\|_\infty^2+1)\int_a^te(s)\,ds,
\]
where the norms of $F$, $G$ and $g_0$ are taken over $J\times\R^3$. Grönwall's inequality then gives~\eqref{eq:localized-vorticity-energy} with $C=36$.
\end{proof}

\begin{lemma}[Local heat gain]\label{lem:local-heat-gain}
Let $B_r\Subset B_R\subset\R^3$ be concentric balls, $a<s<b$, and let $m=1$ or $m=2$. If $z\in L^2((a,b);H^m(B_R))$ and
\[
 \dt z-\Delta z=G\quad\hbox{in }B_R\times(a,b),
 \qquad G\in L^2((a,b);H^{m-1}(B_R)),
\]
then, for each $a<s<b$,
\[
 z\in C([s,b];H^m(B_r))\cap L^2((s,b);H^{m+1}(B_r))
\]
and
\begin{equation}\label{eq:local-heat-gain}
 \|z\|_{C([s,b];H^m(B_r))}+\|z\|_{L^2((s,b);H^{m+1}(B_r))}
 \le C\left(\|z\|_{L^2((a,b);H^m(B_R))}
              +\|G\|_{L^2((a,b);H^{m-1}(B_R))}\right),
\end{equation}
where $C$ depends only on $m$, the radii and their gap, and $s-a,b-a$.
\end{lemma}

\begin{proof}
Choose a smooth spatial cutoff supported in $B_R$ and equal to $1$ on $B_r$, and a smooth time cutoff that vanishes near $a$ and equals $1$ on $[s,b]$. Multiplying by their product $\eta$ gives, after extension by zero in space,
\[
 \dt(\eta z)-\Delta(\eta z)
 =\eta G+(\dt\eta-\Delta\eta)z-2\nabla\eta\cdot\nabla z\,.
\]
The right side belongs to $L^2((a,b);H^{m-1}(\R^3))$: the last term uses $z\in L^2H^m$, and all cutoff derivatives are bounded. The field $\eta z$ vanishes for $t$ near $a$. For a smooth zero-initial solution $w$ of $\dt w-\Delta w=H$ with compact support in space, the energies of the spatial derivatives $\partial^\beta w$, $|\beta|\le m$, give
\[
 \sup_{t\in[a,b]}\|w(t)\|_{H^m}^2+\int_a^b\|w(t)\|_{H^{m+1}}^2\,dt
 \le C(m,b-a)\int_a^b\|H(t)\|_{H^{m-1}}^2\,dt\,.
\]
At integer order we work with the derivative energies $\sum_{|\beta|\le m}\|\partial^\beta w(t)\|_2^2$ in place of the Bessel norm: the Fourier weights $(1+|\xi|^2)^m$ and $\sum_{|\beta|\le m}\xi^{2\beta}$ are comparable, so these energies define a norm equivalent to $\|w(t)\|_{H^m}^2$ with constants depending only on $m$. Each $\partial^\beta w$ satisfies its own energy identity, obtained by pairing its equation with $\partial^\beta w$. For $\beta=\beta'k$ we write $\partial^\beta H=\partial_k\partial^{\beta'}H$ and move $\partial_k$ onto $\partial^\beta w$, so that only $m-1$ derivatives of $H$ enter; Young's inequality absorbs half of $\|\nabla\partial^\beta w\|_2^2$. The remaining term $\|\partial^\beta w\|_2^2$ is handled by Grönwall's inequality, which produces a factor depending on $b-a$ (of size $e^{b-a}$); for long intervals this dependence cannot be dropped, so the constant is $C(m,b-a)$.

The same argument applies to $w=\eta z$, with the energy estimate of Lemma~\ref{lem:localized-vorticity-energy} in place of the smooth computation. Write $H=\eta G+(\dt\eta-\Delta\eta)z-2\nabla\eta\cdot\nabla z$. Testing with derivatives of test functions shows that each $\partial^\beta w$, $|\beta|\le m$, solves $\dt\partial^\beta w-\Delta\partial^\beta w=\partial^\beta H$ in distributions on $\R^3\times(a,b)$. For $|\beta|\le m-1$ this source belongs to $L^2$; for $\beta=\beta'k$ it is the divergence $\partial_k(\partial^{\beta'}H)$ of the $L^2$ field $\partial^{\beta'}H$. Since $\partial^\beta w$ vanishes for $t$ near $a$, its pairings with spatial test functions have absolutely continuous versions with value $0$ at $a$, so its initial trace is zero. Lemma~\ref{lem:localized-vorticity-energy} with $v=0$ therefore gives $\partial^\beta w\in C([a,b];L^2(\R^3))\cap L^2((a,b);H^1(\R^3))$, with
\[
 \sup_{t\in[a,b]}\|\partial^\beta w(t)\|_2^2+\int_a^b\|\nabla\partial^\beta w\|_2^2\,dt
 \le C\int_a^b\bigl(\|H\|_{H^{m-1}}^2+\|\partial^\beta w\|_2^2\bigr)\,dt\,,
\]
where $C$ depends on $b-a$. The weak-derivative relations between the continuous curves of $\partial^\beta w$ and $\partial^{\beta j}w$ hold at almost every time, and by continuity at every time. Summing over $|\beta|\le m$, using that $\eta=1$ on $B_r\times[s,b]$ and bounding the right side by $\|z\|_{L^2((a,b);H^m(B_R))}^2+\|G\|_{L^2((a,b);H^{m-1}(B_R))}^2$ gives $z\in C([s,b];H^m(B_r))\cap L^2((s,b);H^{m+1}(B_r))$ and~\eqref{eq:local-heat-gain}. There is no terminal boundary term to discard: the estimate includes the value at $b$ in its supremum.

The bootstrap variant at the end of the proof of Theorem~\ref{thm:vorticity-regularity} uses only the following first-order form of this gain, which needs no prior information on $\nabla w$ and gives bounds up to the top time. Let $\rho_\varepsilon\ge0$ be a smooth space-time mollifier with integral one, supported in $B_\varepsilon\times(-\varepsilon,\varepsilon)$. For a field $w$ on $B_R\times(a,b)$ define its backward mollification
\[
 w_\varepsilon(x,t)=\iint\rho_\varepsilon(y,\sigma)\,w(x-y,t-5\varepsilon-\sigma)\,dy\,d\sigma\,.
\]
It uses only values of $w$ at times in $(t-6\varepsilon,t-4\varepsilon)$. Hence it is smooth on $B_{R-\varepsilon}\times(a+6\varepsilon,b]$, including the top time, and it commutes with weak derivatives. Suppose $w,F_1,F_2,F_3\in L^2(B_R\times(a,b))$ and $\dt w-\Delta w=\operatorname{div}F$ in distributions. Then $\dt w_\varepsilon-\Delta w_\varepsilon=\operatorname{div}F_\varepsilon$ pointwise. Fix $r<R'<R$ and $a<a'<s<b$. Choose a cutoff $\eta$ supported in $B_{R'}$ in space, vanishing near $a'$ in time, and equal to $1$ on $B_r\times[s,b]$. Pair the equation with $\eta^2w_\varepsilon$ on $B_{R'}\times(a',t)$ and absorb the cross terms by Young's inequality. For small $\varepsilon$ this gives
\[
 \sup_{s<t<b}\|w_\varepsilon(t)\|_{L^2(B_r)}^2+\int_s^b\|\nabla w_\varepsilon(t)\|_{L^2(B_r)}^2\,dt
 \le C\bigl(\|w_\varepsilon\|_{L^2(B_{R'}\times(a',b))}^2+\|F_\varepsilon\|_{L^2(B_{R'}\times(a',b))}^2\bigr),
\]
where $C$ depends only on the gaps. As $\varepsilon\to0$, $w_\varepsilon\to w$ and $F_\varepsilon\to F$ in $L^2(B_{R'}\times(a',b))$. Thus $\nabla w_\varepsilon$ is bounded in $L^2(B_r\times(s,b))$. Every weak limit is the weak gradient of $w$ there, so
\[
 \|\nabla w\|_{L^2(B_r\times(s,b))}
 \le C\bigl(\|w\|_{L^2(B_R\times(a,b))}+\|F\|_{L^2(B_R\times(a,b))}\bigr).
\]
The same estimate for $w_\varepsilon-w_\delta$ shows that $w_\varepsilon(t)$ is Cauchy in $L^2(B_r)$, uniformly in $t\in(s,b)$, as $\varepsilon,\delta\to0$.
\end{proof}

\begin{lemma}[Local div--curl recovery]\label{lem:local-div-curl}
Let $B_r\Subset B_R\subset\R^3$ be concentric balls, let $m\in\{0,1,2\}$, and suppose $v\in H^m(B_R;\R^3)$, $\operatorname{div}v=0$, and $\curl v=\zeta\in H^m(B_R;\R^3)$. Then
\begin{equation}\label{eq:local-div-curl}
 \|v\|_{H^{m+1}(B_r)}
 \le C\left(\|\zeta\|_{H^m(B_R)}+\|v\|_{H^m(B_R)}\right),
\end{equation}
where $C$ depends only on $m,r,R$. The same estimate holds after taking an $L^\infty$ or $L^2$ norm in time.
\end{lemma}

\begin{proof}
For a compactly supported smooth vector field $W$ on $\R^3$, two integrations by parts give
\[
 \int_{\R^3}|\nabla W|^2
 =\int_{\R^3}(\operatorname{div}W)^2+\int_{\R^3}|\curl W|^2\,,
\]
the integrated form of $|\xi|^2|\widehat W(\xi)|^2=|\xi\cdot\widehat W(\xi)|^2+|\xi\times\widehat W(\xi)|^2$. Spatial derivatives commute with $\operatorname{div}$ and $\curl$. Applying the identity to $W=\partial^\alpha V$ for every derivative $\partial^\alpha$ of order $|\alpha|\le m$ and adding $\|V\|_{H^m}^2$ proves
\[
 \|V\|_{H^{m+1}(\R^3)}
 \le C_m\big(\|\operatorname{div}V\|_{H^m}
          +\|\curl V\|_{H^m}+\|V\|_{H^m}\big)
\]
for compactly supported smooth $V$. Now let $V\in H^m(\R^3;\R^3)$ have compact support, and let $\operatorname{div}V$ and $\curl V$, taken in distributions, belong to $H^m$. For a standard mollifier $\rho_\varepsilon$ the fields $V_\varepsilon=\rho_\varepsilon\ast V$ are smooth with compact support, $\operatorname{div}V_\varepsilon=\rho_\varepsilon\ast\operatorname{div}V$ and $\curl V_\varepsilon=\rho_\varepsilon\ast\curl V$. Since $\rho_\varepsilon\ast f\to f$ in $H^m$ for $f\in H^m$, the smooth estimate applied to $V_\varepsilon-V_\delta$ shows that $V_\varepsilon$ is Cauchy in $H^{m+1}(\R^3)$. The limits of the derivatives of order $m+1$ are the weak derivatives of $V$, so $V\in H^{m+1}(\R^3)$ and the estimate holds for $V$. Choose a cutoff $\chi$ supported in $B_R$ and equal to $1$ on $B_r$, with an intermediate concentric ball for its derivatives. For $V=\chi v$, extended by zero, testing the equations for $v$ with $\chi\varphi$ gives in distributions on $\R^3$
\[
 \operatorname{div}V=\nabla\chi\cdot v,
 \qquad \curl V=\chi\zeta+\nabla\chi\times v\,.
\]
By the product rule these expressions and $V$ are bounded in $H^m(\R^3)$ by $C(\|\zeta\|_{H^m(B_R)}+\|v\|_{H^m(B_R)})$. The estimate for compactly supported fields and $V=v$ on $B_r$ give~\eqref{eq:local-div-curl}. Integrating the pointwise-in-time estimate, or taking its essential supremum, gives the time-dependent versions.

The bootstrap variant at the end of the proof of Theorem~\ref{thm:vorticity-regularity} uses the case $m=0$ in time-integrated form. In components, the identity above reads
\[
 \int|\nabla V|^2=\int(\operatorname{div}V)^2
   +\frac12\sum_{i,k=1}^3\int(\partial_iV_k-\partial_kV_i)^2
\]
for compactly supported smooth $V$. With the cutoff above, for smooth $v$ and $r<R'$ it bounds $\|\nabla v\|_{L^2(B_r)}$ by the $L^2(B_{R'})$ norms of $\operatorname{div}v$, of the antisymmetric derivatives $\partial_iv_k-\partial_kv_i$, and of $v$. Now let $Y\in L^2(B_R\times(a,b);\R^3)$ satisfy $\operatorname{div}Y=0$ and $\partial_iY_k-\partial_kY_i\in L^2(B_R\times(a,b))$ in distributions. Apply the smooth bound at each time to the backward mollifications $Y_\varepsilon$ from the proof of Lemma~\ref{lem:local-heat-gain}. They are divergence free, and their antisymmetric derivatives are the mollified ones. Integrate over $(s,b)$ with $a<s<b$. The gradients are then bounded in $L^2(B_r\times(s,b))$, and a weak limit gives
\[
 \|\nabla Y\|_{L^2(B_r\times(s,b))}
 \le C\Bigl(\sum_{i,k=1}^3\|\partial_iY_k-\partial_kY_i\|_{L^2(B_R\times(a,b))}
       +\|Y\|_{L^2(B_R\times(a,b))}\Bigr).
\]
No time-continuous representative of $Y$ is needed.
\end{proof}

\begin{lemma}[Products at the two bootstrap levels]\label{lem:vorticity-products}
For compactly supported $v\in H^2(\R^3;\R^3)\cap L^\infty$, and for $v\in H^3(\R^3;\R^3)$, respectively,
\begin{align}
 \|v\otimes v\|_{H^2}&\le C\|v\|_\infty\|v\|_{H^2},\label{eq:vorticity-product-h2}\\
 \|v\otimes v\|_{H^3}&\le C\|v\|_{H^2}\|v\|_{H^3}\,.
 \label{eq:vorticity-product-h3}
\end{align}
The corresponding local estimates hold after inserting a cutoff; their constants also depend on the distance between the two balls.
\end{lemma}

\begin{proof}
The derivatives of order zero and one in~\eqref{eq:vorticity-product-h2} are bounded by the product rule and $\|v\|_\infty$. For the second derivatives, the only additional term is bounded by $C\|\nabla v\|_4^2$. Integration by parts for compactly supported smooth $v$ gives
\[
\begin{aligned}
 \int_{\R^3}|\nabla v|^4
 &=-\sum_{a,j=1}^3\int_{\R^3}v_a\,\partial_j\big(|\nabla v|^2\partial_jv_a\big)\\
 &\le C\|v\|_\infty\int_{\R^3}|\nabla v|^2|\nabla^2v|\\
 &\le C\|v\|_\infty\|\nabla v\|_4^2\|\nabla^2v\|_2\,.
\end{aligned}
\]
If $\|\nabla v\|_4\ne0$, division yields
$\|\nabla v\|_4^2\le C\|v\|_\infty\|\nabla^2v\|_2$; the zero case is immediate. Approximation proves the first estimate. For~\eqref{eq:vorticity-product-h3}, the same argument applied to two compactly supported smooth scalar functions $f,g$ gives the bilinear form $\|fg\|_{H^2}\le C(\|f\|_\infty\|g\|_{H^2}+\|g\|_\infty\|f\|_{H^2})$: the mixed term $\nabla f\otimes\nabla g$ is bounded by $\|\nabla f\|_4\|\nabla g\|_4$ and the inequality above. With $H^2\hookrightarrow L^\infty$ this shows that $H^2(\R^3)$ is an algebra, $\|fg\|_{H^2}\le C\|f\|_{H^2}\|g\|_{H^2}$. The derivatives of order at most three of $v_iv_j$ are those of order at most two of $v_iv_j$ and of $\partial_k(v_iv_j)=\partial_kv_i\,v_j+v_i\,\partial_kv_j$. Hence
\[
 \|v_iv_j\|_{H^3}
 \le C\Bigl(\|v\|_{H^2}^2+\sum_{k=1}^3\|\partial_kv\|_{H^2}\|v\|_{H^2}\Bigr)
 \le C\|v\|_{H^2}\|v\|_{H^3}\,,
\]
since $\|\partial_kv\|_{H^2}\le\|v\|_{H^3}$ and $\|v\|_{H^2}\le\|v\|_{H^3}$. Smooth approximation and cutoffs give the stated local forms.

The bootstrap variant at the end of the proof of Theorem~\ref{thm:vorticity-regularity} uses these bounds only in the following local, time-integrated forms. For smooth $f$ with $|f|\le M$ on $B_R$ and $r<R$, the integration by parts above, with weight $\eta^4$ for a cutoff $\eta$ equal to one on $B_r$, and Young's inequality give
\[
 \int_{B_r}|\nabla f|^4\le CM^2\int_{B_R}\bigl(|\nabla f|^2+|\nabla^2f|^2\bigr).
\]
Apply this at each time to the components of the backward mollifications $u_\varepsilon$ of $u$; these are also bounded by $M$. Integrate in time and pass to the limit by Fatou's lemma. If $|u|\le M$ and $Du,D^2u\in L^2$ on a box, then $Du\in L^4$ on a smaller box, and $\|Du\|_4^4\le CM^2(\|Du\|_2^2+\|D^2u\|_2^2)$. At the next level, fix a time. The embedding $H^2\hookrightarrow L^\infty$ for $\nabla u_\varepsilon(t)$ and the smooth case of Lemma~\ref{lem:local-div-curl} bound $\sup_{B_r}|\nabla u_\varepsilon(t)|^2$ by the squared $L^2$ norms of $u_\varepsilon(t)$, $\omega_\varepsilon(t)$, $\nabla\omega_\varepsilon(t)$ and $\nabla^2\omega_\varepsilon(t)$ on a larger ball. Also, by the same div--curl estimate and the uniform-in-time bounds of the first-order heat gain in the proof of Lemma~\ref{lem:local-heat-gain} for $\omega$ and $\nabla\omega$, the sum $\|D^2u_\varepsilon(t)\|_2+\|\nabla\omega_\varepsilon(t)\|_2$ on that ball is bounded uniformly in $t$. Integrating in time and applying Fatou's lemma give
\[
 \iint|Du|^2\bigl(|D^2u|^2+|\nabla\omega|^2\bigr)<\infty
\]
on a smaller box, with a bound by the data. These are the products in the first and second spatial derivatives of the flux $u_j\omega_i-\omega_ju_i$.
\end{proof}

\begin{theorem}[Vorticity regularity near a bounded-velocity point]\label{thm:vorticity-regularity}
Let $z_0=(x_0,t_0)$, and let $U=\Omega\times I$ be an open neighborhood of $\overline{Q_1^-(z_0)}$. Suppose $q>\nf52$, $M\ge0$, $(u,Du,p,0)$ is suitable on $U$, and $|u|\le M$ almost everywhere on $Q_1^-(z_0)$. Set
\[
 E=\left(\iint_{Q_1^-(z_0)}|Du|^2\,dx\,dt\right)^{1/2}\,.
\]
Then there is a constant $C=C(M,E)$ such that, on $Q_{1/2}^-(z_0)$,
\begin{align}
 \operatorname*{ess\,sup}_{Q_{1/2}^-(z_0)}\big(|u|+|Du|\big)&\le C,\label{eq:vorticity-regularity-gradient}\\
 \|\omega\|_{W^{2,1}_2(Q_{1/2}^-(z_0))}
 +\sup_{t\in[t_0-1/4,t_0]}\|\widetilde\omega(\cdot,t)\|_{H^2(B_{1/2}(x_0))}
 +\|\widetilde\omega\|_{C(\overline{Q_{1/2}^-(z_0)})}&\le C,
 \label{eq:vorticity-regularity-omega}
\end{align}
where $\omega=\curl u$ is the weak vorticity and $\widetilde\omega$ is a representative continuous on the closed past half-cylinder. In addition,
\begin{equation}\label{eq:vorticity-regularity-inequality}
 |\dt\omega-\Delta\omega|
 \le C\big(|\omega|+|\nabla\omega|\big)
 \quad\text{almost everywhere on }Q_{1/2}^-(z_0)\,.
\end{equation}
The bounds for $u$ and $Du$ are essential bounds on the open past half-cylinder, uniform as $t\uparrow t_0$; no time-continuous representative of $u$ or $Du$ is asserted.
\end{theorem}

\begin{proof}
The suitable-solution energy clause gives $E<\infty$ and $u\in L^2_tH^1_x$ on compact subcylinders. The weak vorticity equation is~\eqref{eq:vorticity-weak-eq}. We use nested radii
\[
 1>\frac78>\frac34>\frac58>\frac9{16}>\frac{17}{32}>\frac{33}{64}>\frac{65}{128}>\frac12\,.
\]
All cutoff constants below depend only on these fixed gaps. We denote by $C(M,E)$ a finite bound depending only on the displayed parameters and the fixed radii; its value may increase from line to line.

\emph{Initial energy level.} On $Q_1^-$, $\|\omega\|_{L^2}\le C E$. Let $J=(t_0-1,t_0)$, and choose a spatial cutoff $\varphi\in C_c^\infty(B_1(x_0))$ equal to one on $\overline{B_{7/8}(x_0)}$ and a time cutoff $\theta$ that vanishes near $t_0-1$ and equals one on $[t_0-(7/8)^2,t_0]$. Set $\chi(x,t)=\varphi(x)\theta(t)$. Define $v=u$ on $B_1(x_0)\times J$ and extend it by zero to $J\times\R^3$; extend $z=\chi\omega$ by zero outside $B_1(x_0)$. Then $\|v\|_\infty\le M$ and $v_jz_i-z_jv_i=\chi C_{ji}$. Multiplying the weak equation by $\chi$ gives
\[
 \dt(\chi\omega_i)-\Delta(\chi\omega_i)
 =-\partial_j\big(\chi C_{ji}+2\omega_i\partial_j\chi\big)
   +\partial_j\chi\,C_{ji}+(\dt\chi+\Delta\chi)\omega_i\,.
\]
Every term on the right is in $L^2_tH^{-1}_x$; the two non-divergence terms are in $L^2_{t,x}$, and the divergence flux is in $L^2_{t,x}$ because $|C|\le2M|\omega|$. The cutoff field has zero initial trace. Lemma~\ref{lem:localized-vorticity-energy}, with the zero-extended multiplier $v$, therefore gives
\begin{equation}\label{eq:vorticity-base-energy}
 \omega\in C([t_0-(7/8)^2,t_0];L^2(B_{7/8}))
       \cap L^2((t_0-(7/8)^2,t_0);H^1(B_{7/8})),
 \qquad \|\omega\|_{L^\infty L^2\cap L^2H^1}\le C(M,E)\,.
\end{equation}
No spatial boundary term occurs because the cutoff is compactly supported in $B_1$; the energy estimate includes the terminal time $t_0$.

\emph{Velocity through order two.} For almost every time, $\operatorname{div}u=0$ by (S2) and $\curl u=\omega$ by definition. Lemma~\ref{lem:local-div-curl} with $m=0$, together with $\|u(t)\|_{L^2(B_{7/8})}\le C M$, gives
\[
\begin{aligned}
 u&\in L^\infty((t_0-(3/4)^2,t_0);H^1(B_{3/4}))\\
  &\cap L^2((t_0-(3/4)^2,t_0);H^2(B_{3/4})),\\
 \|u\|_{L^\infty H^1\cap L^2H^2}&\le C(M,E)\,.
\end{aligned}
\]
Choose a spatial cutoff supported in $B_{3/4}$ and equal to one on $B_{5/8}$. Applying~\eqref{eq:vorticity-product-h2} to the cutoff velocity gives $u\otimes u\in L^2_tH^2_x(B_{5/8})$, with norm at most $C(M,E)$. Since $\operatorname{div}u=0$, the weak vorticity equation is also
\[
 \dt\omega_i-\Delta\omega_i
 =-\epsilon_{ik\ell}\partial_k\partial_j(u_\ell u_j)\,.
\]
This identity follows by taking the curl of the weak momentum equation; it is valid distributionally because $u\otimes u\in L^1_{\rm loc}$ and its localized derivatives just obtained are in $L^2_tH^2_x(B_{5/8})$. Its right side belongs to $L^2_tL^2_x(B_{5/8})$. Apply Lemma~\ref{lem:local-heat-gain} with $m=1$, outer ball $B_{5/8}$, inner ball $B_{9/16}$, initial time $t_0-(3/4)^2$, and terminal time $t_0$. The prior $L^2_tH^1_x$ bound follows from~\eqref{eq:vorticity-base-energy}. We obtain
\[
\begin{aligned}
 \omega&\in C([t_0-(9/16)^2,t_0];H^1(B_{9/16}))\\
  &\cap L^2((t_0-(9/16)^2,t_0);H^2(B_{9/16})),\\
 \|\omega\|_{C H^1\cap L^2H^2}&\le C(M,E)\,.
\end{aligned}
\]
A first application of Lemma~\ref{lem:local-div-curl} with $m=1$ from $B_{9/16}$ to $B_{17/32}$ uses the $L^\infty_tH^1_x$ bound on $\omega$ and gives $u\in L^\infty_tH^2_x(B_{17/32})$. A separate application with $m=2$ on the same pair of balls uses $\omega\in L^2_tH^2_x(B_{9/16})$ from the heat gain and $u\in L^2_tH^2_x(B_{9/16})$ from the earlier $m=0$ step; it gives $u\in L^2_tH^3_x(B_{17/32})$. Combining these outputs yields
\[
\begin{aligned}
 u&\in L^\infty((t_0-(9/16)^2,t_0);H^2(B_{17/32}))\\
  &\cap L^2((t_0-(9/16)^2,t_0);H^3(B_{17/32})),\\
 \|u\|_{L^\infty H^2\cap L^2H^3}&\le C(M,E)\,.
\end{aligned}
\]

\emph{One more vorticity level.} A cutoff of $u$ supported in $B_{17/32}$ and equal to one on $B_{33/64}$ is bounded in $L^\infty_tH^2_x$ and $L^2_tH^3_x$. Estimate~\eqref{eq:vorticity-product-h3} therefore gives $u\otimes u\in L^2_tH^3_x(B_{33/64})$, with norm bounded by $C(M,E)$. The right side of the last vorticity equation is now in $L^2_tH^1_x(B_{33/64})$. The prior estimate gives $\omega\in L^2_tH^2_x$ there, so Lemma~\ref{lem:local-heat-gain} with $m=2$, outer ball $B_{33/64}$, inner ball $B_{65/128}$, initial time $t_0-(17/32)^2$ and terminal time $t_0$ yields
\[
\begin{aligned}
 \omega&\in C([t_0-(65/128)^2,t_0];H^2(B_{65/128}))\\
  &\cap L^2((t_0-(65/128)^2,t_0);H^3(B_{65/128})),\\
 \|\omega\|_{C H^2\cap L^2H^3}&\le C(M,E)\,.
\end{aligned}
\]
Applying Lemma~\ref{lem:local-div-curl} with $m=2$ from $B_{65/128}$ to $B_{1/2}$ gives
\[
 u\in L^\infty((t_0-1/4,t_0);H^3(B_{1/2})),
 \qquad \|u\|_{L^\infty H^3}\le C(M,E)\,.
\]
The Sobolev embedding $H^3(B_{1/2})\hookrightarrow W^{1,\infty}(B_{1/2})$ proves the gradient bound in~\eqref{eq:vorticity-regularity-gradient}. The same level gives $\|\widetilde\omega(t)\|_{H^2(B_{1/2})}\le C(M,E)$ for every $t$ in the closed time interval. The continuous embedding $H^2(B_{65/128})\hookrightarrow C(\overline{B_{1/2}})$ and continuity in time with values in $H^2$ supply the representative $\widetilde\omega\in C(\overline{Q_{1/2}^-})$.

Finally, on $Q_{1/2}^-$ the vorticity equation gives
\[
 \dt\omega=\Delta\omega-\partial_j(u_j\omega_i-\omega_j u_i)\,.
\]
The level-two bounds put $\omega,\nabla\omega,\nabla^2\omega$ in $L^2$ and the right side in $L^2$. Indeed, $\Delta\omega\in L^2$; expanding the flux derivative is legitimate since $u$ and $Du$ are essentially bounded, $\operatorname{div}u=\operatorname{div}\omega=0$, and it is bounded pointwise by $|u||\nabla\omega|+|Du||\omega|$. Thus $\omega\in W^{2,1}_2(Q_{1/2}^-)$ and
\[
 |\dt\omega-\Delta\omega|
 \le |u|\,|\nabla\omega|+|Du|\,|\omega|
 \le C(M,E)(|\omega|+|\nabla\omega|)
\]
a.e. This proves~\eqref{eq:vorticity-regularity-omega}--\eqref{eq:vorticity-regularity-inequality}.

\emph{A variant of the bootstrap.} The same conclusions also follow from $L^2$ space-time bounds alone, on nested boxes $B_\rho(x_0)\times(t_0-\sigma,t_0)$ with the same top time, without Lemma~\ref{lem:localized-vorticity-energy}. All smoothing uses the backward mollifications from the proof of Lemma~\ref{lem:local-heat-gain}, which are smooth up to $t_0$. Each box below is slightly smaller than the previous one, and every bound depends only on $M$, $E$ and the fixed gaps. Write $C_{ji}=u_j\omega_i-\omega_ju_i$ as in Lemma~\ref{lem:vorticity-weak-eq}. Then $|C|\le2M|\omega|$ and $|\omega|\le\sqrt2|Du|$.

The vorticity solves $\dt\omega_i-\Delta\omega_i=-\partial_jC_{ji}$ with $\omega,C\in L^2$, so the first-order heat gain gives $\nabla\omega\in L^2$. For each $j$, the field $\partial_ju$ is divergence free, and its antisymmetric derivatives are given by $\partial_j\omega$. The integrated div--curl bound therefore gives $D^2u\in L^2$, and the $L^4$ bound gives $Du\in L^4$. Differentiating the weak equation shows that $\partial_m\omega$ solves the heat equation with flux $\partial_mC$. Since $|\partial_mC|\le2|Du||\omega|+2M|\nabla\omega|$, this flux is in $L^2$, and the first-order heat gain gives $\nabla^2\omega\in L^2$. By the integrability of $|Du|^2(|D^2u|^2+|\nabla\omega|^2)$, the flux $\partial_k\partial_mC$ in the equation for $\partial_k\partial_m\omega$ is also in $L^2$.

Each of $\omega$, $\nabla\omega$ and $\nabla^2\omega$ therefore solves a heat equation with $L^2$ flux. Their backward mollifications are Cauchy in $L^2$ of a ball, uniformly in time. By $H^2\hookrightarrow L^\infty$ on a slightly larger ball, the mollifications of $\omega$ are uniformly Cauchy on $\overline{Q_{1/2}^-(z_0)}$. The limit is the continuous representative $\widetilde\omega$, and the same uniform bounds give $\sup_t\|\widetilde\omega(t)\|_{H^2(B_{1/2})}+\|\widetilde\omega\|_{C(\overline{Q_{1/2}^-})}\le C(M,E)$. Next, fix a time. The embedding $H^2\hookrightarrow L^\infty$ for $\nabla u_\varepsilon(t)$ and the smooth div--curl estimate at that time bound $|\nabla u_\varepsilon|$ on $Q_{1/2}^-$ by $C(M,E)$. These use $|u_\varepsilon|\le M$ and the uniform-in-time bounds of the mollified $\omega$, $\nabla\omega$ and $\nabla^2\omega$. Since $\nabla u_\varepsilon\to Du$ in $L^2$, a subsequence converges almost everywhere, which gives~\eqref{eq:vorticity-regularity-gradient}. Finally, $\dt\omega_i=\Delta\omega_i-\partial_jC_{ji}$ has right side in $L^2$. Hence $\omega\in W^{2,1}_2(Q_{1/2}^-)$ with the stated bound, and the pointwise flux estimate above gives~\eqref{eq:vorticity-regularity-inequality}. This variant does not use time continuity of $u$ or $Du$.
\end{proof}

\begin{remark}[Departure from \cite{ESS2003}]
ESS Lemma~2.2 and Remark~2.3~\cite[Lemma~2.2, Remark~2.3, pp.~216, 220--221]{ESS2003} refer the higher spatial orders to a smooth-solution estimate and linear Stokes regularity; that citation does not by itself justify the passage from a suitable weak solution to the regularity used here. We give the finite weak-to-regular bootstrap needed for this theorem, following the vorticity and div--curl route of Lemari\'e-Rieusset~\cite[Theorems~13.1--13.2, pp.~397--403]{LemarieRieusset2016}. The proof retains the dependence on the local gradient energy $E$ needed to start the vorticity energy estimate, and establishes joint continuity only for $\omega$, which is the trace used below.
\end{remark}

\begin{remark}[Dependence of the bounds]
Every constant $C$ in the theorem depends only on $(M,E)$ and the fixed radii in the proof. To make the dependence on $E=\|Du\|_{L^2(Q_1^-)}$ uniform in a family, take a larger cylinder $\mathcal Q=B_R(x_0)\times(t_0-T_-,t_0+T_+)\Subset U$ containing $\overline{Q_1^-(z_0)}$ in its interior, and choose $\psi\in C_c^\infty(\mathcal Q)$ with $\psi\ge0$ and $\psi=1$ on $\overline{Q_1^-(z_0)}$. Set $M_{\mathcal Q}=\operatorname*{ess\,sup}_{\mathcal Q}|u|$. With $f=0$, the local energy inequality (S4) and H\"older's inequality give
\[
 E^2\le\frac12\Bigl[M_{\mathcal Q}^2\bigl(\|\dt\psi\|_\infty+\|\Delta\psi\|_\infty\bigr)|\mathcal Q|
 +\|\nabla\psi\|_\infty\bigl(M_{\mathcal Q}^3|\mathcal Q|+2M_{\mathcal Q}|\mathcal Q|^{1/3}\|p\|_{L^{3/2}(\mathcal Q)}\bigr)\Bigr].
\]
Thus fixed larger-cylinder geometry and uniform bounds for $\sup_{\mathcal Q}|u|$ and $\|p\|_{L^{3/2}(\mathcal Q)}$ give a uniform $E$, and hence a uniform theorem constant when $M$ is also uniform. Finiteness of $E$ for each suitable solution follows from (S1) in Definition~\ref{def:sws}. The argument proves continuity up to the top face for $\omega$; it does not assert joint continuity of $u$ or $Du$ in time.
\end{remark}
\section{Stability of suitable weak solutions and the blow-up limit}\label{sec:blowup}

The hypotheses of Theorem~\ref{thm:ess-local} give energy bounds and an $L^\infty_tL^3_x$ bound, but they do not include the local energy inequality required by Definition~\ref{def:sws}. We first recover an energy identity and a representative continuous under spatial $L^{3/2}$ tests. We then fix a pressure decomposition on the original cylinder; its harmonic part has enough interior control to vanish under blow-up scaling.

\subsection{Energy equality and weak continuity}

\begin{lemma}[Local energy equality]\label{lem:lei-L4}
Suppose that $u$, $Du$, and $p$ satisfy the hypotheses of Theorem~\ref{thm:ess-local}, and set
\begin{equation}\label{eq:lei-L4-M}
M:=\esssup_{t\in(-1,0)}\norm{u(\cdot,t)}_{L^3(B_1)}<\infty\,.
\end{equation}
For every $r\in(0,1)$, $u\in L^4(B_r\times(-1,0))$ and, for every $\psi\in C_c^\infty(B_r\times(-1,0))$,
\begin{equation}\label{eq:local-energy-equality}
2\iint_{B_r\times(-1,0)}\abs{Du}^2\psi\,dx\,dt
=\iint_{B_r\times(-1,0)}\left[\abs{u}^2(\dt\psi+\Delta\psi)
 +(\abs{u}^2+2p)u\cdot\nabla\psi\right]dx\,dt\,.
\end{equation}
Consequently $(u,Du,p,0)$ is a suitable weak solution on $B_r\times(-1,0)$, in the sense of Definition~\ref{def:sws} with $q=3$.
\end{lemma}

\begin{proof}
For almost every $t\in(-1,0)$, each component of $u(\cdot,t)$ belongs to $H^1(B_r)$ with weak gradient given by the corresponding row of $Du(\cdot,t)$. The ball Sobolev--Poincar\'e inequality applied after subtracting the component means, followed by H\"older's inequality for those means, gives the scale-explicit estimate
\begin{equation}\label{eq:local-sobolev-L6}
\norm{u(\cdot,t)}_{L^6(B_r)}\le C\left(\norm{Du(\cdot,t)}_{L^2(B_r)}
 +r^{-1}\norm{u(\cdot,t)}_{L^2(B_r)}\right)\,,
\end{equation}
where $C$ is an absolute constant in dimension three for three components (CKN library). Interpolation between $L^3$ and $L^6$ gives
\begin{equation}\label{eq:local-interpolate-L4}
\int_{B_r}\abs{u(x,t)}^4\,dx\le\norm{u(\cdot,t)}_{L^3(B_r)}^2
 \norm{u(\cdot,t)}_{L^6(B_r)}^2\le CM^2\left(\norm{Du(\cdot,t)}_{L^2(B_r)}^2
 +r^{-2}\norm{u(\cdot,t)}_{L^2(B_r)}^2\right)\,.
\end{equation}
Integrating in time and using the energy bound in Theorem~\ref{thm:ess-local} proves $u\in L^4(B_r\times(-1,0))$. In particular $F:=u\otimes u$ belongs to $L^2$ on this cylinder.

Fix a compact subset of $B_r\times(-1,0)$ containing the support of $\psi$ in its interior, and choose a nonnegative smooth space-time mollifier $\rho$ of integral one, with $\rho_\epsilon(z)=\epsilon^{-4}\rho(z/\epsilon)$. For sufficiently small $\epsilon$, convolution is taken only where its translated support remains in $B_1\times(-1,0)$. Write $u_\epsilon=u*\rho_\epsilon$, $F_\epsilon=F*\rho_\epsilon$, and $p_\epsilon=p*\rho_\epsilon$ there. Spatial divergence commutes with this convolution, so $\nabla\cdot u_\epsilon=0$. The local $L^s$ convergence of mollification for finite $s$ follows from translation continuity in $L^s$: $\norm{g*\rho_\epsilon-g}_{L^s}\le\int\rho_\epsilon(h)\norm{g(\cdot-h)-g}_{L^s}\,dh\to0$ after restricting to a slightly larger compact subset.

Convolving the weak momentum identity (S3) gives
\begin{equation}\label{eq:mollified-momentum}
\dt u_\epsilon+\nabla\cdot F_\epsilon-\Delta u_\epsilon+\nabla p_\epsilon=0\,.
\end{equation}
Set $R_\epsilon=F_\epsilon-u_\epsilon\otimes u_\epsilon$. Testing \eqref{eq:mollified-momentum}, rewritten with $u_\epsilon\otimes u_\epsilon$ on the left and $-\nabla\cdot R_\epsilon$ on the right, against $2\psi u_\epsilon$ and integrating by parts yields
\begin{align}\label{eq:mollified-energy}
2\iint\abs{Du_\epsilon}^2\psi
={}&\iint\left[\abs{u_\epsilon}^2(\dt\psi+\Delta\psi)
 +(\abs{u_\epsilon}^2+2p_\epsilon)u_\epsilon\cdot\nabla\psi\right]\notag\\
&+2\iint R_\epsilon:Du_\epsilon\,\psi
 +2\iint R_{\epsilon,ij}(u_\epsilon)_i\partial_j\psi\,.
\end{align}
All these integrations by parts have no boundary terms because $\psi$ is compactly supported in the region where the convolutions are defined. The pressure term reduces to $p_\epsilon u_\epsilon\cdot\nabla\psi$ because $\nabla\cdot u_\epsilon=0$.

The approximate-identity convergence gives $u_\epsilon\to u$ in $L^4$, $Du_\epsilon\to Du$ in $L^2$, and $F_\epsilon\to F$ in $L^2$ on the support of $\psi$. Also $u_\epsilon\otimes u_\epsilon\to u\otimes u$ in $L^2$, so $R_\epsilon\to0$ in $L^2$. Thus the two error integrals in \eqref{eq:mollified-energy} tend to zero by Cauchy--Schwarz. Since $p_\epsilon\to p$ in $L^{3/2}$ and $u_\epsilon\to u$ in $L^3$ on this finite-measure set, $p_\epsilon u_\epsilon\to pu$ in $L^1$. The remaining terms converge in $L^1$ by the stated $L^4$ and $L^2$ convergences. Passing to the limit in \eqref{eq:mollified-energy} proves \eqref{eq:local-energy-equality}.

The energy, pressure, measurability, and slice weak-gradient clauses of Definition~\ref{def:sws} on $B_r\times(-1,0)$ follow by restricting the hypotheses of Theorem~\ref{thm:ess-local} to compactly contained boxes. Conditions (S2) and (S3) restrict to tests supported there, and (S4) follows from \eqref{eq:local-energy-equality} for nonnegative tests. This proves suitability with $q=3$ and $f=0$.
\end{proof}

\begin{lemma}[Weakly continuous $L^3$ representative]\label{lem:weak-cont-L3}
Under the hypotheses of Theorem~\ref{thm:ess-local}, there is a map
\begin{equation}\label{eq:weak-cont-domain}
\bar u\colon[-(3/4)^2,0]\longrightarrow L^3(B_{3/4};\R^3)
\end{equation}
that agrees with $u(\cdot,t)$ for almost every $t$ and satisfies, for every $w\in L^{3/2}(B_{3/4};\R^3)$, continuity of
\begin{equation}\label{eq:weak-cont-pairing}
t\longmapsto\int_{B_{3/4}}\bar u(x,t)\cdot w(x)\,dx\,,
\end{equation}
together with
\begin{equation}\label{eq:weak-cont-bound}
\sup_{-(3/4)^2\le t\le0}\norm{\bar u(\cdot,t)}_{L^3(B_{3/4})}\le M\,,
\end{equation}
where $M$ is defined in \eqref{eq:lei-L4-M}.
\end{lemma}

\begin{proof}
Fix $\varphi\in C_c^\infty(B_{3/4};\R^3)$. Testing (S3) with $\eta(t)\varphi(x)$ shows that the distributional derivative of $t\mapsto\int u(x,t)\cdot\varphi(x)\,dx$ is the locally integrable function
\begin{equation}\label{eq:weak-cont-derivative}
\int_{B_{3/4}}\left[u_i u_j\partial_j\varphi_i
 -Du_{ij}\partial_j\varphi_i+p\,\nabla\cdot\varphi\right](x,t)\,dx\,.
\end{equation}
The first term has absolute value at most $M^2\norm{\nabla\varphi}_{L^3(B_{3/4})}$ for almost every $t$. For $s<t$, Cauchy--Schwarz in space-time bounds the increment of the second term by $\norm{Du}_{L^2(Q)}\norm{\nabla\varphi}_{L^2(B_{3/4})}|t-s|^{1/2}$. H\"older's inequality bounds the integrated pressure term by $\norm{p}_{L^{3/2}(Q)}\norm{\nabla\cdot\varphi}_{L^3(B_{3/4})}|t-s|^{1/3}$. Therefore the scalar pairing has an absolutely continuous representative, denoted by $\ell_\varphi$, on $[-(3/4)^2,0]$, and this representative obeys
\begin{align}\label{eq:weak-cont-modulus}
\abs{\ell_\varphi(t)-\ell_\varphi(s)}
\le{}&M^2\norm{\nabla\varphi}_{L^3(B_{3/4})}|t-s|
 +\norm{Du}_{L^2(Q)}\norm{\nabla\varphi}_{L^2(B_{3/4})}|t-s|^{1/2}\notag\\
&+\norm{p}_{L^{3/2}(Q)}\norm{\nabla\cdot\varphi}_{L^3(B_{3/4})}|t-s|^{1/3}\,.
\end{align}
The value at $t=0$ is the limit from below; the lower endpoint is interior to the original time interval and the same absolute-continuity argument applies there.

Choose a full-measure set of times on which $u(\cdot,t)\in L^3(B_{3/4})$ with norm at most $M$ and on which its pairings agree with the absolutely continuous representatives for a countable dense family of smooth tests. At each time in this set, Hölder's inequality bounds the pairing by $M\norm{w}_{L^{3/2}}$. Density extends this bound to every $w\in L^{3/2}(B_{3/4};\R^3)$, and the duality $(L^{3/2})^*=L^3$ represents the time-slice pairing by an element of $L^3$ with norm at most $M$.

For any $t$ in the closed interval, choose good times converging to $t$. Since $B_{3/4}$ has finite measure, H\"older's inequality makes these slices uniformly bounded in $L^2(B_{3/4})$. Their pairings against the countable $L^2$-dense family of smooth tests converge by \eqref{eq:weak-cont-modulus}; these limits determine an $L^2$ element $\bar u(\cdot,t)$ and weak $L^2$ convergence of the slices to it. Passing the slice bound to this limit gives $|\int_{B_{3/4}}\bar u(\cdot,t)\cdot w|\le M\norm{w}_{L^{3/2}}$ for every $L^2$ test $w$. Applying the $L^3$ norming estimate to $w_N=\min(\|\bar u(\cdot,t)\|,N)\bar u(\cdot,t)\in L^2$ and letting $N\to\infty$ gives $\bar u(\cdot,t)\in L^3$ with norm at most $M$. The smooth-test pairings then extend by density to every $L^{3/2}$ test; on the full-measure set they agree with those of $u$, so this representative agrees with $u$ almost everywhere and satisfies \eqref{eq:weak-cont-bound}.
\end{proof}

\subsection{The pressure split}

For the pressure estimates below we use the whole-space pressure operator of Definition~\ref{def:riesz-pressure}, with the sign convention fixed by (S3). The input tensor is extended by zero outside $B_1$ before applying that operator.

\begin{lemma}[Pressure split and rescaling]\label{lem:pressure-split}
Suppose that $u$, $Du$, and $p$ satisfy the hypotheses of Theorem~\ref{thm:ess-local}, and let $M$ be as in \eqref{eq:lei-L4-M}. For almost every $t\in(-1,0)$ define
\begin{equation}\label{eq:pressure-split-defs}
F_{ij}(x,t):=\mathbf1_{B_1}(x)u_i(x,t)u_j(x,t),\qquad
p_1(\cdot,t):=P[F](\cdot,t),\qquad
p_2:=p-p_1\quad\hbox{on }B_1\times(-1,0)\,.
\end{equation}
Then $p_1\in L^\infty((-1,0);L^{3/2}(\R^3))$, $p_2\in L^{3/2}(B_1\times(-1,0))$, and for almost every $t$ the slice $p_2(\cdot,t)$ is weakly harmonic in $B_1$. There is an absolute constant $C$ such that
\begin{align}\label{eq:pressure-split-bounds}
\esssup_{t\in(-1,0)}\norm{p_1(\cdot,t)}_{L^{3/2}(\R^3)}&\le CM^2\,,\notag\\
\int_{-1}^0\norm{p_2(\cdot,t)}_{L^\infty(B_{3/4})}^{3/2}\,dt
&\le C\left(\norm{p}_{L^{3/2}(Q)}^{3/2}+M^3\right)\,.
\end{align}
Extend $u$ by zero outside $Q$, extend $p_1$ by zero in time, and extend $p_2$ by zero outside $B_1\times(-1,0)$; denote these extensions by tildes. Given $z_0=(x_0,t_0)\in\overline{B_{1/2}}\times[-1/4,0]$ and $R>0$, set
\begin{equation}\label{eq:pressure-split-rescalings}
\begin{aligned}
v^R(x,t)&:=R\widetilde u(x_0+Rx,t_0+R^2t),\\
p_1^R(x,t)&:=R^2\widetilde p_1(x_0+Rx,t_0+R^2t),\\
p_2^R(x,t)&:=R^2\widetilde p_2(x_0+Rx,t_0+R^2t)\,.
\end{aligned}
\end{equation}
For almost every $t\in\R$ these satisfy
\begin{align}\label{eq:pressure-split-scaling}
\int_{\R^3}\abs{v^R(x,t)}^3\,dx
 &=\int_{\R^3}\abs{\widetilde u(x,t_0+R^2t)}^3\,dx\,,\notag\\
\int_{\R^3}\abs{p_1^R(x,t)}^{3/2}\,dx
 &=\int_{\R^3}\abs{\widetilde p_1(x,t_0+R^2t)}^{3/2}\,dx\,.
\end{align}
For every bounded measurable $\Omega\subset\R^3$, all sufficiently small $R$ also satisfy $x_0+R\Omega\subset B_{3/4}$ and
\begin{equation}\label{eq:pressure-split-p2-scaling}
\int_{\R}\norm{p_2^R(\cdot,t)}_{L^\infty(\Omega)}^{3/2}\,dt
=R\int_{\R}\norm{\widetilde p_2(\cdot,s)}_{L^\infty(x_0+R\Omega)}^{3/2}\,ds
\le CR\left(\norm{p}_{L^{3/2}(Q)}^{3/2}+M^3\right)\,.
\end{equation}
In particular, $p_2^R\to0$ in $L^{3/2}(\R;L^\infty(\Omega))$ as $R\downarrow0$.
\end{lemma}

\begin{proof}
For almost every $t$, the zero-extended tensor $F(\cdot,t)$ belongs to $L^{3/2}(\R^3)$ and
\begin{equation}\label{eq:pressure-split-cz}
\sum_{i,j=1}^3\norm{F_{ij}(\cdot,t)}_{L^{3/2}(\R^3)}
\le 9\norm{u(\cdot,t)}_{L^3(B_1)}^2\,.
\end{equation}
The slice estimate and joint measurability in Definition~\ref{def:riesz-pressure} give the first bound in \eqref{eq:pressure-split-bounds}. Since $p\in L^{3/2}(Q)$ and $p_1$ is in that space on $B_1\times(-1,0)$, their difference $p_2$ also belongs to $L^{3/2}$ there.

For $\zeta\in C_c^\infty(Q)$, use $\nabla\zeta$ as the vector test in (S3). The time term vanishes by (S2) applied to $\partial_t\zeta$. In the viscous term, integration by parts in $x_j$ gives $-\iint u_i\partial_i\Delta\zeta$, which vanishes by (S2) applied to $\Delta\zeta$. These integrations have no boundary terms because $\zeta$ is compactly supported in $Q$. The remaining terms yield
\begin{equation}\label{eq:pressure-poisson}
\iint_Q p\,\Delta\zeta\,dx\,dt
=-\sum_{i,j=1}^3\iint_Q u_i u_j\partial_i\partial_j\zeta\,dx\,dt\,.
\end{equation}
By Definition~\ref{def:riesz-pressure}, $-\Delta p_1=\sum_{i,j}\partial_i\partial_jF_{ij}$ in spatial distributions. Since $F=u\otimes u$ on $B_1$, subtracting this identity from \eqref{eq:pressure-poisson} gives $\Delta_xp_2=0$ in distributions on $B_1\times(-1,0)$. To obtain the slice assertion, test with products $\zeta(x)\eta(t)$, first for a countable dense family of $\zeta\in C_c^\infty(B_1)$; Fubini's theorem gives a common full-measure set of times on which the harmonic identity holds for that family. The $L^{3/2}$ slice bound for $p_2$ extends the identity to every spatial test function, so $p_2(\cdot,t)$ is weakly harmonic for almost every $t$.

The CKN weakly harmonic interior estimate gives, for almost every $t$,
\begin{equation}\label{eq:pressure-harmonic-sup}
\norm{p_2(\cdot,t)}_{L^\infty(B_{3/4})}
\le C_{17}\norm{p_2(\cdot,t)}_{L^{3/2}(B_1)}\,,
\end{equation}
where $C_{17}$ is absolute in dimension three (CKN library). The triangle inequality and the first estimate in~\eqref{eq:pressure-split-bounds} give
\begin{align*}
\int_{-1}^0\norm{p_2(\cdot,t)}_{L^\infty(B_{3/4})}^{3/2}\,dt
&\le C\int_{-1}^0\int_{B_1}\left(\abs{p(x,t)}^{3/2}+\abs{p_1(x,t)}^{3/2}\right)dx\,dt\\
&\le C\left(\norm{p}_{L^{3/2}(Q)}^{3/2}+M^3\right)\,,
\end{align*}
which proves the second estimate in \eqref{eq:pressure-split-bounds}.

The first scaling identity in \eqref{eq:pressure-split-scaling} follows by changing variables $y=x_0+Rx$ and using $|R\widetilde u|^3dx=|\widetilde u|^3dy$. The second follows from the same change of variables and $|R^2\widetilde p_1|^{3/2}dx=|\widetilde p_1|^{3/2}dy$. For the last identity, the same spatial change and then $s=t_0+R^2t$ give the factor $R^3R^{-2}=R$ in \eqref{eq:pressure-split-p2-scaling}. Because $x_0\in\overline{B_{1/2}}$ and $\Omega$ is bounded, $x_0+R\Omega\subset B_{3/4}$ for all sufficiently small $R$, with the threshold depending only on $\sup_{x\in\Omega}|x|$. The second estimate of \eqref{eq:pressure-split-bounds} then proves \eqref{eq:pressure-split-p2-scaling}; its right side tends to zero with $R$.
\end{proof}

\begin{remark}[Departure from \cite{ESS2003}]\label{rem:pb1-departure}
ESS define their local $p_1$ by the weak Dirichlet problem in (3.11)--(3.12) on printed pp.~222--223 of~\cite{ESS2003}, and take $p_2$ to be the resulting harmonic remainder; here $p_1$ is the whole-space Riesz pressure of the zero extension of $u\otimes u$. The new $p_2$ is harmonic on $B_1$, and the needed interior bound and rescaled vanishing are proved directly from this fixed decomposition and the assumed $L^{3/2}(Q)$ pressure norm. Thus the proof does not transfer (3.19) on printed p.~223 of~\cite{ESS2003} to a pressure remainder redefined after rescaling. ESS derive the time continuity in (3.8) through the stronger estimate (3.10) on printed p.~222; here the pairing modulus follows directly from (S3), the energy bound, and the $L^{3/2}$ pressure bound. In the local energy equality, the full pressure remains in $L^{3/2}$: pairing it with the mollified velocity converges in $L^1$ by the $L^{3/2}$--$L^3$ H\"older inequality, so no $L^2$ hypothesis on $p$ is imposed.
\end{remark}

\subsection{The blow-up limit}\label{subsec:blowup-limit}

The compactness argument needs uniform control on each fixed region of the rescaled past. The local energy inequality supplies the gradient bound; the pressure split gives the pressure bound and, after blow-up, identifies its limit with the whole-space pressure of the limiting velocity.

We use the open past cylinder $Q_r^-(z_0)$ fixed in Section~\ref{sec:notation}.
The open top is useful when \(t_0=0\), since all rescaled times then remain in the original time interval.

\begin{proposition}[Blow-up limit]\label{prop:blowup-limit}
Suppose \(u_{\mathrm{orig}},Du_{\mathrm{orig}},p_{\mathrm{orig}}\) satisfy the hypotheses of Theorem~\ref{thm:ess-local} on \(D=B_1\times(-1,0)\). Let
\begin{equation*}
 K=\overline{B_{1/2}}\times[-1/4,0],
\end{equation*}
and let \(z_0=(x_0,t_0)\in K\) be a bad point in the sense of Definition~\ref{def:good-point}, with its top-face interpretation. Let \(R_k\to0\), and discard finitely many terms so that \(Q_{R_k}^-(z_0)\subset D\) for every \(k\).

For almost every \(t\in(-1,0)\), let
\begin{equation*}
 F(x,t)=\mathbf 1_{B_1}(x)\,u_{\mathrm{orig}}(x,t)\otimes u_{\mathrm{orig}}(x,t),\qquad
 p_1(\cdot,t)=P[F(\cdot,t)],\qquad p_2=p_{\mathrm{orig}}-p_1\quad\hbox{on }D,
\end{equation*}
as in Lemma~\ref{lem:pressure-split}. Extend \(\tilde u_{\mathrm{orig}}=\mathbf1_Du_{\mathrm{orig}}\), \(\widetilde{Du}_{\mathrm{orig}}=\mathbf1_DDu_{\mathrm{orig}}\), and \(\tilde p_{\mathrm{orig}}=\mathbf1_Dp_{\mathrm{orig}}\) by zero to \(\mathbb R^3\times\mathbb R\). Extend \(p_1\) by zero in time only, retaining its whole-space spatial values, and extend \(p_2\) by zero outside \(D\); call these extensions \(\tilde p_1,\tilde p_2\). Define
\begin{align*}
 v^k(x,t)&=R_k\tilde u_{\mathrm{orig}}(x_0+R_kx,t_0+R_k^2t),&
 Dv^k(x,t)&=R_k^2\widetilde{Du}_{\mathrm{orig}}(x_0+R_kx,t_0+R_k^2t),\\
 p^k(x,t)&=R_k^2\tilde p_{\mathrm{orig}}(x_0+R_kx,t_0+R_k^2t),&
 p_1^k(x,t)&=R_k^2\tilde p_1(x_0+R_kx,t_0+R_k^2t),\\
 &&p_2^k(x,t)&=R_k^2\tilde p_2(x_0+R_kx,t_0+R_k^2t).
\end{align*}
Let \(D_k=\{(x,t):(x_0+R_kx,t_0+R_k^2t)\in D\}\). Then \(p_1^k=P[v^k\otimes v^k]\), and \(p^k=p_1^k+p_2^k\) on \(D_k\). For time-slice assertions, use the \(L^3(B_{3/4})\) equivalence class \(\bar u(\cdot,\tau)\) from Lemma~\ref{lem:weak-cont-L3} whenever \(\tau=t_0+R_k^2t\in[-(3/4)^2,0]\) and the spatial region under consideration rescales into \(B_{3/4}\). The weak-continuity lemma identifies its smooth-test pairings with the absolutely continuous representatives from the momentum equation. This choice agrees with the zero-extended spacetime field almost everywhere and applies throughout each fixed bounded cylinder \(B_R\times[a,0]\) for all sufficiently large \(k\); in particular it specifies the terminal slice at \(t=0\).

There is a subsequence, still denoted \(k\), and fields \(u,Du,q\) on \(\mathbb R^3\times(-\infty,0)\) with the following properties.
\begin{enumerate}[label=(\alph*)]
\item For every bounded cylinder \(C=B_R\times(a,0)\), with \(R<\infty\) and \(a<0\), a constant depending only on \(C\) and the norms in Theorem~\ref{thm:ess-local} satisfies
\begin{equation*}
 \iint_C\bigl(|p^k|^{3/2}+|Dv^k|^2\bigr)\,dx\,dt\le C_C
\end{equation*}
for all sufficiently large \(k\).

\item On each such cylinder, \(v^k\) has the pairing modulus required in \cite[Lemma~11.1]{CKNLean}: for every compact \(K_0\subset\mathbb R^3\), every finite interval \(J=[a,0]\) with \(a<0\), every \(w\in C_c^\infty(\mathbb R^3;\mathbb R^3)\) supported in \(K_0\), and all \(s,t\in J\),
\begin{align*}
\left|\int_{\mathbb R^3}(v^k(x,t)-v^k(x,s))\cdot w(x)\,dx\right|
\le C_{K_0,J}\Bigl(&|t-s|\bigl(\|\Delta w\|_{L^2(\mathbb R^3)}
                  +\|\nabla w\|_{L^\infty(\mathbb R^3)}\bigr)\\
 &+|t-s|^{1/3}\|\operatorname{div}w\|_{L^\infty(\mathbb R^3)}\Bigr),
\end{align*}
where the constants are independent of \(k,s,t,w\) for all sufficiently large \(k\). For \(J\subset(-\infty,0)\) this is the hypothesis of \cite[Lemma~11.1]{CKNLean}; the same estimate extends to \(t=0\) using the trace just defined. Also the \(L^{10/3}(C)\) norms are uniformly bounded for all sufficiently large \(k\) on every bounded cylinder \(C=B_R\times(a,0)\).

\item The subsequence satisfies
\begin{equation*}
 v^k\longrightarrow u\quad\hbox{strongly in }L^3(C)
\end{equation*}
for every bounded cylinder \(C=B_R\times(a,0)\), and \(v^k(\cdot,t)\rightharpoonup u(\cdot,t)\) weakly in \(L^2_{\mathrm{loc}}(\mathbb R^3)\) for every \(t<0\). Also \(Dv^k\rightharpoonup Du\) weakly in \(L^2(C)\) for every such \(C\), where \(Du\) is the weak spatial gradient of \(u\). Moreover,
\begin{align*}
 u&\in L^\infty((-\infty,0);L^3(\mathbb R^3)),\\
 p_1^k&\longrightarrow q\quad\hbox{strongly in }L^{3/2}(C),\\
 p_2^k&\longrightarrow0\quad\hbox{strongly in }L^{3/2}(C),\\
 p^k&\longrightarrow q\quad\hbox{strongly in }L^{3/2}(C).
\end{align*}

\item The limit \((u,Du,q,0)\) is a suitable weak solution with exponent \(3\) on every \(B_\rho\times(a,0)\), where \(\rho<\infty\) and \(a<0\). The limiting pressure is the whole-space associated pressure
\begin{equation*}
 q=P[u\otimes u]\quad\hbox{on }\mathbb R^3\times(-\infty,0).
\end{equation*}

\item There is a terminal slice \(u(\cdot,0)\in L^2_{\mathrm{loc}}(\mathbb R^3)\) such that \(v^k(\cdot,0)\rightharpoonup u(\cdot,0)\) weakly in \(L^2_{\mathrm{loc}}(\mathbb R^3)\), and
\begin{equation*}
 \int_{B_1(x)}|u(y,0)|^2\,dy=0\qquad\hbox{for every }x\in\mathbb R^3.
\end{equation*}

\item For every \(k\),
\begin{equation*}
 \iint_{Q_1^-(0)}\bigl(|v^k|^3+|p^k|^{3/2}\bigr)\,dx\,dt\ge\varepsilon_0/8.
\end{equation*}
If \(u=0\) almost everywhere on \(\mathbb R^3\times(-T,0)\) for some \(T>1\), then
\begin{equation*}
 \lim_{k\to\infty}\iint_{Q_1^-(0)}
       \bigl(|v^k|^3+|p^k|^{3/2}\bigr)\,dx\,dt=0.
\end{equation*}
\end{enumerate}
\end{proposition}

\begin{proof}
Put \(M=\esssup_{t\in(-1,0)}\|u_{\mathrm{orig}}(\cdot,t)\|_{L^3(B_1)}\). Lemma~\ref{lem:lei-L4} makes \((u_{\mathrm{orig}},Du_{\mathrm{orig}},p_{\mathrm{orig}},0)\) suitable on \(B_{3/4}\times(-1,0)\), while Lemma~\ref{lem:weak-cont-L3} supplies the all-time local \(L^3\) representative up to \(t=0\). For each fixed bounded cylinder in the rescaled variables, its inverse image lies in \(B_{3/4}\times(-1,0)\) for all sufficiently large \(k\), except for its top face, which has measure zero. Since \(t_0+R_k^2t\le t_0\le0\) whenever \(t\le0\), these rescalings use only past times.

\paragraph{Local energy and pressure bounds.}
The velocity scaling and Lemma~\ref{lem:pressure-split} give the bound $M$ for the global spatial $L^3$ norm of $v^k(\cdot,t)$ for almost every rescaled time. For almost every such $t$, both $p_1^k(\cdot,t)$ and $P[v^k(\cdot,t)\otimes v^k(\cdot,t)]$ belong to $L^{3/2}(\mathbb R^3)$ and solve the Poisson equation with right-hand side $\sum_{i,j=1}^3\partial_i\partial_j(v_i^kv_j^k)$. Their difference is harmonic, hence smooth, and belongs to $L^{3/2}(\mathbb R^3)$. The mean-value property gives $|h(x)|\le |B_R|^{-2/3}\|h\|_{L^{3/2}(\mathbb R^3)}$ for this difference $h$ and every $R>0$; letting $R\to\infty$ shows that $h=0$. Hence $p_1^k=P[v^k\otimes v^k]$ on almost every slice, and the $L^{3/2}$ boundedness of $P$ gives
\begin{equation*}
 \|p_1^k(\cdot,t)\|_{L^{3/2}(\mathbb R^3)}\le C M^2.
\end{equation*}
For every bounded spatial set \(\Omega\),
\begin{equation*}
 \int_{\mathbb R}\|p_2^k(\cdot,t)\|_{L^\infty(\Omega)}^{3/2}\,dt
 \le C R_k\bigl(\|p_{\mathrm{orig}}\|_{L^{3/2}(D)}^{3/2}+M^3\bigr).
\end{equation*}
Here and below \(C\) may change from line to line and depends only on the displayed bounded sets and the fixed data norms. In particular, \(p^k\) is uniformly bounded in \(L^{3/2}(C)\) on each bounded cylinder \(C\). For a fixed bounded spatial set \(\Omega\), Hölder's inequality and the velocity scaling give, uniformly for \(t\) in any fixed interval \([a,0]\),
\begin{equation}\label{eq:blowup-slice-bound}
 \int_\Omega |v^k(x,t)|^2\,dx
 \le R_k^{-1}\int_{x_0+R_k\Omega}|\bar u(y,t_0+R_k^2t)|^2\,dy
 \le |\,\Omega\,|^{1/3}M^2.
\end{equation}
This holds for all sufficiently large \(k\), since \(x_0+R_k\Omega\subset B_{3/4}\) and the time arguments lie in \([-(3/4)^2,0]\).

Fix \(C=B_R\times(a,0)\). Choose a nonnegative spatial cutoff \(\zeta\in C_c^\infty(B_{R+1})\) that equals one on \(B_R\). For \(b<0\), choose time cutoffs \(\eta_b\in C_c^\infty((a-1,0))\), \(0\le\eta_b\le1\), equal to one on \([a,b]\), with \(\|\eta_b'\|_{L^1}\le2\), uniformly as \(b\uparrow0\). Apply the local energy inequality to \(\psi(x,t)=\zeta(x)\eta_b(t)\) for the rescaled suitable solution. The terms with \(\nabla\zeta\) are bounded by the local \(L^\infty_tL^3_x\) bound and the local \(L^{3/2}\) pressure bound; the term with \(\eta_b'\) is bounded by \(2\sup_t\int_{\operatorname{supp}\zeta}|v^k|^2\). The term with \(\Delta\zeta\) has the same slice \(L^2\) bound. These bounds do not depend on \(b\) or \(k\). Thus \(\iint_{B_R\times(a,b)}|Dv^k|^2\le C_C\); letting \(b\uparrow0\) proves (a), including its top-face endpoint.

The weak momentum identity (S3), tested with a spatial field \(w\) and a scalar time test, gives the absolutely continuous pairing identity
\begin{equation*}
 \int_{\mathbb R^3}(v^k(x,t)-v^k(x,s))\cdot w(x)\,dx
 =\int_s^t\int_{\mathbb R^3}
 \bigl(v^k\cdot\Delta w+(v^k\otimes v^k):\nabla w
                         +p^k\operatorname{div}w\bigr)\,dx\,d\tau.
\end{equation*}
The integration by parts in the diffusion term has no spatial boundary contribution because \(w\) is compactly supported. The first term is bounded by \(C|t-s|\|\Delta w\|_2\) using the local slice \(L^2\) bound. The second is bounded by \(C|t-s|\|\nabla w\|_\infty\) using the \(L^\infty_tL^3_x\) bound. For the last term, the local \(L^{3/2}\) pressure bound and Hölder's inequality give
\begin{equation*}
 \int_s^t\int_{\operatorname{supp}w}|p^k|\,|\operatorname{div}w|
 \le C_{K_0,J}|t-s|^{1/3}\|\operatorname{div}w\|_\infty.
\end{equation*}
The identity first holds for \(s,t<0\). The trace from Lemma~\ref{lem:weak-cont-L3} and the integrability of the right side let \(t\uparrow0\), proving the same modulus on \([a,0]\).

For the local \(L^{10/3}\) bound, choose a spatial cutoff \(\eta\) supported in a bounded neighborhood of \(B_R\) and equal to one there. The three-dimensional Gagliardo--Nirenberg inequality applied to \(\eta v^k(\cdot,t)\), extended by zero, gives
\begin{equation*}
 \|\eta v^k(\cdot,t)\|_{10/3}^{10/3}
 \le C\|\eta v^k(\cdot,t)\|_2^{4/3}
       \|\nabla(\eta v^k)(\cdot,t)\|_2^2.
\end{equation*}
The local slice bound and (a), integrated over \((a,0)\), give the asserted uniform \(L^{10/3}\) bound.

For each integer $m\ge1$, choose $N(m)$ so that the local energy, pairing-modulus, and $L^{10/3}$ bounds above hold on $\overline{B}_{m+1}\times[-(m+1),0]$ for every $k\ge N(m)$. Choose a strictly increasing sequence $\nu(j)$ with $\nu(j)\ge\max_{m\le j}N(m)$. Let $\chi_j$ be a smooth spatial cutoff equal to one on $\overline{B}_j$ and supported in $B_{j+1}$, and let $\theta_j$ be a smooth time ramp vanishing before $-(j+1)$ and equal to one on $[-j,0]$. Apply \cite[Lemma~11.1]{CKNLean} once on $\mathbb R^3\times(-\infty,0)$ to the cutoff sequence $\widehat v^j=\theta_j\chi_jv^{\nu(j)}$. The preceding bounds and the cutoff derivative estimates give the required uniform hypotheses on each compact cylinder; the finitely many early indices at each stage are controlled by their individual moduli. Since $\widehat v^j=v^{\nu(j)}$ on any fixed cylinder for all sufficiently large $j$, the resulting all-slice weak $L^2_{\rm loc}$ convergence, local strong $L^2$ convergence, and local weak-gradient convergence transfer to the original rescaled fields. \cite[Corollary~11.3]{CKNLean} then gives strong $L^3$ convergence on every compact subcylinder below $t=0$. Fix \(R<\infty\) and \(a<0\). For sufficiently small \(\delta>0\), the uniform \(L^{10/3}\) bound and H\"older's inequality give
\begin{equation*}
 \|v^k\|_{L^3(B_R\times(a,a+\delta))}^3
 \le (|B_R|\delta)^{1/10}\|v^k\|_{L^{10/3}(B_R\times(a,a+\delta))}^3,
\end{equation*}
and the identical estimate holds on \(B_R\times(-\delta,0)\) and for \(u\) by Fatou's lemma. Both time-boundary strips therefore have \(L^3\) norms tending to zero uniformly in \(k\). On \(B_R\times(a+\delta,-\delta)\) the strong convergence follows from \cite[Corollary~11.3]{CKNLean}, since \(\overline{B_R}\times[a+\delta,-\delta]\) is compactly contained in a slightly larger space-time cylinder. Combining the interior convergence with the two strip estimates proves strong \(L^3\) convergence on \(B_R\times(a,0)\). The local weak-gradient convergence upgrades to weak convergence in \(L^2(B_R\times(a,0))\): (a) gives a uniform \(L^2\) bound on that cylinder, and compactly supported tests in its interior are dense in \(L^2\), so the local weak limit determines the weak limit on the whole cylinder. The distributional gradient identity passes to the limit by the strong local \(L^2\) convergence and identifies \(Du\) as the weak spatial gradient of \(u\).

By strong local $L^3$ convergence, take a diagonal subsequence of the zero-extended fields that converges almost everywhere on every bounded space-time cylinder. Fubini's theorem gives a full-measure set of times on which this subsequence converges almost everywhere in space; intersect it with the common full-measure set on which every slice $v^k(\cdot,t)$ has global $L^3$ norm at most $M$. For each such time, Fatou's lemma on each ball gives the same bound for $u(\cdot,t)$ there. Monotone convergence over an increasing sequence of balls yields $\|u(\cdot,t)\|_{L^3(\mathbb R^3)}\le M$, and hence $u\in L^\infty((-\infty,0);L^3)$.

\paragraph{The pressure limit.}
Set \(F_k=v^k\otimes v^k\) and \(F=u\otimes u\). The just-proved strong \(L^3\) convergence gives \(F_k\to F\) strongly in \(L^{3/2}(B_L\times(a,0))\) for each fixed \(L<\infty\) and \(a<0\). Fix \(R<L/2\). The near-field pressure satisfies
\begin{equation*}
 \|P[\mathbf1_{B_L}(F_k-F)]\|_{L^{3/2}(B_R\times(a,0))}
 \le C\|F_k-F\|_{L^{3/2}(B_L\times(a,0))}\longrightarrow0
\end{equation*}
by the \(L^{3/2}\) bound for the double Riesz transforms in Definition~\ref{def:riesz-pressure}. Away from the support, their kernel is the second derivative of the Newtonian kernel and obeys \(|K_{ij}(x-y)|\le C|x-y|^{-3}\). Thus, for \(|x|\le R\), \(L\ge2R\), and almost every \(t\),
\begin{align*}
 |P[\mathbf1_{\mathbb R^3\setminus B_L}F_k](x,t)|
 &\le C\int_{|y|>L}|x-y|^{-3}|F_k(y,t)|\,dy\\
 &\le C\bigl\||x-\cdot|^{-3}\bigr\|_{L^3(\mathbb R^3\setminus B_L)}
          \|F_k(\cdot,t)\|_{L^{3/2}(\mathbb R^3)}
 \le C L^{-2}M^2.
\end{align*}
The same estimate holds for \(F\). Taking the \(L^{3/2}\) norm on \(B_R\times(a,0)\), the far-field contribution is at most \(C|B_R|^{2/3}|a|^{2/3}L^{-2}M^2\). First let \(k\to\infty\) at fixed \(L\), then let \(L\to\infty\). This proves \(p_1^k=P[F_k]\to P[F]\) strongly in \(L^{3/2}(B_R\times(a,0))\). Lemma~\ref{lem:pressure-split} gives \(p_2^k\to0\) in \(L^{3/2}_tL^\infty_x\) on bounded spatial sets, so \(p_2^k\to0\) in \(L^{3/2}(B_R\times(a,0))\) as well. Since \(p^k=p_1^k+p_2^k\) on every fixed cylinder for all sufficiently large \(k\), the pressure limit is \(q=P[u\otimes u]\).

\paragraph{Suitability of the limit.}
Fix \(\rho<\infty\) and \(a<0\). On every open box compactly contained in \(B_\rho\times(a,0)\), its inverse image lies in \(B_{3/4}\times(-1,0)\) for all sufficiently large \(k\). Restriction by \cite[Lemma~11.5]{CKNLean} and the Navier--Stokes rescaling \(v^k=R_ku_{\mathrm{orig}}\), \(Dv^k=R_k^2Du_{\mathrm{orig}}\), \(p^k=R_k^2p_{\mathrm{orig}}\) make the tail of the sequence suitable there with exponent \(3\) and zero force. On each such box, (a) and the slice bound \eqref{eq:blowup-slice-bound} give the uniform essential-supremum \(L^2\) and space-time gradient bounds; (b)--(c) give strong \(L^3\), weak \(L^2\) gradients, and strong local \(L^{3/2}\) pressure convergence. Apply \cite[Theorem~11.4]{CKNLean} on each box to the corresponding tail. These conclusions for all compactly contained boxes prove suitability of \((u,Du,q,0)\) on \(B_\rho\times(a,0)\).

\paragraph{The terminal slice.}
For a fixed compact spatial set and a smooth compactly supported test \(w\), the modulus in (b) is uniform up to \(t=0\). Since the pairings converge for every \(t<0\), it follows that the sequence \(\int v^k(x,0)\cdot w(x)\,dx\) is Cauchy: for \(\delta>0\), compare each endpoint pairing with its value at \(-\delta\), use (b) twice, pass \(k,\ell\to\infty\) at \(-\delta\), and then let \(\delta\downarrow0\). The local \(L^2\) bounds make the limiting functionals bounded on \(L^2\); density and the Riesz representation theorem give a compatible slice \(u(\cdot,0)\in L^2_{\mathrm{loc}}\) and weak \(L^2_{\mathrm{loc}}\) convergence of \(v^k(\cdot,0)\) to it. The same modulus shows that this is the weak trace of \(u\) from times below zero.

Fix \(x\in\mathbb R^3\). Since \(t_0\in[-1/4,0]\), Lemma~\ref{lem:weak-cont-L3} gives \(\bar u(\cdot,t_0)\in L^3(B_{3/4})\). For all sufficiently large \(k\), \(B(x_0+R_kx,R_k)\subset B_{3/4}\), and change of variables followed by Hölder's inequality gives
\begin{align*}
 \int_{B_1(x)}|v^k(y,0)|^2\,dy
 &=R_k^{-1}\int_{B(x_0+R_kx,R_k)}|\bar u(y,t_0)|^2\,dy\\
 &\le |B_1|^{1/3}
       \left(\int_{B(x_0+R_kx,R_k)}|\bar u(y,t_0)|^3\,dy\right)^{2/3}\\
 &\longrightarrow0.
\end{align*}
The final limit follows from absolute continuity of the integral of the fixed \(L^3(B_{3/4})\) function on sets whose measure tends to zero. Weak convergence and lower semicontinuity of the \(L^2(B_1(x))\) norm prove the zero terminal-slice assertion in (e), for every center \(x\).

Finally, by Lemma~\ref{lem:good-open-glue} and equation~\eqref{eq:bad-point-lower-bound}, the bad-point lower bound applies at every admissible open-top radius. At radius \(R_k\) it gives, for every \(k\),
\begin{equation*}
 R_k^{-2}\iint_{Q_{R_k}^-(z_0)}
       (|u_{\mathrm{orig}}|^3+|p_{\mathrm{orig}}|^{3/2})\,dx\,dt\ge\varepsilon_0/8.
\end{equation*}
The Navier--Stokes scaling changes this exactly into the lower bound in (f). If \(u=0\) on \(\mathbb R^3\times(-T,0)\) with \(T>1\), then \(q=P[u\otimes u]=0\) there. Strong \(L^3(Q_1^-)\) convergence gives \(\int_{Q_1^-}|v^k|^3\to0\); the pressure convergence in (c) gives \(\int_{Q_1^-}|p^k|^{3/2}\to0\). This proves the final assertion.

\end{proof}

\begin{remark}[Departure from \cite{ESS2003}]
The rescaling and the roles of the local energy and pressure bounds are those of (3.13)--(3.20), printed pp.~223--225, of~\cite{ESS2003}. The compactness argument gives strong local space-time \(L^2\), strong \(L^3\), and weak convergence of every velocity slice; it does not use the strong \(C([a,b];L^2_{\mathrm{loc}})\) convergence asserted in (3.23), printed p.~225, of~\cite{ESS2003}. The endpoint conclusion replacing (3.34), printed p.~226, follows from the uniform pairing modulus, weak slice convergence at \(t=0\), and lower semicontinuity. The pressure split is the fixed whole-space Riesz split of Lemma~\ref{lem:pressure-split}; its harmonic remainder decays by that lemma's scaling estimate, proved for this split, rather than by transferring (3.19) to a newly defined pressure remainder.
\end{remark}
\section{Proof of the local theorem}\label{sec:thm-local}

\begin{lemma}[Flatness from a zero neighborhood]\label{lem:flat-from-vanishing}
Let $R',T',r_0,M>0$, $x_\ast\in\R^3$, and let
$w\colon B_{R'}(x_\ast)\times[0,T')\to\R^3$ be continuous. Suppose
$\abs{w}\le M$ on this cylinder and
\[
 w=0\quad\text{on }B_{r_0}(x_\ast)\times(0,r_0^2),
\]
where $r_0\le R'$ and $r_0^2\le T'$. Then, for every integer $k\ge1$,
\begin{equation}\label{eq:flat-from-vanishing}
 \abs{w(x,s)}\le M r_0^{-k}
       \bigl(\abs{x-x_\ast}+\sqrt{s}\bigr)^k
 \quad\text{for }(x,s)\in B_{R'}(x_\ast)\times(0,T').
\end{equation}
For $k=0$ the same conclusion holds with the factor on the right interpreted
as $M$.
\end{lemma}

\begin{proof}
Put $d=\abs{x-x_\ast}+\sqrt{s}$. If $d<r_0$, then
$x\in B_{r_0}(x_\ast)$ and $0<s<r_0^2$, so $w(x,s)=0$. If $d\ge r_0$,
then $M r_0^{-k}d^k\ge M\ge\abs{w(x,s)}$. This proves the assertion,
including the case $k=0$.
\end{proof}

\begin{lemma}[Exterior vorticity at the terminal face]\label{lem:vorticity-top-extension}
Let $T_2>2$ and $R_2>0$. Suppose $(u,Du,q,0)$ is suitable with exponent
$3$ on $\R^3\times(-4T_2,0)$, $u\in L^\infty((-4T_2,0);L^3(\R^3))$,
and its weakly continuous local trace satisfies $u(\cdot,0)=0$. Suppose
there are constants $M,E<\infty$ such that for every
$z_1=(x_1,t_1)$ with $\abs{x_1}>R_2$ and $t_1\in(-2,0)$,
$(u,Du,q,0)$ is suitable on an open neighborhood of
$\overline{Q_1^-(z_1)}$, $\abs{u}\le M$ on $Q_1^-(z_1)$, and
\[
 \left(\iint_{Q_1^-(z_1)}\abs{Du}^2\,dx\,dt\right)^{1/2}\le E.
\]
Then $\omega=\curl u$ has a representative continuous on
$\{\abs{x}>R_2\}\times(-2,0]$. On every bounded spatial set whose closure
lies in $\{\abs{x}>R_2\}$, this representative belongs to
$W^{2,1}_2$ up to $t=0$, satisfies
\begin{equation}\label{eq:exterior-vorticity-inequality}
 \abs{\dt\omega-\Delta\omega}
 \le C(M,E)\bigl(\abs{\omega}+\abs{\nabla\omega}\bigr)
 \quad\text{a.e. on }\{\abs{x}>R_2\}\times(-2,0),
\end{equation}
and has $\omega(\cdot,0)=0$ there. In addition, its representative satisfies
\[
 |\omega(x,t)|\le C_\omega(M,E)
 \quad (|x|>R_2,\ -2<t\le0),
\]
with a constant independent of the exterior point. The coefficient in the
inequality depends only on $M,E$; the local $W^{2,1}_2$ norm may also depend
on the bounded set.
\end{lemma}

\begin{proof}
Apply Theorem~\ref{thm:vorticity-regularity} to each indicated unit past
cylinder. Its bounds depend only on the common $M$ and $E$. In particular,
on $B_{1/2}(x_1)\times(t_1-1/4,t_1)$ they give a common bound for
$\omega$ in $W^{2,1}_2$, for $\omega$ in $L^\infty_tH^2_x$, and for the
continuous representative on the closed half-cylinder. Representatives
from two cylinders agree on overlaps: they are continuous and represent the
same weak curl almost everywhere, so if they differed at a point of an open
overlap they would differ on a set of positive measure. They therefore glue
to a continuous representative on the exterior open region below $t=0$.
The common continuous bound on each half-cylinder also gives a bound
independent of the exterior point. Indeed, for $|x|>R_2$ and $-2<t<0$,
choose $t_1$ with $t<t_1<\min\{t+1/4,0\}$. Then
$(x,t)\in B_{1/2}(x)\times(t_1-1/4,t_1)$, so the half-cylinder estimate
gives $|\omega(x,t)|\le C_0(M,E)$. Continuity up to $t=0$ gives the same
bound there.

Fix a ball $B$ whose closure is contained in $\{\abs{x}>R_2\}$. Cover
$\overline B$ by finitely many balls $B_{1/2}(x_j)$ with
$\abs{x_j}>R_2$. For the initial time portion, the interval $(-2,-7/4)$
is contained in the half-cylinder time interval centered at $-7/4$.
The compact interval $[-7/4,-1/4]$ is covered by finitely many intervals
$(s_\ell-1/4,s_\ell)$ with centers $s_\ell\in(-2,0)$, since this family
is an open cover of that compact interval. For the terminal portion set
$\tau_n=-1/(n+4)$ for $n\ge1$ and $I_n=(\tau_n-1/4,\tau_n)$. Then
$\tau_n\uparrow0$ and, for every fixed $t\in(-1/4,0)$,
$\mathbf1_{I_n}(t)\to1$. For each spatial center $x_j$ and each of the
finitely many initial time centers, and for each $(x_j,\tau_n)$, the
half-cylinder estimates give with one constant $C=C(M,E)$
\[
 \esssup_{t\in I}\norm{\omega(\cdot,t)}_{H^2(B\cap B_{1/2}(x_j))}
 +\norm{\dt\omega}_{L^2((B\cap B_{1/2}(x_j))\times I)}
 \le C,
\]
where $I$ denotes the corresponding time interval. They also bound each
spatial derivative through order two in $L^2$ on that patch and interval.
The finite covers give the asserted bounds on the initial portion. On the
terminal portion, Fatou's lemma applied to the indicators $\mathbf1_{I_n}$
and to the squares of $\dt\omega$ and of each spatial derivative through
order two gives finite $L^2$ norms on $B\times(-1/4,0)$. Since the $I_n$
cover this terminal portion and the $H^2$ bounds are uniform, countability
also gives $\esssup_{-1/4<t<0}\norm{\omega(t)}_{H^2(B)}\le C_B$.
Combining the two portions shows that on $B\times(-2,0)$ these derivatives
have finite $L^2$ norms, that
$\esssup_{-2<t<0}\norm{\omega(t)}_{H^2(B)}\le C_B$, and that
$\dt\omega\in L^2(B\times(-2,0))$. Thus
$\omega\in W^{1,2}((-2,0);L^2(B))$. The fundamental theorem for
Bochner--Sobolev functions gives
\[
 \norm{\omega(t)-\omega(s)}_{L^2(B)}
 \le C_B\abs{t-s}^{1/2}\qquad(-2<s<t<0).
\]
Interpolate between $L^2(B)$ and $H^2(B)$ at order $7/4$ and use
$H^{7/4}(B)\hookrightarrow L^\infty(B)$. The uniform $H^2$ bound then
gives
\[
 \norm{\omega(t)-\omega(s)}_{L^\infty(B)}
 \le C_B\abs{t-s}^{1/16}.
\]
Consequently $\omega$ has a unique continuous extension to $t=0$ on $B$.
The spatial embedding $H^2(B)\hookrightarrow C^{0,\alpha}(\overline B)$
for every $\alpha<1/2$ shows that this extension is jointly continuous in
space and time. The local $W^{2,1}_2$ bounds and the inequality
\eqref{eq:exterior-vorticity-inequality} pass to the union by countable
covering; the endpoint is a null time slice for the inequality.

It remains to identify the trace. For every smooth compactly supported
vector test $\varphi$ in $B$ and almost every $t<0$, the definition of weak
curl gives
\[
 \int_B\omega(x,t)\cdot\varphi(x)\,dx
 =\int_B u(x,t)\cdot\curl\varphi(x)\,dx.
\]
The left side converges to the pairing with the continuous extension at
$t=0$. By the weakly continuous $L^3_{\mathrm{loc}}$ trace and
$u(\cdot,0)=0$, the right side tends to zero. Hence the extension has zero
curl-trace as a distribution on $B$. Since it is continuous in space, it
vanishes at every point of $B$ at $t=0$. The ball was arbitrary.
\end{proof}

\begin{lemma}[Small pressure mass near a regular point]\label{lem:pressure-smallness-at-regular-point}
Let $(u,Du,q,0)$ be suitable with exponent $3$ on an open set containing
$\overline{Q_{r_0}^-(z_0)}$, and suppose $q=P[u\otimes u]$ there. If $u$ has
a parabolically H\"older representative on a neighborhood of $z_0$, then
\begin{equation}\label{eq:regular-point-pressure-smallness}
 \lim_{r\downarrow0}r^{-2}\iint_{Q_r^-(z_0)}
       (\abs{u}^3+\abs{q}^{3/2})\,dx\,dt=0.
\end{equation}
\end{lemma}

\begin{proof}
Shrink $r_0$ so that the closure of the resulting cylinder is inside the
H\"older neighborhood. On this compact cylinder, $\abs{u}$ and the spatial
H\"older seminorm of $u$ are bounded by a constant $H$. The velocity term
in \eqref{eq:regular-point-pressure-smallness} is at most
$C H^3r^5$, so after division by $r^2$ it tends to zero.

Choose $\chi\in C_c^\infty(B_{r_0}(x_0))$ equal to one on
$B_{r_0/2}(x_0)$ and put
$F=\chi u\otimes u$, extended by zero in space. For almost every time,
$F(\cdot,t)$ is uniformly compactly supported and uniformly
$C^\alpha$ for some $\alpha>0$. The kernel formula for the double Riesz
transforms defining $P$, its spherical cancellation, and
$\abs{K(x-y)}\le C\abs{x-y}^{-3}$ give
\[
 \norm{P[F(\cdot,t)]}_{L^\infty(B_{r_0/2}(x_0))}\le C(H,r_0,\chi).
\]
Indeed, in the principal-value integral subtract $F(x,t)$ near $x$; the
remaining near-field integrand is bounded by
$C H\abs{x-y}^{\alpha-3}$, which is integrable in dimension three, and the
far-field integral is bounded because $F$ has fixed compact support.

Set $h=q-P[F]$ on $B_{r_0/2}(x_0)$. The pressure identity from (S3) and the
definition of $P$ imply $\Delta_x h=0$ there for almost every time. Also
$h\in L^{3/2}$ on a slightly larger compact cylinder. The harmonic interior
estimate therefore gives a function $g\in L^1$ of time such that
\[
 \norm{h(\cdot,t)}_{L^\infty(B_{r_0/4}(x_0))}^{3/2}\le g(t)
 \quad\text{for a.e. }t.
\]
For $r<r_0/4$, the pressure mass of $P[F]$ on $Q_r^-(z_0)$ is at most
$C r^5$. The pressure mass of $h$ is at most
\[
 C r^3\int_{t_0-r^2}^{t_0}g(t)\,dt=o(r^3),
\]
by absolute continuity of the integral of $g$. After division by $r^2$ both
terms tend to zero. The inequality
$\abs{a+b}^{3/2}\le C(\abs a^{3/2}+\abs b^{3/2})$ completes the proof.
\end{proof}

\begin{lemma}[Liouville theorem for a curl-free $L^3$ field]\label{lem:liouville-L3}
If $v\in L^3(\R^3;\R^3)$ and $\dv v=0$, $\curl v=0$ in distributions,
then $v=0$ almost everywhere.
\end{lemma}

\begin{proof}
The vector identity
$\Delta v=\nabla(\dv v)-\curl(\curl v)$, valid in distributions,
shows that every component of $v$ is distributionally harmonic. For
completeness, mollify one component on each relatively compact ball. The
mollifications are smooth harmonic functions on smaller balls. A smooth
harmonic function equals its convolution with any radial smooth unit-mass
kernel supported in such a ball, by the mean-value property. Differentiating
this identity bounds every derivative on a compact subball by a constant
times the $L^1$ norm on a slightly larger ball. The mollifications have
uniform such bounds because they converge to $v_i$ in $L^1_{\mathrm{loc}}$.
Arzel\`a--Ascoli, applied successively to derivatives and then diagonally on
an exhaustion, gives a smooth harmonic representative equal to $v_i$
almost everywhere. This proves the needed form of Weyl's lemma.

For this representative, the mean-value identity and H\"older's inequality
give, for every $x\in\R^3$ and $R>0$,
\[
 \abs{v_i(x)}\le \frac1{\abs{B_R}\,}\int_{B_R(x)}\abs{v_i(y)}\,dy
 \le \abs{B_R}^{-1/3}\norm{v_i}_{L^3(\R^3)}.
\]
Since $\abs{B_R}^{-1/3}=\abs{B_1}^{-1/3}R^{-1}$, letting $R\to\infty$
shows $v_i(x)=0$. This holds for each component.
\end{proof}

\begin{proof}[Proof of Theorem~\ref{thm:ess-local}]
We argue by contradiction. By Lemma~\ref{lem:lei-L4}, the hypotheses of
Theorem~\ref{thm:ess-local} make $(u,Du,p,0)$ suitable with exponent $3$
on $B_1\times(-1,0)$: every compactly supported test in this open cylinder
is supported in some $B_r\times(-1,0)$, $r<1$, where that lemma gives the
local energy inequality. Let
\[
 K=\overline{B_{1/2}}\times[-1/4,0].
\]
If $u$ has no H\"older representative on $K$, the contrapositive of
Lemma~\ref{lem:good-open-glue} (with its amended top-face definition of
good points) gives $z_0=(x_0,t_0)\in K$ that is not good. Choose any
sequence $R_k\downarrow0$ small enough that the open-top cylinders
$Q^-_{R_k}(z_0)$ lie in $B_1\times(-1,0)$. Proposition~\ref{prop:blowup-limit}, using
\eqref{eq:bad-point-lower-bound}, supplies a subsequential blow-up limit
$(U,DU,q,0)$ on $\R^3\times(-\infty,0)$. Clause (c) of that
proposition gives $U\in L^\infty((-\infty,0);L^3(\R^3))$; clause (d) gives
suitability with exponent $3$ on every bounded cylinder in the open past and
$q=P[U\otimes U]$; clause (e) gives a terminal slice $U(\cdot,0)=0$ in
$L^2_{\mathrm{loc}}$; and clause (f) gives the lower bound
\[
 \iint_{Q_1^-(0)}\bigl(\abs{v^k}^3+\abs{p^k}^{3/2}\bigr)\,dx\,dt
       \ge\varepsilon_0/8\quad\text{for every }k,
\]
with convergence of these integrals to zero if $U=0$ on
$\R^3\times(-T,0)$ for any $T>1$. Here $v^k,p^k$ are the rescalings
in Proposition~\ref{prop:blowup-limit}; its top-face cylinders are open
at $t=0$, so no future values are used.

\paragraph{Exterior regularity.}
Fix $T_2>2$. The bounds $U\in L^\infty L^3$ and
$q=P[U\otimes U]$, together with the $L^{3/2}$ boundedness of the double
Riesz transforms, imply
$U\in L^3(\R^3\times(-4T_2,0))$ and
$q\in L^{3/2}(\R^3\times(-4T_2,0))$. Hence
\[
 \iint_{(\R^3\setminus B_A)\times(-4T_2,0)}
       (\abs U^3+\abs q^{3/2})\,dx\,dt\longrightarrow0
       \quad(A\to\infty).
\]
Let $\varepsilon_{\rm A}$ be the threshold in
Lemma~\ref{lem:thmA-rescaled}, and set $r_{\rm e}=1/4$. Choose $A$ so
large that this tail is less than $\varepsilon_{\rm A}r_{\rm e}^2$.
For every center $z=(x,t)$ with $\abs x>A+1$ and
$-2T_2-1<t<0$, choose a top time $t'>t$ with
$t'<0$ and $t'-t<r_{\rm e}^2/8$. The closed full cylinder
$\overline{Q_{r_{\rm e}}^-(x,t')}$ is contained in the open solution domain
and in the region used for the tail estimate: its spatial ball lies outside
$B_A$, and its time interval lies in $(-4T_2,0)$ because $T_2>2$.
Its normalized integral is less than $\varepsilon_{\rm A}$. Theorem
\ref{lem:thmA-rescaled} gives a representative bounded by
$C_4/r_{\rm e}$ on its half-cylinder, which contains $(x,t)$. This proves
an essential bound $\abs U\le M$ on
$\{\abs x>A+1\}\times(-2T_2-1,0)$, with $M=C_4/r_{\rm e}$.

Set $R_2=A+3$ and $R_3=R_2+2$. For every
$z_1=(x_1,t_1)$ with $\abs{x_1}>R_3$ and $t_1\in(-2T_2,0)$, the
cylinder $Q_1^-(z_1)$ lies in the bounded-velocity region just obtained and
in $\R^3\times(-4T_2,0)$. Its local gradient energy has a bound independent
of $x_1,t_1$. To see this, apply the local energy inequality with a fixed
spatial cutoff equal to one on $B_1(x_1)$ and supported in $B_2(x_1)$, and
with a time cutoff that rises from zero to one on $(t_1-2,t_1-1)$,
equals one on $[t_1-1,t_1]$, and falls to zero on
$(t_1,t_1+\delta_1)$, where
$\delta_1=\min\{-t_1/2,1/8\}$. This interval lies before $0$, and the
total variation of the cutoff is $2$, uniformly in $t_1$. Its support lies
in $(-4T_2,0)$ because $t_1>-2T_2$ and $T_2>2$. The terms with $U^2$ are
bounded by the global $L^\infty_tL^3_x$ bound on balls of fixed volume; the
terms with $U^3$ and $qU$ are bounded by the global $L^\infty_tL^3_x$ and
$L^\infty_tL^{3/2}_x$ bounds and H\"older's inequality. The local energy
inequality with this test bounds the gradient integral on the whole open-top
cylinder by a constant uniform in $x_1,t_1$. The integrals over
$B_1(x_1)\times(t_1-1,s)$ increase to this full integral as $s\uparrow t_1$,
so monotone convergence includes the top endpoint. Thus
\[
 \left(\iint_{Q_1^-(z_1)}\abs{DU}^2\right)^{1/2}\le E
\]
for a constant $E$ independent of the exterior center. For each such
$z_1$, the closure of $Q_1^-(z_1)$ is contained in an open cylinder compactly
contained in $\R^3\times(-4T_2,0)$; suitability restricts to that
neighborhood. All hypotheses of Theorem~\ref{thm:vorticity-regularity}
therefore hold with the same $M,E$.

We first verify the trace hypothesis of Lemma~\ref{lem:vorticity-top-extension}.
For every smooth compactly supported test field $\psi$, clause (b) of
Proposition~\ref{prop:blowup-limit} gives the rescaled fields a uniform
pairing modulus on intervals $[a,0]$. By clause (c), their pairings
converge at every negative time to those of $U$. By clause (e), their
pairings at time zero converge to those of the terminal slice $U(\cdot,0)$,
which is zero in $L^2_{\rm loc}$. Passing to the limit in the modulus
therefore gives
\[
 \left|\int U(x,t)\cdot\psi(x)\,dx
       -\int U(x,0)\cdot\psi(x)\,dx\right|
 \le C_\psi\bigl(|t|+|t|^{1/3}\bigr)\longrightarrow0
 \quad(t\uparrow0).
\]
Thus $U$ has the required weakly continuous local trace and that trace is
zero. The uniform $L^\infty_tL^3_x$ bound in clause (c) and density of
smooth compactly supported tests in $L^{3/2}$ extend this pairing
continuity to weak continuity in local $L^3$. Lemma~\ref{lem:vorticity-top-extension}
now gives a continuous
representative of $\omega=\curl U$ on
$\{\abs{x}>R_3\}\times(-2,0]$, locally in $W^{2,1}_2$ up to the top, with
a uniform inequality
\[
 \abs{\dt\omega-\Delta\omega}
 \le c_{\rm ext}(\abs\omega+\abs{\nabla\omega})
\]
and terminal value $\omega(\cdot,0)=0$.

\paragraph{Backward uniqueness outside a ball.}
Fix a unit vector $e$ and choose $O\in SO(3)$ with $Oe_3=e$. For
$(y,s)\in\R^3_+\times[0,1)$ set
\[
 x=R_3e+\sqrt{2}\,Oy,\qquad t=-2s,\qquad
 w(y,s)=A_\omega^{-1}\omega(R_3e+\sqrt{2}\,Oy,-2s),
\]
where $A_\omega=\max\{1,\|\omega\|_{L^\infty(\{|x|>R_3\}\times(-2,0))}\}$.
Then $x\cdot e=R_3+\sqrt2 y_3>R_3$ and the map sends
$\{x\cdot e>R_3\}\times(-2,0]$ onto
$\R^3_+\times[0,1)$. The weak chain rule gives
\[
 (\partial_s+\Delta_y)w
 =-2A_\omega^{-1}(\dt\omega-\Delta_x\omega)
       (R_3e+\sqrt2 Oy,-2s),
 \qquad
 \abs{\nabla_x\omega}/A_\omega=2^{-1/2}\abs{\nabla_y w}.
\]
Thus
$\abs{\partial_s w+\Delta_yw}\le c_1(\abs w+\abs{\nabla_yw})$ with
$c_1=2c_{\rm ext}$, after increasing it if necessary. The zero terminal
trace gives $w(y,0)=0$; the exterior bound gives $\abs w\le1\le e^0$;
and the local $W^{2,1}_2$ estimates give all required $L^2$ bounds on
bounded subsets. Theorem~\ref{thm:bu} yields $w=0$ on $\R^3_+\times(0,1)$.
As $e$ was arbitrary, choosing $e=x/\abs x$ shows
\begin{equation}\label{eq:omega-exterior-zero}
 \omega=0\quad\text{on }(\R^3\setminus\overline{B_{R_3}})\times(-2,0).
\end{equation}

\paragraph{Regular time slices and unique continuation.}
By Proposition~\ref{prop:blowup-limit}(d), the limit is suitable on every
bounded subcylinder of $\R^3\times(-2,0)$. Every compactly supported test
in the full open domain $\R^3\times(-2,0)$ is supported in one such bounded
cylinder, so these local suitability statements give the global CKN
suitability predicate. Clause (c) gives
$U\in L^\infty((-2,0);L^3(\R^3))$, and clause (d) gives the global
pressure identity $q=P[U\otimes U]$. Boundedness of the double Riesz
transforms on $L^{3/2}(\R^3)$ therefore gives
$q\in L^\infty((-2,0);L^{3/2}(\R^3))$, the required global pressure
integrability; the local gradient and energy conditions are those in (d).
Thus the limit lies in the exponent-three CKN class on the full open
domain. CKN Theorem~C therefore says that the singular set $S$ there has
zero one-dimensional parabolic Hausdorff measure. By
Lemma~\ref{lem:time-projection}, for almost every $t_\ast\in(-2,0)$ there
is an interval $I$ about $t_\ast$ such that
$\overline{B_{4R_3}}\times I$ consists of regular points. Shrink $I$ so
that its closure is compactly contained in $(-2,0)$.

We record the small-scale input at each regular point. If $z$ is regular,
choose $r_0>0$ so that $\overline{Q_{r_0}^-(z)}$ lies in its regularity
neighborhood and in the open solution domain. By Proposition~\ref{prop:blowup-limit}(d), the limit is suitable with exponent $3$ there;
by clauses (c)--(d), its pressure is the global associated pressure
$q=P[U\otimes U]$. Lemma~\ref{lem:pressure-smallness-at-regular-point}
applies directly and gives
\[
 r^{-2}\iint_{Q_r^-(z)}(\abs U^3+\abs q^{3/2})\longrightarrow0.
\]
For small $r$, Lemma~\ref{lem:thmA-rescaled} also gives a Hölder
representative and an explicit bound on a smaller cylinder. Since the
closed cylinder is inside the open domain, Theorem
\ref{thm:vorticity-regularity} applies there after parabolic rescaling.
Cover the compact cylinder $\overline{B_{4R_3}}\times\overline I$ by
finitely many of these smaller regular neighborhoods. The resulting
vorticity representatives glue; on every smaller cylinder they satisfy
$W^{2,1}_2$ bounds, are continuous, and obey
\begin{equation}\label{eq:omega-local-uc-inequality}
 \abs{\dt\omega-\Delta\omega}
       \le c_I(\abs\omega+\abs{\nabla\omega}).
\end{equation}
The constants and the bound $\abs\omega\le M_I$ are uniform on this
compact cylinder.

Define
\[
 V=\{x\in B_{4R_3}:\text{ for some }r>0,
       \ \omega=0\text{ on }B_r(x)\times I\}.
\]
This set is relatively open. It is nonempty by
\eqref{eq:omega-exterior-zero}, since the exterior zero region intersects
$B_{4R_3}$. We show it is relatively closed. Let
$x\in\overline V\cap B_{4R_3}$. Choose $R'>0$ with
$\overline B_{2R'}(x)\subset B_{4R_3}$, and choose $x_\ast\in V$ so close
to $x$ that $\abs{x-x_\ast}<R'/2$. By the definition of $V$, there is
$r>0$ such that $\omega=0$ on $B_r(x_\ast)\times I$.
Fix any $t_\ast\in I$. Choose $T'>0$ with
$T'<t_\ast-\inf I$, and choose
$r_0>0$ so small that $r_0<\min\{r,R'\}$ and $r_0^2<T'$. On
$B_{R'}(x_\ast)\times(t_\ast-T',t_\ast]$, the vorticity is continuous,
belongs to $W^{2,1}_2$ below the top, is bounded by $M_I$, and satisfies
\eqref{eq:omega-local-uc-inequality}. Reverse time by setting
$w(y,s)=\omega(x_\ast+y,t_\ast-s)$ for $y\in B_{R'}$ and
$0\le s<T'$. Then
\[
 \abs{\partial_s w+\Delta_yw}
 \le c_I(\abs w+\abs{\nabla_yw}),\qquad
 w=0\text{ on }B_{r_0}(0)\times(0,r_0^2).
\]
Lemma~\ref{lem:flat-from-vanishing} gives the pointwise infinite-order
bound required by Theorem~\ref{thm:uc}, with positive constants
$C_k=\max\{1,M_Ir_0^{-k}\}$; its $W^{2,1}_2$ hypotheses follow
from the local vorticity bounds. Theorem~\ref{thm:uc} implies
$\omega(\cdot,t_\ast)=0$ on $B_{R'}(x_\ast)$, which contains a neighborhood
of $x$. This argument applies to every $t_\ast\in I$, with the same
$R'$ and the same spatial neighborhood of $x$. Hence $x\in V$.
Connectedness of $B_{4R_3}$ gives $V=B_{4R_3}$, and therefore
$\omega(\cdot,t_\ast)=0$ throughout $B_{4R_3}$ for every
$t_\ast\in I$. Together with \eqref{eq:omega-exterior-zero}, this proves
$\omega(\cdot,t_\ast)=0$ on $\R^3$ for every time $t_\ast$ in the
full-measure set supplied by Theorem~C and Lemma~\ref{lem:time-projection}.

For almost every such $t$, $U(\cdot,t)\in L^3(\R^3)$,
$\dv U(\cdot,t)=0$, and $\curl U(\cdot,t)=\omega(\cdot,t)=0$ in spatial
distributions. Lemma~\ref{lem:liouville-L3} gives
$U(\cdot,t)=0$ for almost every $t\in(-2,0)$, hence
$U=0$ almost everywhere on $\R^3\times(-2,0)$.

This contradicts Proposition~\ref{prop:blowup-limit}(f): the bad-point
lower bound \eqref{eq:bad-point-lower-bound} gives a lower bound
$\varepsilon_0/8$ for every rescaled integral on $Q_1^-(0)$, whereas the
last clause of (f), with $T=2>1$, says those same integrals tend to zero.
The contradiction proves that $u$ has a parabolically H\"older
representative on $K$, which is the conclusion of
Theorem~\ref{thm:ess-local}.
\end{proof}

\begin{remark}[Departure from \cite{ESS2003}]\label{rem:thm-local-departure}
The proof follows the blow-up, exterior backward-uniqueness, and interior
unique-continuation structure of Theorem~1.4 in~\cite[pp.~222--229]{ESS2003}. The endpoint good-point criterion is stated on open-top
past cylinders, and the pressure split is the whole-space Riesz split of
Lemma~\ref{lem:pressure-split}; the endpoint lower bound and pressure
convergence are supplied by Proposition~\ref{prop:blowup-limit}. We give
the terminal vorticity extension and the infinite-order vanishing step
explicitly, and use CKN Theorem~C for the regular-time projection. These
choices make the top-face trace and both applications of unique continuation
explicit in the regularity classes used here.
\end{remark}

\section{Proof of the global theorem}\label{sec:thm-global}

The pressure needed for the local theorem is the canonical whole-space pressure
of the velocity tensor. Its mixed-norm bound follows directly from the
slice-wise Riesz estimate.

\begin{lemma}[Pressure bound under the critical velocity hypothesis]
\label{lem:assoc-pressure-L3}
Suppose $(u,Du)$ satisfies the hypotheses of
Theorem~\ref{thm:ess-global}. Let $p_*=P[u\otimes u]$ be the pressure witness
selected in the proof of Theorem~\ref{thm:assoc-pressure}. Then
$p_*\in L^{5/3}(\R^3\times(0,T))$, the identity (S3) of
Definition~\ref{def:sws} holds with $p=p_*$ and $f=0$, and
$p_*\in L^\infty((0,T);L^{3/2}(\R^3))$. There is a constant $C$ depending
only on the dimension and the $L^{3/2}$ double-Riesz-transform bound such that
\[
 \esssup_{t\in(0,T)}\norm{p_*(\cdot,t)}_{L^{3/2}(\R^3)}
 \le C\left(\esssup_{t\in(0,T)}
                 \norm{u(\cdot,t)}_{L^3(\R^3)}\right)^2.
\]
\end{lemma}

\begin{proof}
The proof of Theorem~\ref{thm:assoc-pressure} takes
$F_{ij}=u_i u_j$ and sets $p_*=P[F]$. \cite[Lemma~8.4]{CKNLean}
gives $F\in L^{5/3}(\R^3\times(0,T))$; the proof there uses
\cite[Lemma~8.2]{CKNLean} to verify the gradient part of (S3). Thus the
existential theorem's chosen witness is exactly $P[u\otimes u]$, with the
stated global integrability and (S3) identity.

For almost every $t$, the hypothesis of Theorem~\ref{thm:ess-global} gives
$u(\cdot,t)\in L^3(\R^3)$. The slice estimate in
Definition~\ref{def:riesz-pressure}, applied with $r=3/2$, yields
\[
 \norm{p_*(\cdot,t)}_{L^{3/2}}
 \le C_{3/2}\sum_{i,j=1}^3
                    \norm{u_i(\cdot,t)u_j(\cdot,t)}_{L^{3/2}}
 \le 9C_{3/2}\norm{u(\cdot,t)}_{L^3}^2.
\]
Taking the essential supremum proves the estimate. The pressure representative
in Definition~\ref{def:riesz-pressure} is jointly measurable, so it also has
the measurability required below.
\end{proof}

\begin{proof}[Proof of Theorem~\ref{thm:ess-global}]
Fix $z_0=(x_0,t_0)\in\R^3\times(0,T)$. Choose $R>0$ so small that
\[
 R^2<\min\left\{\frac{8t_0}{7},\,8(T-t_0)\right\},
 \qquad t_1=t_0+\frac{R^2}{8}.
\]
Then $0<t_0-7R^2/8<t_1<T$. Set
$I_R=(t_0-7R^2/8,t_1)$ and define, on $Q=B_1\times(-1,0)$,
\begin{align*}
 \widetilde u(y,s)&=R u(x_0+Ry,t_1+R^2s),\\
 \widetilde{Du}_{ij}(y,s)&=R^2Du_{ij}(x_0+Ry,t_1+R^2s),\\
 \widetilde p(y,s)&=R^2p_*(x_0+Ry,t_1+R^2s).
\end{align*}
The affine map $(y,s)\mapsto(x_0+Ry,t_1+R^2s)$ sends $Q$ onto
$B_R(x_0)\times I_R$, whose closure is contained in
$\R^3\times(0,T)$. Composition with this map preserves measurability.

The change of variables gives
\begin{align*}
 \esssup_{s\in(-1,0)}\int_{B_1}|\widetilde u(y,s)|^2\,dy
 &\le R^{-1}\esssup_{t\in(0,T)}\int_{\R^3}|u(x,t)|^2\,dx<\infty,\\
 \iint_Q|\widetilde u|^2\,dy\,ds
 &=R^{-3}\iint_{B_R(x_0)\times I_R}|u|^2\,dx\,dt<\infty,\\
 \iint_Q|\widetilde{Du}|^2\,dy\,ds
 &=R^{-1}\iint_{B_R(x_0)\times I_R}|Du|^2\,dx\,dt<\infty,\\
 \iint_Q|\widetilde p|^{3/2}\,dy\,ds
 &=R^{-2}\iint_{B_R(x_0)\times I_R}|p_*|^{3/2}\,dx\,dt<\infty.
\end{align*}
The first three right-hand sides are finite by (LH1) of
Definition~\ref{def:leray-hopf} and the last is finite by
Lemma~\ref{lem:assoc-pressure-L3}. Also,
\[
 \esssup_{s\in(-1,0)}\int_{B_1}|\widetilde u(y,s)|^3\,dy
 \le \esssup_{t\in(0,T)}\int_{\R^3}|u(x,t)|^3\,dx<\infty.
\]

For almost every $s\in(-1,0)$, the physical time $t=t_1+R^2s$ is a slice on
which $Du(\cdot,t)$ is the weak spatial gradient of $u(\cdot,t)$. If
$\psi\in\Cc(B_1)$ and $\psi_R(x)=\psi((x-x_0)/R)$, that weak-gradient identity
and the change of variables give, componentwise,
\[
 \int_{B_1}\widetilde u_i(y,s)\partial_j\psi(y)\,dy
 =-R^{-1}\int_{B_R(x_0)}Du_{ij}(x,t)\psi_R(x)\,dx
 =-\int_{B_1}\widetilde{Du}_{ij}(y,s)\psi(y)\,dy.
\]
Thus $\widetilde{Du}$ is the weak gradient on almost every time slice.

It remains to check (S2) and (S3), including the test scaling. For
$\eta\in\Cc(Q)$ scalar and $\varphi\in\Cc(Q;\R^3)$, put
\[
 \eta_R(x,t)=\eta\left(\frac{x-x_0}{R},\frac{t-t_1}{R^2}\right),\qquad
 \varphi_R(x,t)=\varphi\left(\frac{x-x_0}{R},\frac{t-t_1}{R^2}\right).
\]
These are smooth compactly supported tests in
$\R^3\times(0,T)$. Since $dy\,ds=R^{-5}dx\,dt$ and
$\nabla_y\eta=R\nabla_x\eta_R$ under this substitution, (LH1) gives
\[
 \iint_Q\widetilde u\cdot\nabla_y\eta\,dy\,ds
 =R^{-3}\iint_{B_R(x_0)\times I_R}u\cdot\nabla_x\eta_R\,dx\,dt=0,
\]
which is (S2). For the vector test,
$\partial_t\varphi_R=R^{-2}\partial_s\varphi$ and
$\partial_{x_j}(\varphi_R)_i=R^{-1}\partial_{y_j}\varphi_i$; all four terms
of (S3) consequently have the same scaling. In full,
\begin{align*}
 &\iint_Q\left[-\widetilde u\cdot\partial_s\varphi
 -\widetilde u_i\widetilde u_j\partial_{y_j}\varphi_i
 +\widetilde{Du}_{ij}\partial_{y_j}\varphi_i
 -\widetilde p\,\dv_y\varphi\right]dy\,ds\\
 &\qquad=R^{-2}\iint_{B_R(x_0)\times I_R}
 \left[-u\cdot\partial_t\varphi_R-u_i u_j\partial_{x_j}(\varphi_R)_i
 +Du_{ij}\partial_{x_j}(\varphi_R)_i-p_*\dv_x\varphi_R\right]dx\,dt=0,
\end{align*}
by (S3) for $p_*$ from Lemma~\ref{lem:assoc-pressure-L3}. This verifies every
hypothesis of Theorem~\ref{thm:ess-local}; its statement requires (S2) and
(S3), and does not require the local energy inequality (S4).

The local theorem supplies a parabolically H\"older representative of
$\widetilde u$ on
$\overline{B_{1/2}\times(-1/4,0)}$. The point $z_0$ has scaled coordinates
$(0,-1/8)$. In physical variables the closed half-cylinder is
\[
 \overline{B_{R/2}(x_0)}\times[t_0-R^2/8,t_1].
\]
It is contained in $\R^3\times(0,T)$, and the representative equals $u$
almost everywhere on its interior. Thus the hypotheses of
Lemma~\ref{lem:regular-point-shift} hold. The open cylinder
$B_{1/4}\times(-3/16,-1/16)$ contains that point and has closure inside the
half-cylinder above. Pulling its representative back by the inverse scaling
therefore gives a H\"older representative of $u$ on the open cylinder $N$
below, since the parabolic metric scales by the factor $R$ under this map:
\[
 N=B_{R/4}(x_0)\times(t_0-R^2/16,t_0+R^2/16)
 \subset\R^3\times(0,T).
\]
The conclusion of Lemma~\ref{lem:regular-point-shift} is the
open-neighborhood condition in the definition of a regular point.
Since $z_0$ was arbitrary, every point of $\R^3\times(0,T)$ is regular.
\end{proof}

\begin{remark}[Departure from \cite{ESS2003}]
The pressure estimate (3.4) of~\cite{ESS2003}, printed p.~221, is obtained
there from the pressure equation after invoking coercive Stokes estimates.
Here the proof of Theorem~\ref{thm:assoc-pressure} selects the canonical
pressure $P[u\otimes u]$, and the slice-wise double-Riesz bound gives (3.4)
directly from the $L^\infty_tL^3_x$ hypothesis. The weak continuity asserted
in (3.1) is not needed for this regularity conclusion. To obtain the
open-neighborhood condition in the definition of a regular point, the local theorem is
applied at the later center $(x_0,t_0+R^2/8)$; this puts $z_0$ strictly inside
its Hölder half-cylinder, as required by the shifted-cylinder argument of
Lemma~\ref{lem:regular-point-shift}. We use the coherent backward cylinder
$B_1\times(-1,0)$ in Theorem~\ref{thm:ess-local}: the printed statement of
ESS Theorem~1.4 on printed p.~214 says $B_1\times(0,1)$, although its displayed
hypotheses and conclusion use negative times. This proof establishes the
regularity conclusion stated as Theorem~\ref{thm:ess-global}.
\end{remark}

% Part V: L^5 bounds and uniqueness
\part*{Part V. $L^5$ bounds and uniqueness}

A finite $L^\infty_tL^3_x$ bound controls the solution at positive times, but it does not itself give space-time integrability near the initial time. We first recover an $L^3$ trace and construct a solution from that trace whose $L^5$ norm is small on a short interval. A comparison argument transfers that estimate to the given weak solution. Away from time zero, the local regularity theorem and the pressure tail give a global velocity bound, which turns the finite energy into the remaining $L^5$ estimate.

\section{The trace and heat flow}\label{sec:pv-trace}

Write $Q_\tau=\mathbb R^3\times(0,\tau)$ and let $S(t)=e^{t\Delta}$ be the heat semigroup with the Fourier convention fixed in \cite[Section~9]{CKNLean}.

\begin{lemma}[The initial trace is in $L^3$]\label{lem:pv-initial-l3}
Let $(u,Du)$ be a Leray--Hopf solution on $\mathbb R^3\times[0,T]$ with datum $a\in\J$, and suppose
\begin{equation*}
M^3:=\operatorname*{ess\,sup}_{t\in(0,T)}\int_{\mathbb R^3}|u(x,t)|^3\,dx<\infty\,.
\end{equation*}
Then $a\in L^3(\mathbb R^3;\mathbb R^3)$ and $\|a\|_{L^3}\le M$.
\end{lemma}

\begin{proof}
The set of times at which $\|u(t)\|_3\le M$ has full measure. Choose $t_n\downarrow0$ in this set. The strong $L^2$ initial trace in Definition~\ref{def:leray-hopf} gives $u(t_n)\to a$ in $L^2$, hence in measure on $\mathbb R^3$. Passing to a subsequence, $u(t_{n_k})\to a$ almost everywhere. Fatou's lemma gives $\int_{\mathbb R^3}|a|^3\,dx\le\liminf_{k\to\infty}\int_{\mathbb R^3}|u(x,t_{n_k})|^3\,dx\le M^3$, so $a\in L^3$ and $\|a\|_{L^3}\le M$. Its membership in $\J$ is already part of the Leray--Hopf hypothesis.
\end{proof}

\begin{lemma}[Critical heat estimate]\label{lem:pv-heat-critical}
For $b\in L^3(\mathbb R^3;\mathbb R^3)$ and $\tau>0$,
\begin{equation}\label{eq:pv-heat-l5}
 \|S(\cdot)b\|_{L^5(Q_\tau)}\le C\|b\|_{L^3}\,,
\end{equation}
where $C$ depends only on the dimension. If in addition $b\in L^2$, then
\begin{equation}\label{eq:pv-heat-l4}
 \|S(\cdot)b\|_{L^4(Q_\tau)}\le C\bigl(\|b\|_{L^2}+\|b\|_{L^3}\bigr)\,.
\end{equation}
The constants do not depend on $\tau$.
\end{lemma}

\begin{proof}
First take $b\in C_c^\infty(\mathbb R^3;\mathbb R^3)$. The Gaussian heat kernel gives a smooth solution $h=S(t)b$ with enough decay to integrate by parts in space. Testing $\partial_t h-\Delta h=0$ against $|h|h$ yields
\begin{equation*}
\frac13\frac{d}{dt}\|h(t)\|_3^3+c\int_{\mathbb R^3}|h|\,|\nabla h|^2\,dx+c\int_{\mathbb R^3}|\nabla(|h|^{3/2})|^2\,dx\le0\,,
\end{equation*}
with a numerical $c>0$. At zeros of $h$, this identity follows first with $|h|$ replaced by $(|h|^2+\eta^2)^{1/2}$ and then by letting $\eta\downarrow0$; the nonnegative terms pass by Fatou's lemma. There is no boundary contribution, since one may first insert a cutoff equal to one on $B_R$ and use the Gaussian decay as $R\to\infty$. With $g=|h|^{3/2}$, integration in time gives
\begin{equation*}
\operatorname*{ess\,sup}_{0<t<\tau}\|g(t)\|_2^2+\int_{Q_\tau}|\nabla g|^2\,dx\,dt\le C\|b\|_3^3\,.
\end{equation*}
The multiplicative Sobolev inequality proved in the argument for \cite[Lemma~8.4]{CKNLean} gives
\begin{equation*}
\int_{Q_\tau}|g|^{10/3}\,dx\,dt\le C\left(\operatorname*{ess\,sup}_{0<t<\tau}\|g(t)\|_2\right)^{4/3}\int_{Q_\tau}|\nabla g|^2\,dx\,dt\le C\|b\|_3^5\,,
\end{equation*}
which is \eqref{eq:pv-heat-l5}. Approximation in $L^3$ extends the estimate to all $b\in L^3$.

For the second estimate, the heat energy identity and $L^2$ contraction give
\begin{equation*}
\operatorname*{ess\,sup}_{0<t<\tau}\|h(t)\|_2+\|\nabla h\|_{L^2(Q_\tau)}\le C\|b\|_2\,.
\end{equation*}
The spatial inequality in \cite[Lemma~8.4]{CKNLean}, applied at almost every time, therefore gives $\|h\|_{L^{10/3}(Q_\tau)}\le C\|b\|_2$. Interpolating this bound with \eqref{eq:pv-heat-l5} yields
\begin{equation*}
\|h\|_{L^4(Q_\tau)}\le \|h\|_{L^{10/3}(Q_\tau)}^{1/2}\|h\|_{L^5(Q_\tau)}^{1/2}\le C\|b\|_2^{1/2}\|b\|_3^{1/2}\le C(\|b\|_2+\|b\|_3)\,,
\end{equation*}
which proves \eqref{eq:pv-heat-l4}. The interpolation inequality is on the product measure space $Q_\tau$ and its constant is one.
\end{proof}

\begin{lemma}[Small heat orbit]\label{lem:pv-small-heat}
If $a\in L^2(\mathbb R^3;\mathbb R^3)\cap L^3(\mathbb R^3;\mathbb R^3)$, then
\begin{equation*}
\|S(\cdot)a\|_{L^5(Q_\tau)}+\|S(\cdot)a\|_{L^4(Q_\tau)}\longrightarrow0\qquad(\tau\downarrow0)\,.
\end{equation*}
\end{lemma}

\begin{proof}
Choose $a_\rho\in C_c^\infty(\mathbb R^3;\mathbb R^3)$ with $a_\rho\to a$ in both $L^2$ and $L^3$. Such a sequence is obtained by first cutting off outside large balls and then convolving with a smooth approximate identity; truncation and translation continuity give convergence in each of the two norms. By Lemma~\ref{lem:pv-heat-critical},
\begin{equation*}
\|S(\cdot)(a-a_\rho)\|_{L^5(Q_\tau)}+\|S(\cdot)(a-a_\rho)\|_{L^4(Q_\tau)}\le C\bigl(\|a-a_\rho\|_2+\|a-a_\rho\|_3\bigr)
\end{equation*}
uniformly in $\tau$. For fixed $\rho$, the Gaussian kernel and Young's convolution inequality give, for $t>0$,
\begin{equation*}
\|S(t)a_\rho\|_5\le Ct^{-3/40}\|a_\rho\|_4\,,\qquad
 \|S(t)a_\rho\|_4\le Ct^{-1/8}\|a_\rho\|_3\,.
\end{equation*}
Integrating the fifth and fourth powers in time gives
\begin{equation*}
\|S(\cdot)a_\rho\|_{L^5(Q_\tau)}+\|S(\cdot)a_\rho\|_{L^4(Q_\tau)}
 \le C\tau^{1/8}(\|a_\rho\|_4+\|a_\rho\|_3)\,.
\end{equation*}
Given $\epsilon>0$, choose $\rho$ so that the approximation error is below $\epsilon/2$, then choose $\tau$ so that the smooth-data term is below $\epsilon/2$. This proves the limit. The pointwise heat bounds are the cases $(s,s_1)=(5,4)$ and $(4,3)$ of (7.3) in~\cite[printed p.~243]{ESS2003}; their time integrals above are computed directly, without using the general mixed-norm assertion (7.4).
\end{proof}

\section{The forced Stokes response}\label{sec:pv-stokes}

The fixed point needs an estimate that controls both the critical $L^5$ norm and an $L^4$ norm used to keep the quadratic forcing in $L^2$. The pressure in this linear equation must be included in the endpoint energy test; it is controlled by the same Riesz bounds as the force.

\begin{lemma}[Zero-data Stokes estimate]\label{lem:pv-stokes}
Let $\tau>0$ and let $F\in L^{5/2}(Q_\tau;\mathbb R^{3\times3})\cap L^2(Q_\tau;\mathbb R^{3\times3})$. There is a unique energy-class pair $(z,q)$ with zero initial datum satisfying, in distributions,
\begin{equation*}
\partial_tz-\Delta z=\operatorname{div}F-\nabla q\,,\qquad \operatorname{div}z=0\,,\qquad z(0)=0\,,
\end{equation*}
where $(\operatorname{div}F)_j=\sum_i\partial_iF_{ij}$ and $q=-P[F]$. It has
\begin{equation*}
z\in C([0,\tau];L^2)\cap L^2(0,\tau;H^1)\cap C([0,\tau];L^3)\cap L^5(Q_\tau)\cap L^4(Q_\tau),\qquad q\in L^2(Q_\tau)\cap L^{5/2}(Q_\tau),
\end{equation*}
and $\partial_tz\in L^2(0,\tau;V')$, where $V=\{\phi\in H^1(\mathbb R^3;\mathbb R^3):\operatorname{div}\phi=0\}$ and $V'$ is its dual on the divergence-free $L^2$ space. Moreover
\begin{align}
 \|z\|_{L^\infty(0,\tau;L^3)}+\|z\|_{L^5(Q_\tau)}&\le C\|F\|_{L^{5/2}(Q_\tau)}\,,\label{eq:pv-stokes-l5}\\
 \|z\|_{L^4(Q_\tau)}&\le C\bigl(\|F\|_{L^{5/2}(Q_\tau)}+\|F\|_{L^2(Q_\tau)}\bigr)\,,\label{eq:pv-stokes-l4}\\
 \|z\|_{L^\infty(0,\tau;L^2)}+\|\nabla z\|_{L^2(Q_\tau)}&\le C\|F\|_{L^2(Q_\tau)}\,.\label{eq:pv-stokes-energy}
\end{align}
The constants depend only on dimension and the pressure-operator bounds at exponents $2$ and $5/2$, and are independent of $\tau$.
\end{lemma}

\begin{proof}
For smooth tensor data, the Fourier formula in \cite[(168)]{CKNLean} gives the divergence-free evolution $z(t)=\int_0^t K(t-s)F(s)\,ds$. The pressure is $q=-P[F]$: the identity $-\Delta P[F]=\operatorname{div}\operatorname{div}F$ gives $\Delta q=\operatorname{div}\operatorname{div}F$, which is the divergence of the displayed Stokes equation. The $L^2$ energy estimate follows by pairing the projected equation with $z$ in the divergence-free energy space, where the Leray projection is orthogonal and the pressure is absent. Fourier truncation justifies this pairing for smooth data. Cauchy--Schwarz and Young's inequality give \eqref{eq:pv-stokes-energy}. The equation then gives the asserted $L^2(0,\tau;V')$ bound for $\partial_tz$.

For the critical estimate, test instead with $|z|z$. First take $F$ smooth and compactly supported in space and time. The Riesz kernel gives $q(x,t)=O(|x|^{-3})$ and $\nabla q(x,t)=O(|x|^{-4})$ outside a fixed ball; the heat representation then gives $z(x,t)=O(|x|^{-4})$ there on a finite time interval. These bounds make all cutoff errors tend to zero as the spatial cutoff radius tends to infinity. A smooth regularization at $z=0$ justifies the chain rule. The diffusion term controls $|z||\nabla z|^2$ and $|\nabla(|z|^{3/2})|^2$. Integration by parts in the force and pressure terms gives
\begin{equation*}
\frac13\frac{d}{dt}\|z(t)\|_3^3+c\int_{\mathbb R^3}|z|\,|\nabla z|^2+c\int_{\mathbb R^3}|\nabla(|z|^{3/2})|^2
 \le C\int_{\mathbb R^3}(|F|^2+|q|^2)|z|\,dx\,,
\end{equation*}
with numerical constants $c,C>0$. Indeed, $|\nabla(|z|z)|\le2|z||\nabla z|$ and $|\operatorname{div}(|z|z)|\le2|z||\nabla z|$; Cauchy--Schwarz absorbs half of the weighted diffusion term and leaves the right side displayed. Cutoff errors vanish for smooth approximants by their spatial decay. Definition~\ref{def:riesz-pressure} gives $\|q\|_{L^{5/2}(Q_\tau)}\le C\|F\|_{L^{5/2}(Q_\tau)}$.

Let $Z=\|z\|_{L^5(Q_\tau)}$, $A=\operatorname*{ess\,sup}_{0<t<\tau}\|z(t)\|_3^3$, and $D=\int_{Q_\tau}|\nabla(|z|^{3/2})|^2$. The zero initial trace and the last differential inequality imply
\begin{equation*}
A+D\le C\bigl(\|F\|_{5/2}^2+\|q\|_{5/2}^2\bigr)Z\le C\|F\|_{5/2}^2Z\,.
\end{equation*}
The multiplicative inequality used in Lemma~\ref{lem:pv-heat-critical}, now with $g=|z|^{3/2}$, gives
\begin{equation*}
Z^5\le C A^{2/3}D\le C(A+D)^{5/3}\le C\|F\|_{5/2}^{10/3}Z^{5/3}\,.
\end{equation*}
If $Z=0$ the desired estimate holds; otherwise division and taking the power $3/10$ give $Z\le C\|F\|_{5/2}$. The same differential inequality now bounds $A^{1/3}$ by $C\|F\|_{5/2}$, proving \eqref{eq:pv-stokes-l5}. These estimates first hold for smooth data and solutions.

To pass to general $F$, choose smooth tensor fields converging to $F$ in both $L^2$ and $L^{5/2}$. The $L^2$ estimate applied to differences makes the responses Cauchy in $C_tL^2_x$ and $L^2_tH^1_x$; the $L^5$ and $L^4$ estimates make them Cauchy in the corresponding space-time norms. Applying the $L^\infty_tL^3_x$ estimate to differences and using continuity in $L^3$ for smooth responses makes them Cauchy in $C_tL^3_x$, so the limit has zero strong $L^3$ trace. The pressure sequence converges in $L^2\cap L^{5/2}$ by Definition~\ref{def:riesz-pressure}. Passing to the limit in the equation against compactly supported tests proves the distributional equation; the bounds pass by completeness and weak lower semicontinuity. The energy-space solution is unique: the difference solves the homogeneous heat equation with zero data, and its $L^2$ energy identity forces it to vanish. The $L^4$ bound follows by interpolating the spatial energy estimate, which gives $\|z\|_{10/3}\le C\|F\|_2$, with \eqref{eq:pv-stokes-l5}; then $\|z\|_4\le C\|F\|_2^{1/2}\|F\|_{5/2}^{1/2}\le C(\|F\|_2+\|F\|_{5/2})$. This proves the lemma. Its endpoint estimate is the special $L^{5/2}\to L^5$ Stokes argument in Theorem~7.3 of~\cite[printed pp.~244--246]{ESS2003}, with the pressure term and approximation limit written out here.
\end{proof}

\begin{remark}[Departure from \cite{ESS2003}]\label{rem:pv-stokes-departure}
The proof of Theorem~7.3 in~\cite[printed pp.~244--246]{ESS2003} supplies the same critical estimate after testing with $|z|z$. Here the zero-data estimate is isolated as a linear response, the pressure is fixed as $-P[F]$ under Definition~\ref{def:riesz-pressure}, and the difference estimates needed for completion are recorded explicitly. No general mixed-norm estimate from Lemma~7.1 is used.
\end{remark}

\section{A solution near the initial time}\label{sec:pv-local}

For $U,V\in L^5(Q_\tau)\cap L^4(Q_\tau)$, let $\mathcal B(U,V)$ be the zero-data Stokes response of the tensor $-U\otimes V$. In this definition we use the bilinear extension of Lemma~\ref{lem:pv-stokes}; it is well-defined because $U\otimes V\in L^{5/2}\cap L^2$ by H\"older's inequality. Write $X_\tau=L^5(Q_\tau)\cap L^4(Q_\tau)$ with norm $\|U\|_{X_\tau}=\|U\|_5+\|U\|_4$. The lemma gives a constant $K$, independent of $\tau$, such that
\begin{equation}\label{eq:pv-bilinear}
 \|\mathcal B(U,V)\|_{X_\tau}\le K\|U\|_{X_\tau}\|V\|_{X_\tau}\,,\qquad
 \|\mathcal B(U,U)-\mathcal B(V,V)\|_{X_\tau}\le K(\|U\|_{X_\tau}+\|V\|_{X_\tau})\|U-V\|_{X_\tau}\,.
\end{equation}

\begin{proposition}[Short-time solution from an $L^3$ trace]\label{prop:pv-local-solution}
For every $a\in L^2(\mathbb R^3;\mathbb R^3)\cap L^3(\mathbb R^3;\mathbb R^3)\cap\J$, there is $\tau>0$ and a Leray--Hopf solution $(U,DU)$ on $\mathbb R^3\times[0,\tau]$ with datum $a$ such that $U\in L^5(Q_\tau)\cap L^4(Q_\tau)$. Moreover $U(t)\to a$ strongly in $L^3$ as $t\downarrow0$, and the associated pressure $P[U\otimes U]$ belongs to $L^2(Q_\tau)\cap L^{5/2}(Q_\tau)$.
\end{proposition}

\begin{proof}
Let $h(t)=S(t)a$. Lemma~\ref{lem:pv-small-heat} permits a choice of $\tau>0$ with $H:=\|h\|_{X_\tau}<1/(8K)$. If $H=0$, take $U=0$; otherwise the map $\Phi(U)=h+\mathcal B(U,U)$ maps the closed ball of radius $2H$ in $X_\tau$ into itself and contracts it. In fact, \eqref{eq:pv-bilinear} gives $\|\Phi(U)-h\|_{X_\tau}\le4KH^2\le H/2$ on that ball and Lipschitz constant at most $4KH<1/2$. The contraction theorem gives a unique fixed point $U$ in that ball. The heat flow and every Stokes response are divergence free, so $\operatorname{div}U=0$ in distributions.

Set $U_0=h$ and $U_{n+1}=h+\mathcal B(U_n,U_n)$. The contraction estimate gives convergence to $U$ in $X_\tau$. Since $U_n\to U$ in $L^4$, the tensors $U_n\otimes U_n$ converge to $U\otimes U$ in $L^2$; the $L^5$ convergence also gives convergence in $L^{5/2}$. Thus the forcing tensors converge in both data spaces of Lemma~\ref{lem:pv-stokes}. The associated Stokes responses converge in $C_tL^2_x\cap L^2_tH^1_x$, so $U\in C([0,\tau];L^2)\cap L^2(0,\tau;H^1)$ with $U(0)=a$ strongly in $L^2$. Their $L^\infty_tL^3_x$ difference estimate gives convergence in $C_tL^3_x$ as well. The heat flow is strongly continuous in $L^3$, and the nonlinear response tends to zero in $L^3$ at time zero; hence $U(t)\to a$ strongly in $L^3$. The pressure convergence from Lemma~\ref{lem:pv-stokes} gives $P[U\otimes U]\in L^2\cap L^{5/2}$, and passing to the limit in the mild equation gives
\begin{equation*}
\partial_tU-\Delta U+\operatorname{div}(U\otimes U)+\nabla P[U\otimes U]=0
\end{equation*}
in distributions. This is the unregularized Navier--Stokes equation; the Fourier multiplier supplies only the linear Stokes response, and the tensor in the fixed point is the actual $U\otimes U$.

It remains to verify the Leray--Hopf clauses, including the energy inequality. The response energy estimate and $U\otimes U\in L^2$ give $\partial_tU\in L^2(0,\tau;V')$. The Hilbert-space energy identity for $L^2(0,\tau;H^1)$ fields with time derivative in $L^2(0,\tau;V')$ follows by time convolution: on each compact subinterval convolve in time, pair the equation with the convolved field, and pass in the $L^2H^1$ and $L^2V'$ norms; continuity of $U$ in $L^2$ supplies the endpoint values. The nonlinear pairing is integrable since $U\otimes U\in L^2$ and $DU\in L^2$. Its spatial integral vanishes: for a cutoff $\chi_R$ equal to one on $B_R$, supported in $B_{2R}$, and with $|\nabla\chi_R|\le C/R$, integration by parts gives a convection error supported on the annulus $A_R=B_{2R}\setminus B_R$ and bounded by $C R^{-1}\int_{Q_\tau\cap A_R}|U|^3$. Hölder's inequality and $|Q_\tau\cap A_R|\le C\tau R^3$ bound this by $C\tau^{1/4}R^{-1/4}\|U\|_{L^4(Q_\tau)}^3$, which tends to zero. The diffusion cutoff error is bounded by $C R^{-1}\|U\|_{L^2(Q_\tau)}\|DU\|_{L^2(Q_\tau)}$ and also tends to zero by the energy bounds. There is no pressure boundary term in this divergence-free pairing. Thus $\int (U\otimes U):\nabla U=0$. We obtain the energy equality
\begin{equation*}
\frac12\|U(t)\|_2^2+\int_0^t\|DU(s)\|_2^2\,ds=\frac12\|a\|_2^2\qquad(0\le t\le\tau)\,,
\end{equation*}
first for almost every $t$ and then for every $t$ by $L^2$ continuity and absolute continuity of the dissipation integral. This implies (LH1), (LH2), and the stronger equality version of (LH4); the distributional equation gives (LH3), and the strong $L^2$ trace gives (LH5). The pressure term does not enter the energy pairing because we work in the divergence-free space $V$. This proves the proposition.
\end{proof}

\section{Weak--strong comparison}\label{sec:pv-serrin}

The short-time solution is a Leray--Hopf solution, so it can be compared with the given weak solution using only the latter's energy inequality. The comparison identity is obtained by mollifying the time slices of both solutions in space and removing the mollification; the $L^5$ integrability of one solution makes every nonlinear term integrable down to time zero, and (LH4) is used for the other solution only from zero.

\begin{lemma}[Weak--strong uniqueness for $L^5$ solutions]\label{lem:pv-serrin-uniqueness}
Let $T>0$, and let $(u,Du)$ and $(v,Dv)$ be Leray--Hopf solutions on $\mathbb R^3\times[0,T]$ with the same datum in $\J$. If $u\in L^5(\mathbb R^3\times(0,T))$, then $u=v$ almost everywhere on $\mathbb R^3\times(0,T)$.
\end{lemma}

\begin{proof}
By Theorem~\ref{thm:assoc-pressure}, each of $u$ and $v$ has an associated pressure in $L^{5/3}(Q_T)$ with which it satisfies the momentum equation against compactly supported smooth vector tests; the pairing of its slices with a fixed test field is continuous on $[0,T]$ and equals the pairing with $a$ at $t=0$. The spatial and temporal exponents of $u$ are both $5$, and $2/5+3/5=1$. The only nontrivial nonlinear integrability check is
\begin{align*}
\int_0^T\|v(t)\|_{10/3}\|Dv(t)\|_2\|u(t)\|_5\,dt
 &\le C\int_0^T\|v(t)\|_2^{2/5}\|Dv(t)\|_2^{8/5}\|u(t)\|_5\,dt\\
 &\le \epsilon\|Dv\|_{L^2(Q_T)}^2+C_\epsilon\|v\|_{L^\infty_tL^2_x}^2\|u\|_{L^5_tL^5_x}^5<\infty\,,
\end{align*}
where the first inequality is spatial interpolation between $L^2$ and $L^6$, and the second is Young's inequality with exponents $5/4$ and $5$. It makes the convection terms $u\cdot(v\cdot\nabla)v$ and $v\cdot(u\cdot\nabla)u$ integrable on $Q_T$, down to time zero. The gradient pairing is integrable by the energy bounds, and the pressures pair integrably with both velocities, which lie in $L^{5/2}(Q_T)$ by interpolation.

We derive the cross-testing identity without any regularity of $u$ beyond $L^5$. Let $u_n(\cdot,t)$ and $v_n(\cdot,t)$ be the spatial mollifications of the slices at radius $1/(n+1)$, and let $\chi_n$ be a smooth cutoff equal to one on $B_{n+1}$ with $|\nabla\chi_n|\le C/(n+1)$. At a fixed point $y$, testing the momentum equation of $u$ with the kernel $\rho_n(y-\cdot)e_k$ and a time cutoff shows that $t\mapsto u_n(y,t)_k$ is the primitive in time of an integrable tested flux, and likewise for $v_n$. The product rule for primitives then expresses $\int\chi_n\,v_n(t)\cdot u_n(t)\,dy-\int\chi_n\,v_n(0)\cdot u_n(0)\,dy$ as a space-time integral of mollified velocities, gradients, convection terms and pressures. The mollified slices are divergence free, so the pressure contributions only pair the mollified pressures with $\nabla\chi_n$. As $n\to\infty$, spatial mollification converges in the specific space-time $L^p$ classes used for these terms, and the product estimates above and the energy bounds give convergence of the cross integrals; the terms carrying $\nabla\chi_n$ vanish by the cutoff derivative bound and the same integrability estimates. The initial pairing converges to $\|a\|_2^2$. For almost every $t\in(0,T)$ this gives
\begin{equation*}
\int_{\mathbb R^3}v(t)\cdot u(t)\,dx-\|a\|_2^2
 =-\int_0^t\!\!\int_{\mathbb R^3}\bigl(u\cdot(v\cdot\nabla)v+\nabla u:\nabla v\bigr)
 -\int_0^t\!\!\int_{\mathbb R^3}\bigl(v\cdot(u\cdot\nabla)u+\nabla v:\nabla u\bigr)\,.
\end{equation*}
No test by an arbitrary weak solution against itself is used: the identity pairs mollified slices, and the limits use only the integrability above. With $v=u$, the transport term $\int u\cdot(u\cdot\nabla)u\,dx$ vanishes for almost every time by the divergence-free transport identity: the Leray--Hopf slice bounds and $u\in L^5(Q_T)$ provide the integrability needed to justify it by cutoff and mollification. Thus the identity is the energy equality for $u$ at almost every time. We use only (LH4) for $v$ on $[0,t]$. Combining these gives, for almost every $t\in(0,T)$,
\begin{equation*}
\frac12\|w(t)\|_2^2+\int_0^t\|Dw(s)\|_2^2\,ds
 \le\int_0^t\int_{\mathbb R^3}u\cdot(w\cdot\nabla)w\,dx\,ds\,,\qquad w=v-u\,.
\end{equation*}
The transport identities $\int b\cdot\nabla(f\cdot g)\,dx=0$ for divergence-free $b$, applied at almost every time with the integrability above, reduce the combination to this form. This is the zero-start energy comparison of~\cite[(8.12)]{RobinsonRodrigoSadowski2016}, printed p.~174, where it is derived from Serrin regularity and their testing lemma (Theorem~8.17 and Lemma~8.18, printed pp.~173--174); the mollified identity above replaces both. The weak solution contributes only its energy inequality from zero.

For almost every $s$, interpolation and H\"older's inequality give
\begin{align*}
\left|\int_{\mathbb R^3}u\cdot(w\cdot\nabla)w\,dx\right|
 &\le\|u(s)\|_5\|w(s)\|_{10/3}\|Dw(s)\|_2\\
 &\le C\|u(s)\|_5\|w(s)\|_2^{2/5}\|Dw(s)\|_2^{8/5}\\
 &\le\frac12\|Dw(s)\|_2^2+C\|u(s)\|_5^5\|w(s)\|_2^2\,.
\end{align*}
The coefficient $s\mapsto\|u(s)\|_5^5$ is integrable on $(0,T)$. Substitution and Gronwall's inequality yield
\begin{equation*}
\|w(t)\|_2^2\le \|w(0)\|_2^2\exp\!\left(C\int_0^t\|u(s)\|_5^5\,ds\right)=0\,.
\end{equation*}
The strong $L^2$ initial traces give $w(0)=0$. Gronwall therefore proves $w(t)=0$ in $L^2$ for almost every $t$, hence $u=v$ almost everywhere on the open cylinder. This proof does not use energy equality for $v$, and it does not use the positive-time regularity of $u$.
\end{proof}

\section{Positive times and the global conclusion}\label{sec:pv-assembly}

We now prove the proposed successor of Theorem~\ref{thm:ess-global}. The only estimate still needed away from $t=0$ is a velocity bound uniform in space. Local regularity gives such a bound on bounded spatial sets; at spatial infinity, the tails of the velocity and its canonical pressure make the rescaled local energy arbitrarily small at a fixed radius.

\begin{lemma}[A global bound on positive-time slabs]\label{lem:pv-positive-bound}
Under the hypotheses of Theorem~\ref{thm:ess-global}, put $M^3=\operatorname*{ess\,sup}_{t\in(0,T)}\int_{\mathbb R^3}|u(x,t)|^3\,dx$. For every $0<\delta<T$ there is
\begin{equation*}
B_\delta<\infty
\end{equation*}
depending on $\delta$, $T$, $M$, $\|a\|_2$, the absolute constants in Lemma~\ref{lem:thmA-top}, the pressure bound in Lemma~\ref{lem:assoc-pressure-L3}, and the given solution through a finite collection of local regularity neighborhoods, such that
\begin{equation*}
\operatorname*{ess\,sup}_{(x,t)\in\mathbb R^3\times(\delta,T)}|u(x,t)|\le B_\delta\,.
\end{equation*}
\end{lemma}

\begin{proof}
Let $p_*=P[u\otimes u]$ be the canonical pressure selected by Theorem~\ref{thm:assoc-pressure}. Lemma~\ref{lem:assoc-pressure-L3} gives the momentum identity (S3), the pressure equation $\Delta p_*=-\operatorname{div}\operatorname{div}(u\otimes u)$, and the bound $p_*\in L^\infty((0,T);L^{3/2}(\mathbb R^3))$. We use this same pressure in the local energy equality, so the suitability pressure is identified with the canonical pressure rather than chosen separately on each cylinder. Fix $x_0\in\mathbb R^3$, $t_0\in(0,T]$, and $\rho>0$ with $t_0-\rho^2>0$. Define the translated and rescaled fields on $B_1\times(-1,0)$ by
\begin{align*}
u_\rho(y,s)&=\rho\,u(x_0+\rho y,t_0+\rho^2s),\\
(Du)_\rho(y,s)&=\rho^2Du(x_0+\rho y,t_0+\rho^2s),\\
(p_*)_\rho(y,s)&=\rho^2p_*(x_0+\rho y,t_0+\rho^2s).
\end{align*}
By (LH1), these fields have the local energy bounds, weak gradient, and divergence condition required by Theorem~\ref{thm:ess-local}. Lemma~\ref{lem:assoc-pressure-L3} supplies the pressure identity (S3) with this same canonical pressure, and the global $L^\infty_tL^3_x$ hypothesis supplies the remaining norm bound. These properties are preserved by the displayed change of variables, so the hypotheses of Theorem~\ref{thm:ess-local} hold for the rescaled fields. Applying Lemma~\ref{lem:lei-L4} with inner radius $3/4$ proves that $(u_\rho,(Du)_\rho,(p_*)_\rho,0)$ is suitable on $B_{3/4}\times(-1,0)$. Scaling back, $(u,Du,p_*,0)$ is suitable on $B_{3\rho/4}(x_0)\times(t_0-\rho^2,t_0)$.

The function $|u|^3+|p_*|^{3/2}$ is integrable on every finite slab, and its integral outside $B_R$ tends to zero as $R\to\infty$.

Fix $r>0$ with $r^2<\delta/4$. Since $u\in L^\infty((0,T);L^3)$ and $p_*\in L^\infty((0,T);L^{3/2})$, the function $|u|^3+|p_*|^{3/2}$ has an integrable spatial tail uniformly after integration over $(\delta/2,T)$. Choose $R$ so large that for every $x_0$ with $|x_0|>R$,
\begin{equation*}
r^{-2}\iint_{B_r(x_0)\times(t_0-r^2,t_0)}(|u|^3+|p_*|^{3/2})\,dx\,dt<\varepsilon_0/8
\end{equation*}
whenever $\delta\le t_0\le T$. Indeed all these cylinders lie in $(\delta/2,T)$, and their spatial sections lie outside $B_{R-r}$; the integral is bounded by the tail integral over $(\delta/2,T)\times(\mathbb R^3\setminus B_{R-r})$, which tends to zero.
For each such $t_0$, the preceding rescaled argument with outer scale $2r$ gives suitability on $B_{3r/2}(x_0)\times(t_0-4r^2,t_0)$, since $4r^2<\delta$. This contains the smallness cylinder displayed above. Lemma~\ref{lem:thmA-top} then bounds a representative by $2C_4/r$ on $B_{r/2}(x_0)\times(t_0-r^2/4,t_0)$, with constants independent of $x_0$ and $t_0$. These cylinders cover the part of $\{|x|>R+r\}\times(\delta,T)$ with $t\le T-r^2/8$ by taking $t_0=t+r^2/8$; the remaining terminal strip is covered by the same estimate with $t_0=T$. Hence the essential supremum on the unbounded spatial part is at most $2C_4/r$.

On the bounded spatial part, Theorem~\ref{thm:ess-global} gives a regular point at every $(x,t)$ with $\delta\le t<T$. The local endpoint theorem, Theorem~\ref{thm:ess-local}, also gives a one-sided H\"older representative at every $(x,T)$ by applying it on a cylinder ending at $T$. These neighborhoods cover the compact set $\overline B_{R+r}\times[\delta,T]$; a finite subcover gives a finite bound there. Taking the larger of this bound and $2C_4/r$ proves the lemma. The dependence is through $M$, the pressure bound and energy bounds on the slab, the chosen tail radius, and the finite collection of local regularity neighborhoods covering the bounded cylinder; in particular it may depend on the given solution $u$.
\end{proof}

\begin{lemma}[Positive-time $L^5$ integrability]\label{lem:pv-positive-l5}
Under the hypotheses of Theorem~\ref{thm:ess-global}, $u\in L^5(\mathbb R^3\times(\delta,T))$ for every $\delta>0$.
\end{lemma}

\begin{proof}
Fix $\delta>0$ and take $B_{\delta/2}$ from Lemma~\ref{lem:pv-positive-bound}. Set $w=|u|^{3/2}$ on $Q_{\delta,T}$. The critical hypothesis gives $\operatorname*{ess\,sup}_{t\in(\delta,T)}\|w(t)\|_2^2\le M^3$. The weak chain rule, obtained by applying the Sobolev chain rule to smooth truncations of $z\mapsto|z|^{3/2}$ and passing by monotone convergence, gives
\begin{equation*}
\int_{Q_{\delta,T}}|\nabla w|^2\,dx\,dt\le C B_\delta\int_{Q_{\delta,T}}|Du|^2\,dx\,dt<\infty\,.
\end{equation*}
Indeed $|\nabla w|\le(3/2)|u|^{1/2}|Du|$ almost everywhere and $|u|\le B_\delta$ there. The multiplicative inequality proved in \cite[Lemma~8.4]{CKNLean}, applied to $w$ on this slab, now yields
\begin{equation*}
\int_{Q_{\delta,T}}|u|^5\,dx\,dt=\int_{Q_{\delta,T}}|w|^{10/3}\,dx\,dt
 \le C\left(\operatorname*{ess\,sup}_{t\in(\delta,T)}\|w(t)\|_2\right)^{4/3}\int_{Q_{\delta,T}}|\nabla w|^2\,dx\,dt<\infty\,.
\end{equation*}
All constants are finite because the Leray--Hopf energy inequality controls the last integral. This proves the assertion.
\end{proof}

\begin{theorem}[The $L^5$ bound and uniqueness]\label{thm:ess-l5-unique}
Let $0<T<\infty$, $a\in\J$, and let $(u,Du)$ be a Leray--Hopf solution on $\mathbb R^3\times[0,T]$ with datum $a$. Suppose
\begin{equation*}
\operatorname*{ess\,sup}_{t\in(0,T)}\int_{\mathbb R^3}|u(x,t)|^3\,dx<\infty\,.
\end{equation*}
Then $u\in L^5(\mathbb R^3\times(0,T))$, and every other Leray--Hopf solution $(v,Dv)$ on the same interval with the same datum satisfies $u=v$ almost everywhere on $\mathbb R^3\times(0,T)$.
\end{theorem}

\begin{proof}
Lemma~\ref{lem:pv-initial-l3} gives $a\in L^2\cap L^3\cap\J$, so Proposition~\ref{prop:pv-local-solution} supplies a Leray--Hopf solution $U$ on $[0,\tau]$ with datum $a$ and $U\in L^5(Q_\tau)$. Put $\sigma=\min\{\tau,T\}$ and restrict $U$ to $[0,\sigma]$. Lemma~\ref{lem:pv-serrin-uniqueness} on this interval identifies $u=U$ almost everywhere on $Q_\sigma$, so $u\in L^5(Q_\sigma)$. Choose $0<\delta<\sigma$. Lemma~\ref{lem:pv-positive-l5} gives $u\in L^5(\mathbb R^3\times(\delta,T))$. The union of these two cylinders is $Q_T$, proving the asserted global $L^5$ bound.

For any other Leray--Hopf solution $v$ with the same datum, apply Lemma~\ref{lem:pv-serrin-uniqueness} on $[0,T]$ to $u$ and $v$. It gives the required almost-everywhere equality. The energy inequality for $v$ is the only energy information used on that side of the comparison.
\end{proof}

\begin{remark}[Departure from \cite{ESS2003}]\label{rem:pv-departure}
The proof follows the positive-time reduction (3.5)--(3.7) and the short-time construction of Theorems~7.3--7.4 in~\cite{ESS2003}, printed pp.~221--222 and 244--248. We derive the initial $L^3$ trace from good times and the strong $L^2$ trace, prove the special heat estimate rather than use the general mixed-norm assertion (7.4), and write the fixed-point limit in the spaces used here. On positive-time cylinders, the translated and rescaled local energy equality of Lemma~\ref{lem:lei-L4} gives suitability with the canonical pressure; the global velocity bound and $L^5$ estimate then follow from the explicit spatial tails. Weak--strong uniqueness is proved at exponents $(5,5)$ by the zero-start energy comparison of Robinson, Rodrigo, and Sadowski~\cite[(8.12), printed p.~174]{RobinsonRodrigoSadowski2016}, with the cross-testing identity derived by spatial mollification instead of their Serrin regularity theorem and testing lemma (Theorem~8.17 and Lemma~8.18). No general mixed-norm estimate from ESS Lemma~7.1(7.4) and no result attributed to Giga is used.
\end{remark}

%======================================================================
\part*{Part VI. The Ladyzhenskaya--Prodi--Serrin theorem}
%======================================================================

\section{Statements}\label{sec:lps-statements}

For $3<s<\infty$ put $\ell(s)=2s/(s-3)$. The notation
$u\in L^{\ell(s)}(0,T;L^s(\R^3))$ refers to the mixed norm of the velocity
on the open slab, whereas the endpoint $s=\infty$ means
$u\in L^2(0,T;L^\infty(\R^3))$. In both cases
$3/s+2/\ell=1$. Equality of two weak solutions below means equality almost
everywhere on $\R^3\times(0,T)$; their prescribed time-slice representatives
need not agree at every point. The following is the whole-space theorem
in the form of \cite[Theorem~1.2, printed p.~213]{ESS2003}; that source
attributes uniqueness to Prodi and Serrin~\cite{Prodi1959,Serrin1963}
and smoothness to Ladyzhenskaya~\cite{Ladyzhenskaya1967}. Serrin's
earlier interior regularity result is~\cite{Serrin1962}. A later full-range source is
\cite[Theorems~8.17 and 8.19, printed pp.~173--174]{RobinsonRodrigoSadowski2016}.
Its Leray--Hopf class includes a strong energy inequality from good positive
times; the present class supplies the needed restart through
Lemma~\ref{lem:lps-energy-equality}.

\begin{theorem}[Ladyzhenskaya--Prodi--Serrin]\label{thm:lps}
Let $0<T<\infty$, let $a\in\J$, and let $(u,Du)$ be a Leray--Hopf
solution on $\R^3\times[0,T]$ with datum $a$. Suppose either that
$u\in L^{\ell(s)}(0,T;L^s(\R^3))$ for some $3<s<\infty$, or that
$u\in L^2(0,T;L^\infty(\R^3))$. Then every Leray--Hopf solution with
datum $a$ on $[0,T]$ equals $u$ almost everywhere on the open slab.
There is a representative $\widetilde u=u$ almost everywhere on that
slab which is $C^\infty$ on $\R^3\times(0,T]$, with all
space-time derivatives at $T$ computed within this domain.
\end{theorem}

Theorem~\ref{thm:ess-l5-unique} supplies the critical $L^5$ norm needed to
apply Theorem~\ref{thm:lps} at $(s,\ell)=(5,5)$. Thus the smoothness
mentioned in Remark~\ref{rem:not-claimed} is proved here as a separate
consequence of that theorem, while its original statement remains the
$L^5$ and uniqueness result.

\begin{corollary}[Smoothness at the $L^\infty L^3$ endpoint]\label{cor:ess-smooth}
Let $0<T<\infty$, $a\in\J$, and let $(u,Du)$ be a Leray--Hopf solution
on $\R^3\times[0,T]$ with datum $a$. If
$\esssup_{0<t<T}\int_{\R^3}|u(x,t)|^3\,dx<\infty$, then
$u\in L^5(\R^3\times(0,T))$, it is unique among Leray--Hopf solutions
with datum $a$, and it has a representative that is $C^\infty$ on
$\R^3\times(0,T]$ in the one-sided sense at $T$.
\end{corollary}

Together, Corollary~\ref{cor:ess-smooth} and Theorem~\ref{thm:lps} cover the
whole range $3\le s\le\infty$ of the regularity criterion.

\begin{corollary}[The regularity criterion for $3\le s\le\infty$]\label{cor:serrin-criterion}
Let $0<T<\infty$, $a\in\J$, and let $(u,Du)$ be a Leray--Hopf solution on
$\R^3\times[0,T]$ with datum $a$. Suppose that one of the following holds:
$\esssup_{0<t<T}\int_{\R^3}|u(x,t)|^3\,dx<\infty$; or
$u\in L^{\ell}(0,T;L^s(\R^3))$ for some $3<s<\infty$ with $\ell=2s/(s-3)$; or
$u\in L^2(0,T;L^\infty(\R^3))$. Then every Leray--Hopf solution with datum $a$ on
$[0,T]$ equals $u$ almost everywhere on $\R^3\times(0,T)$, and $u$ has a
representative that is $C^\infty$ on $\R^3\times(0,T]$ in the one-sided sense at $T$.
\end{corollary}

\begin{proof}
In the first case apply Corollary~\ref{cor:ess-smooth}; in the other two, apply
Theorem~\ref{thm:lps}.
\end{proof}

\section{Energy equality and comparison}\label{sec:lps-energy}

The weak equation has tests supported away from the time endpoints. A Serrin bound makes the convection term integrable when the solution tests its own equation, and the energy equality is obtained from the mollified cross-testing identity of Part~V with the solution in both slots. Definition~\ref{def:leray-hopf} may prescribe nonmeasurable slices on exceptional times. The energy equality and its positive-time restart are therefore asserted only at almost every measurable time, without changing that definition. A whole-space energy identity under $L^4$ control also appears in
\cite[Lemma~4.3, printed pp.~72--73]{Tsai2018}. For comparison,
\cite[Theorem~12.4, printed pp.~359--361]{LemarieRieusset2016}
treats weak--strong uniqueness. The proof below works with almost-every-time
slices, as required by the measurability clauses of the Leray--Hopf definition.

\begin{lemma}[Mixed-norm interpolation]\label{lem:lps-mixed-interpolation}
Under the hypotheses of Theorem~\ref{thm:lps}, set
$q=2s/(s-2)$ and $r=2s/3$ when $s<\infty$, and set $(q,r)=(2,\infty)$
when $s=\infty$. Then $u\in L^r(0,T;L^q(\R^3))$ and
$u\in L^4(\R^3\times(0,T))$. The same energy-class estimate
$w\in L^r_tL^q_x$ holds for any Leray--Hopf solution $w$ on $[0,T]$.
\end{lemma}

\begin{proof}
For almost every $t$, the weak gradient clause of (LH1) and the
whole-space homogeneous Sobolev inequality give
$\|w(t)\|_6\le C_S\|Dw(t)\|_2$. For finite $s$, put
$\theta=3/s$. The identities
$1/q=(1-\theta)/2+\theta/6$ and $r\theta=2$ give
\begin{equation}\label{eq:lps-energy-interpolation}
\int_0^T\|w(t)\|_q^r\,dt
\le C_S^2\left(\esssup_{0<t<T}\|w(t)\|_2\right)^{r-2}
       \int_0^T\|Dw(t)\|_2^2\,dt<\infty\,.
\end{equation}
The vector Sobolev inequality follows componentwise and the equivalence of
the coordinate and Euclidean norms changes only $C_S$. At $s=\infty$ the
assertion is exactly $w\in L^\infty_tL^2_x$.
For finite $s$, $1/s+1/q=1/2$ and $1/\ell(s)+1/r=1/2$.
H\"older's inequality first in space and then in time therefore gives
$\|u^2\|_{L^2_{t,x}}\le
\|u\|_{L^{\ell(s)}_tL^s_x}\|u\|_{L^r_tL^q_x}$.
The same calculation at $s=\infty$ uses
$L^2_tL^\infty_x$ and $L^\infty_tL^2_x$. Hence $u^2\in L^2$,
which is the asserted space-time $L^4$ bound.
\end{proof}

\begin{lemma}[Good positive times]\label{lem:lps-good-times}
If $(u,Du)$ is a Leray--Hopf solution on $[0,T]$, then there is a
full-measure set $G\subset(0,T)$ such that for every $t_0\in G$,
$u(\cdot,t_0)\in H^1(\R^3;\R^3)\cap\J$, with weak gradient
$Du(\cdot,t_0)$; this slice is measurable as an $L^2$ field.
Every nonempty open subinterval of $(0,T)$ contains a time in $G$.
\end{lemma}

\begin{proof}
Fubini applied to (LH1) gives finite $L^2$ norms of $u(t)$ and
$Du(t)$ for almost every $t$ and identifies $Du(t)$ as the weak
spatial gradient there. Intersect these times with those on which the
weak-divergence clause of (LH1) holds; by
Lemma~\ref{lem:J-weak-div}, $u(t)\in\J$. This intersection has full
measure. An open interval has positive Lebesgue measure, so it meets
that set.
\end{proof}

\begin{lemma}[Energy chain rule for square-integrable time derivatives]\label{lem:lps-Hilbert-triple}
Let $a<b$, put $Q=\R^3\times(a,b)$, and let $w,r\in L^2(Q)$ satisfy
$\dt w=r$ in the sense of distributions on $Q$.
\begin{enumerate}[label=(\roman*)]
\item The function $w$ has a representative $\bar w\in C([a,b];L^2(\R^3))$,
and for $a\le s\le t\le b$
\[
\|\bar w(t)\|_2^2-\|\bar w(s)\|_2^2
=2\int_s^t\!\!\int_{\R^3}w\,r\,dx\,d\tau .
\]
\item Let $j\in\{1,2,3\}$ and suppose in addition that the weak spatial
derivatives $g=\partial_jw$ and $h=\partial_jg$ on $Q$ belong to $L^2(Q)$.
Then $g$ has a representative $\bar g\in C([a,b];L^2(\R^3))$, and for
$a\le s\le t\le b$
\[
\|\bar g(t)\|_2^2-\|\bar g(s)\|_2^2
=-2\int_s^t\!\!\int_{\R^3}h\,r\,dx\,d\tau .
\]
\end{enumerate}
The same statements hold componentwise for vector fields.
\end{lemma}

\begin{proof}
(i) Regard $w$ as an element of $L^2(a,b;L^2(\R^3))$. Its distributional
time derivative $r$ lies in the same space, so
$w\in W^{1,2}(a,b;L^2(\R^3))$. Such a function agrees for almost every
time with the absolutely continuous curve
$\bar w(t)=\bar w(a)+\int_a^tr\,d\tau$, which is continuous on the closed
interval, endpoints included. The product rule for absolutely continuous
Hilbert-space valued curves gives the displayed identity.

(ii) Let $\rho_\varepsilon$ be a standard spatial mollifier, and for
$\varepsilon,\varepsilon'>0$ put $\kappa=\rho_\varepsilon-\rho_{\varepsilon'}$;
convolutions act in $x$ at almost every time. The field
$d=\kappa*g=(\partial_j\kappa)*w$ lies in $L^2(Q)$, because
$\|\kappa*g(t)\|_2\le\|\kappa\|_1\|g(t)\|_2$, and
$\dt d=(\partial_j\kappa)*r=\partial_j(\kappa*r)\in L^2(Q)$. Part~(i)
applies to $d$, with representative $\bar d$: for $s\le t$,
\[
\|\bar d(t)\|_2^2-\|\bar d(s)\|_2^2
=2\int_s^t\!\!\int\partial_j(\kappa*r)\,(\kappa*g)\,dx\,d\tau
=-2\int_s^t\!\!\int(\kappa*r)\,(\kappa*h)\,dx\,d\tau ,
\]
where $\kappa*g\in H^1$ has $\partial_j(\kappa*g)=\kappa*h$ at almost every
time, and the derivative was moved onto this factor. Hence
$|\|\bar d(t)\|_2^2-\|\bar d(s)\|_2^2|\le2\|\kappa*r\|_{L^2(Q)}\|\kappa*h\|_{L^2(Q)}$.
Averaging over $s\in(a,b)$ gives for every $t\in[a,b]$
\[
\|\bar d(t)\|_2^2\le(b-a)^{-1}\|d\|_{L^2(Q)}^2
+2\|\kappa*r\|_{L^2(Q)}\|\kappa*h\|_{L^2(Q)} .
\]
As $\varepsilon,\varepsilon'\to0$ the fields $\rho_\varepsilon*g$,
$\rho_\varepsilon*r$ and $\rho_\varepsilon*h$ converge in $L^2(Q)$ to $g$, $r$
and $h$, so the right side tends to zero. The curves $\bar g_\varepsilon$
representing $\rho_\varepsilon*g$ are therefore uniformly Cauchy in
$C([a,b];L^2(\R^3))$, and their limit $\bar g$ is a representative of $g$.
The same computation with $\kappa$ replaced by $\rho_\varepsilon$ gives the
identity for $\bar g_\varepsilon$, and it passes to the limit
$\varepsilon\to0$.
\end{proof}

\begin{lemma}[Energy equality at almost every time]\label{lem:lps-energy-equality}
Under the hypotheses of Theorem~\ref{thm:lps}, there is a full-measure
set $E\subset(0,T)$ of times at which the slice $u(\cdot,t)$ is a
measurable $L^2$ field such that for every $t\in E$,
\begin{equation}\label{eq:lps-energy-equality}
\frac12\|u(t)\|_2^2+
\int_0^t\|Du(\tau)\|_2^2\,d\tau
=\frac12\|a\|_2^2\,.
\end{equation}
Consequently, for $t_0,t\in E$ with $t_0<t$,
\begin{equation}\label{eq:lps-energy-restart}
\frac12\|u(t)\|_2^2+
\int_{t_0}^t\|Du(\tau)\|_2^2\,d\tau
=\frac12\|u(t_0)\|_2^2\,.
\end{equation}
No equality is asserted for arbitrary prescribed time slices.
\end{lemma}

\begin{proof}
Put $Q_T=\R^3\times(0,T)$. Lemma~\ref{lem:lps-mixed-interpolation} gives
$u\in L^4(Q_T)$, hence, with $u\in L^2(Q_T)$, also $u\in L^{5/2}(Q_T)$ and
$u\in L^3(Q_T)$. By Theorem~\ref{thm:assoc-pressure} there is
$p\in L^{5/3}(Q_T)$ with which $u$ satisfies the momentum equation against
compactly supported smooth vector tests, and the pairing of the slices of
$u$ with a fixed test field is continuous on $[0,T]$ and equals the pairing
with $a$ at $t=0$ by (LH2) and (LH5). These are the inputs of the mollified
cross-testing calculation in the proof of
Lemma~\ref{lem:pv-serrin-uniqueness}, which we apply with $u$ in both slots
and with $L^4$ in place of $L^5$. Its products are integrable:
$|u|^2|Du|\in L^1(Q_T)$ because $1/4+1/4+1/2=1$, and the pressure pairs
integrably with $u\in L^{5/2}(Q_T)$ because $3/5+2/5=1$. The mollified
slices converge in $L^4(Q_T)$ and their gradients in $L^2(Q_T)$, and the
cutoff terms tend to zero by the bound $C/R$ on the cutoff derivative and
the same integrability. The initial pairing converges to $\|a\|_2^2$.
For almost every $t\in(0,T)$ we obtain
\begin{equation*}
\|u(t)\|_2^2-\|a\|_2^2
=-2\int_0^t\!\!\int u_iu_j\,\partial_ju_i\,dx\,d\tau
 -2\int_0^t\|Du(\tau)\|_2^2\,d\tau .
\end{equation*}
No time derivative of $u$ in a negative Sobolev space and no chain rule in
time are used.

It remains to show that the transport integral vanishes. By Fubini's
theorem and Lemma~\ref{lem:lps-good-times}, for almost every $\tau$ the
slice $u(\tau)$ lies in $H^1(\R^3)\cap\J\cap L^4(\R^3)$, with weak gradient
$Du(\tau)$. For such a slice $|u|^2\in W^{1,4/3}_{\rm loc}$, because
$|u||Du|\in L^{4/3}$. Let $\chi_R$ be a smooth cutoff equal to one on
$B_R$, supported in $B_{2R}$, with $|\nabla\chi_R|\le C/R$. Since
$u\in L^4=(L^{4/3})^*$, the weak-divergence identity extends from
smooth compactly supported tests to $\chi_R|u|^2/2\in W^{1,4/3}$, and gives
$\int\chi_Ru_iu_j\partial_ju_i\,dx=-\frac12\int|u|^2u\cdot\nabla\chi_R\,dx$.
The right side is bounded by
$CR^{-1}\int_{B_{2R}\setminus B_R}|u|^3\le CR^{-1}|B_{2R}|^{1/4}\|u\|_4^3
\to0$. Since $u_iu_j\partial_ju_i\in L^1(\R^3)$ for such a slice, dominated
convergence gives $\int u_iu_j\partial_ju_i\,dx=0$. The space-time integral
of $u_iu_j\partial_ju_i$ is defined because $|u|^2|Du|\in L^1(Q_T)$, and it
vanishes because its slice integrals do. Let $E$ be the set of good times of
Lemma~\ref{lem:lps-good-times} at which the displayed identity holds. Then
$E$ has full measure and the identity is \eqref{eq:lps-energy-equality};
subtraction gives \eqref{eq:lps-energy-restart}.
\end{proof}

The comparison below is stated for two Leray--Hopf solutions with the same
datum at time zero. A comparison after a positive restart time is obtained in
Lemma~\ref{lem:lps-continuation} by concatenating $u$ with the strong
solution and applying this lemma to the concatenation. The nonlinear estimate
uses the same exponent $q$ as Lemma~\ref{lem:lps-mixed-interpolation};
the gradient of the difference remains in $L^2$.

\begin{lemma}[Mixed-norm comparison]\label{lem:lps-comparison}
Let $0<T<\infty$, $a\in\J$, and let $(u,Du)$ and $(v,Dv)$ be Leray--Hopf
solutions on $\R^3\times[0,T]$ with datum $a$. Suppose that $u$ satisfies a
mixed-norm hypothesis of Theorem~\ref{thm:lps} on $(0,T)$. Then $v=u$
almost everywhere on $\R^3\times(0,T)$.
\end{lemma}

\begin{proof}
Put $Q=\R^3\times(0,T)$. By Theorem~\ref{thm:assoc-pressure}, $u$ and $v$
have pressures $p_u,p_v\in L^{5/3}(Q)$ with which they satisfy the
momentum equation against compactly supported smooth vector tests, and the
pairing of their slices with a fixed test field is continuous on $[0,T]$ and
equals the pairing with $a$ at $t=0$. By
Lemma~\ref{lem:lps-energy-equality}, $u$ satisfies
\eqref{eq:lps-energy-equality} for almost every $t$, and by (LH4) the field
$v$ satisfies the corresponding inequality with $Dv$.

For either energy field $w=u,v$, the interpolation in
\eqref{eq:lps-energy-interpolation} gives $w\in L^r_tL^q_x$.
H\"older with $1/s+1/q+1/2=1$ and
$1/\ell+1/r+1/2=1$ makes each of
$|u||v||Du|$ and $|u||v||Dv|$ integrable.
It also controls $|u|^2|Du|$; independently,
$u\in L^4(Q)$ makes $|u|^2|Dz|$ integrable for $z=v-u$.
At $s=\infty$ the same bounds use $(\ell,q,r)=(2,2,\infty)$.
These are precisely the nonlinear factors retained below; no term
$|v|^2|Du|$ is needed.

Use the pressure-inclusive weak equations with a time cutoff, spatially
mollified fields, and a spatial cutoff $\chi_R$. Before removing the
mollifier, write the convection of $v$ tested against $u_n$ in
advective form $\int(v\cdot\nabla)v\cdot u_n$, using
$\dv v=0$. Thus one never takes a limit in
$\int(v\otimes v):\nabla u_n$. For the $u$ equation tested against
$v_n$, retain $u\otimes u:\nabla v_n$: the tensor is in $L^2(Q)$
because $u\in L^4(Q)$, and $\nabla v_n\to\nabla v$ strongly in
$L^2(Q)$, for every $s$. For the $v$ equation tested against
$u_n$, the advective factor $(v\cdot\nabla)v$ belongs to
$L^{\ell'}_tL^{s'}_x$ by energy interpolation. When $s<\infty$,
spatial mollification gives $u_n\to u$ strongly in
$L^\ell_tL^s_x$, so this second term also converges. At
$s=\infty$, strong convergence in $L^2_tL^\infty_x$ is not
asserted: $u_n\to u$ almost everywhere and
$|u_n(x,t)|\le\|u(t)\|_\infty$. Dominated convergence applies
to the second term because
$\|u(t)\|_\infty\|v(t)\|_2\|Dv(t)\|_2\in L^1_t$.

The cutoff remainders also vanish. Besides $|u|^3$ and
$|u||v|^2$, the transport integrations produce $|v||u|^2$.
Here $u\in L^{5/2}(Q)$ by interpolation of $L^2$ and $L^4$,
$v\in L^{10/3}(Q)$ by the energy interpolation, and the mixed
H\"older estimate with one $u$ in $L^\infty_tL^2_x$ handles
$|v||u|^2$ on the finite time interval. The cutoff tests are not
divergence free. Their pressure errors are bounded by
\[
\frac C R\int_Q |p_v||u|+|p_u||v|,
\]
which is finite because $p_u,p_v\in L^{5/3}(Q)$, while $u$ and $v$ lie in
$L^{5/2}(Q)$ (by interpolation of $L^2$ with $L^4$ and with $L^{10/3}$
respectively) and $3/5+2/5=1$.
Every cutoff error tends to zero as $R\to\infty$.
The common datum gives the initial pairing $\|a\|_2^2$ by (LH2) and (LH5);
take the terminal limit through times where the weak slices hold. The result
is, for almost every $t\in(0,T)$,
\begin{equation*}
\langle u(t),v(t)\rangle+2\int_{0}^t\!\int Du:Dv
=\|a\|_2^2
-\int_{0}^t\!\int
  \bigl((u\cdot\nabla)u\cdot v+(v\cdot\nabla)v\cdot u\bigr).
\end{equation*}
This is the mollified cross-testing calculation of
Lemma~\ref{lem:pv-serrin-uniqueness} with its integrable products
made explicit for the mixed exponents.

Set $z=v-u$ and let $I$ denote the nonlinear integral in the last display.
The cancellation is carried out only
with integrable terms. First,
$\int(u\cdot\nabla)u\cdot u=0$, using
$|u|^2|Du|\in L^1$ and cutoff remainder $|u|^3$.
Second, $\int(v\cdot\nabla)u\cdot u=0$, using
$|v||u||Du|\in L^1$ and remainder $|v||u|^2$.
Third, $\int(u\cdot\nabla)z\cdot u
=-\int(u\cdot\nabla)u\cdot z$, using
$|u|^2|Dz|$ and $|u||z||Du|$.
Expanding $v=u+z$ in $I$ and applying these three identities gives
$I=\int_0^t\!\int(z\cdot\nabla)z\cdot u$, where
$(z\cdot\nabla)z\cdot u=\sum_ku_k\sum_jz_j\partial_jz_k$. The cross-testing
identity therefore takes the relative form
\begin{equation}\label{eq:lps-cross-testing}
\langle v(t),u(t)\rangle-\|a\|_2^2
=-\int_0^t\!\int(z\cdot\nabla)z\cdot u\,dx\,d\tau
 -2\int_0^t\!\int Du:Dv\,dx\,d\tau
\end{equation}
for almost every $t\in(0,T)$.
Now expand $\|z(t)\|_2^2=\|v(t)\|_2^2-2\langle v(t),u(t)\rangle+\|u(t)\|_2^2$
and insert, for almost every $t$, the energy inequality (LH4) for $v$,
\eqref{eq:lps-cross-testing}, and the energy equality
\eqref{eq:lps-energy-equality} for $u$:
\[
\|z(t)\|_2^2\le
\Bigl(\|a\|_2^2-2\!\int_0^t\!\|Dv\|_2^2\Bigr)
-2\Bigl(\|a\|_2^2-\!\int_0^t\!\!\int(z\cdot\nabla)z\cdot u
-2\!\int_0^t\!\!\int Du:Dv\Bigr)
+\Bigl(\|a\|_2^2-2\!\int_0^t\!\|Du\|_2^2\Bigr),
\]
which equals
$2\int_0^t\!\int(z\cdot\nabla)z\cdot u-2\int_0^t\|Dz\|_2^2$. Thus, for
almost every $t$,
\begin{equation}\label{eq:lps-relative-energy}
\frac12\|z(t)\|_2^2+\int_{0}^{t}\|Dz(\tau)\|_2^2\,d\tau
\le\int_{0}^{t}\int_{\R^3}(z\cdot\nabla)z\cdot u\,dx\,d\tau
\le\int_{0}^{t}\int_{\R^3}|u|\,|z|\,|Dz|\,dx\,d\tau\,.
\end{equation}
For finite $s$, whole-space Sobolev interpolation gives
$\|z\|_q\le C\|z\|_2^{1-3/s}\|Dz\|_2^{3/s}$.
Young's inequality with conjugate exponents
$2/(1+3/s)$ and $2/(1-3/s)=\ell$ yields
\begin{equation}\label{eq:lps-relative-young}
\|u\|_s\|z\|_q\|Dz\|_2
\le\frac12\|Dz\|_2^2+C_s\|u\|_s^\ell\|z\|_2^2\,.
\end{equation}
At $s=\infty$, the same result follows from
$\|u\|_\infty\|z\|_2\|Dz\|_2
\le\frac12\|Dz\|_2^2+C\|u\|_\infty^2\|z\|_2^2$.
Put $m(\tau)=\|u(\tau)\|_s^\ell$ (with $m=\|u\|_\infty^2$ at $s=\infty$) and
$\phi(t)=\|z(t)\|_2^2$. By hypothesis $m\in L^1(0,T)$, and $\phi$ is
essentially bounded by the energy class. Inserting the last estimate in
\eqref{eq:lps-relative-energy} and absorbing the gradient term gives, in
integrated form, $\phi(t)\le K\int_0^tm\phi\,d\tau$ for almost every $t$,
with $K=2C_s$ (an absolute constant at $s=\infty$). The function
$\Phi(t)=\int_0^tm\phi\,d\tau$ is absolutely continuous with
$\Phi'=m\phi\le Km\Phi$ almost everywhere and $\Phi(0)=0$, so
$\Phi=0$, and then $\phi=0$ almost everywhere. Hence $z=0$ almost
everywhere on $Q$. This is the whole-space
relative-energy mechanism of
\cite[Theorem~8.19 and (8.12), printed p.~174]{RobinsonRodrigoSadowski2016};
no smoothness of $u$ was used.
\end{proof}

\section{Strong solutions and continuation}\label{sec:lps-strong}

The Fourier--$L^2$ mild construction of the CKN paper~\cite{CKNLean} supplies smooth solutions
when the transport velocity is mollified. The next proposition takes the
mollifier to zero while retaining an $H^1$ bound independent of its
radius. This is the whole-space strong-solution construction underlying
\cite[Theorem~6.15, printed pp.~144--145]{RobinsonRodrigoSadowski2016};
it avoids the expanding-ball compactness passage of that source.

\begin{lemma}[Regularized Sobolev start]\label{lem:lps-regularized-Hk-start}
Let $b\in\J$ and $\varepsilon>0$, let $U_\varepsilon$ and $p_\varepsilon$ be the
solution and pressure of \cite[Theorem~9.4]{CKNLean} with datum $b$, and
put $V_\varepsilon=\Jeps U_\varepsilon$. For every finite $S>0$:
\begin{enumerate}[label=(\alph*)]
\item for every integer $k\ge0$, $U_\varepsilon\in C([0,S];H^k)$ and
$U_\varepsilon(0)=\Jeps b$ in $H^k$;
\item for $t>0$ the time derivative of $U_\varepsilon$ exists classically and
$\dt U_\varepsilon=R_\varepsilon-\nabla p_\varepsilon$, where
$R_\varepsilon=\Delta U_\varepsilon-(V_\varepsilon\cdot\nabla)U_\varepsilon$
is a continuous map from $(0,S]$ into $L^2(\R^3;\R^3)$;
\item the function $y_\varepsilon(t)=\|\nabla U_\varepsilon(t)\|_2^2$ is
continuous on $[0,S]$ and differentiable on $(0,S)$, and for $0<t<S$
\begin{equation}\label{eq:lps-regularized-Hm-identity}
y_\varepsilon'(t)
=-2\|\nabla^2U_\varepsilon(t)\|_2^2
 +2\big\langle(V_\varepsilon(t)\cdot\nabla)U_\varepsilon(t),
   \Delta U_\varepsilon(t)\big\rangle .
\end{equation}
\end{enumerate}
The constants of these statements may depend on $\varepsilon$; the transport
bound in Proposition~\ref{prop:lps-local-strong} does not. Identity
\eqref{eq:lps-regularized-Hm-identity} is the only case of the
higher-order energy identities that is used below.
\end{lemma}

\begin{proof}
(a) \cite[Lemma~9.2]{CKNLean} puts the initial state
$\Jeps b$ in every $H^k$. For each fixed $k$, the $H^k$ heat and
projected Stokes multiplier bounds in the proof of
\cite[Theorem~9.4]{CKNLean}, together with
\cite[(174)]{CKNLean}, make the mild map contractive on a short
interval in $C_tH^k$. The initial heat term converges in $H^k$ at
$t=0$ by dominated convergence of the Fourier multiplier. The Stokes
integral tends to zero there because its kernel is bounded by
$C_\varepsilon(t-s)^{-1/2}$ and its norm by
$2C_\varepsilon\sqrt t\sup_{s\le t}\|U_\varepsilon(s)\|_{H^k}$.
Mild uniqueness identifies this path with the $L^2$ solution. At each
restart the $L^2$ energy bound fixes a positive step length for the
regularized equation; the finite restart argument in the proof of
\cite[Theorem~9.4]{CKNLean} extends the $H^k$ path to $S$.

(b) The pointwise equation is (R3) of \cite[Theorem~9.4]{CKNLean}. By (a)
with $k=2$, $t\mapsto\Delta U_\varepsilon(t)$ is continuous into $L^2$. By
\cite[Lemma~9.2]{CKNLean},
$\|V_\varepsilon(t)-V_\varepsilon(s)\|_\infty\le C_\varepsilon
\|U_\varepsilon(t)-U_\varepsilon(s)\|_2$ and
$\|V_\varepsilon(s)\|_\infty\le C_\varepsilon\|b\|_2$, so
\[
\|(V_\varepsilon(t)\cdot\nabla)U_\varepsilon(t)
 -(V_\varepsilon(s)\cdot\nabla)U_\varepsilon(s)\|_2
\le\|V_\varepsilon(t)-V_\varepsilon(s)\|_\infty\|\nabla U_\varepsilon(t)\|_2
 +\|V_\varepsilon(s)\|_\infty\|\nabla U_\varepsilon(t)-\nabla U_\varepsilon(s)\|_2 ,
\]
which tends to zero as $s\to t$ by (a) with $k=1$. Hence $R_\varepsilon$ is
continuous into $L^2$.

(c) Continuity of $y_\varepsilon$ on $[0,S]$ is (a) with $k=1$. For $f,g\in H^2$
the symmetry of the Laplacian gives
$\|\nabla f\|_2^2-\|\nabla g\|_2^2=\langle f-g,-\Delta(f+g)\rangle$.
Let $0<a<c\le S$, take $f=U_\varepsilon(c)$, $g=U_\varepsilon(a)$, and put
$w=-\Delta(U_\varepsilon(a)+U_\varepsilon(c))\in L^2$; it is weakly
divergence free, because both slices are. By (R2),
$\dt U_\varepsilon\in L^2(\R^3\times(a,c))$ is continuous, so
$U_\varepsilon(c)-U_\varepsilon(a)=\int_a^c\dt U_\varepsilon\,dr$
pointwise, and Fubini's theorem against the fixed field $w$ gives
$\langle U_\varepsilon(c)-U_\varepsilon(a),w\rangle
=\int_a^c\langle R_\varepsilon(r)-\nabla p_\varepsilon(r),w\rangle\,dr$.
For almost every $r$ we have $p_\varepsilon(r)\in L^2$ and
$\nabla p_\varepsilon(r)=R_\varepsilon(r)-\dt U_\varepsilon(r)\in L^2$, so
$p_\varepsilon(r)\in H^1$ and $\langle\nabla p_\varepsilon(r),w\rangle
=-\langle p_\varepsilon(r),\dv w\rangle=0$. Therefore
\[
y_\varepsilon(c)-y_\varepsilon(a)
=\int_a^c\big\langle R_\varepsilon(r),
 -\Delta(U_\varepsilon(a)+U_\varepsilon(c))\big\rangle\,dr .
\]
Only the pressure-free field $R_\varepsilon$ needs continuity in time; no time
derivative of $\nabla U_\varepsilon$ is taken. Divide by $c-a$ and let
$a\uparrow t$ and $c\downarrow t$ for a fixed $0<t<S$. Continuity of $R_\varepsilon$
from (b) and of $\Delta U_\varepsilon$ from (a) give
$y_\varepsilon'(t)=2\langle R_\varepsilon(t),-\Delta U_\varepsilon(t)\rangle$,
which is \eqref{eq:lps-regularized-Hm-identity} because
$\|\Delta U_\varepsilon(t)\|_2=\|\nabla^2U_\varepsilon(t)\|_2$ by Plancherel.
This argument starts at $t=0$ in $H^1$ because the mollified initial state is
in that space; it is not inferred from (R1)--(R2) alone.
\end{proof}

\begin{proposition}[Local strong solution]\label{prop:lps-local-strong}
For every $b\in H^1(\R^3;\R^3)\cap\J$ there is a
$\tau_b>0$, depending only on $\|b\|_{H^1}$, a solenoidal solution $U$
of the unregularized equation on $[0,\tau_b]$, and a pressure
$p\in L^2(\R^3\times(0,\tau_b))$ for which the momentum equation with
pressure holds in the weak sense against compactly supported smooth vector
tests, such that
\begin{equation}\label{eq:lps-strong-class}
U\in C([0,\tau_b];H^1)\cap L^2(0,\tau_b;H^2),
\qquad \dt U\in L^2(0,\tau_b;L^2),\qquad U(0)=b\,.
\end{equation}
It has energy equality from every time and defines a Leray--Hopf
solution on this interval. A choice of $\tau_b$ bounded below by
$c(1+\|b\|_{H^1})^{-4}$ is possible, with absolute $c>0$.
\end{proposition}

\begin{proof}
Apply \cite[Theorem~9.4]{CKNLean} with datum $b$ and mollifier
radius $\varepsilon\downarrow0$. Lemma~\ref{lem:lps-regularized-Hk-start}
adds the $H^1$ continuity at zero and the differentiated gradient-energy
identity \eqref{eq:lps-regularized-Hm-identity} needed here; these are not
clauses (R1)--(R5) of that theorem. Write $y=\|\nabla U_\varepsilon\|_2^2$
and $d=\|\nabla^2U_\varepsilon\|_2^2=\|\Delta U_\varepsilon\|_2^2$. The
transport velocity $\Jeps U_\varepsilon$ is divergence free, and Young's
inequality for convolution and the homogeneous $H^1\to L^6$ inequality give
$\|\Jeps U_\varepsilon\|_6\le C\|\nabla U_\varepsilon\|_2$. By H\"older's
inequality and the interpolation of $\nabla U_\varepsilon$ between $L^2$ and
$L^6$, whose $L^6$ norm is controlled by $\|\nabla^2U_\varepsilon\|_2$,
\[
|\langle(\Jeps U_\varepsilon\cdot\nabla)U_\varepsilon,\Delta U_\varepsilon\rangle|
\le\|\Jeps U_\varepsilon\|_6\|\nabla U_\varepsilon\|_3\|\Delta U_\varepsilon\|_2
\le C\,y^{3/4}d^{3/4}
\]
with an absolute constant, independent of $\varepsilon$ and of the mollifier
profile. Inserting this in \eqref{eq:lps-regularized-Hm-identity} and using
Young's inequality with exponents $4/3$ and $4$ gives, for $t>0$,
\begin{equation}\label{eq:lps-uniform-H1}
\frac{d}{dt}\|\nabla U_\varepsilon\|_2^2
 +\|\Delta U_\varepsilon\|_2^2
\le C\|\nabla U_\varepsilon\|_2^6\,.
\end{equation}
The initial field is $\Jeps b$, so its $L^2$ and $H^1$ norms do not
exceed those of $b$. Set $X_0=1+\|b\|_{H^1}^2$. Since $y$ is continuous on $[0,S]$ and
satisfies $y'\le Cy^3$ on $(0,S)$ for every $S$, comparison of
$X'=CX^3$ with $X(0)=X_0$ (applied on $[s,\tau]$ and letting $s\downarrow0$)
shows that for
$\tau_b=cX_0^{-2}$, with $c$ chosen so that $2cC<1/2$,
$\sup_{0\le t\le\tau_b}\|\nabla U_\varepsilon(t)\|_2^2
\le2X_0$. Integrating \eqref{eq:lps-uniform-H1} over $[s,\tau_b]$, using this
bound and letting $s\downarrow0$ controls $\int_0^{\tau_b}\|\Delta U_\varepsilon\|_2^2$
independently of $\varepsilon$. The global $L^2$ energy identity adds
$\sup_t\|U_\varepsilon(t)\|_2^2\le\|b\|_2^2$.

The nonlinear term belongs uniformly to $L^2_tL^2_x$:
$\|(\Jeps U_\varepsilon\cdot\nabla)U_\varepsilon\|_2
\le\|\Jeps U_\varepsilon\|_\infty\|\nabla U_\varepsilon\|_2
\le C\|U_\varepsilon\|_{H^2}\|U_\varepsilon\|_{H^1}$.
The projected equation therefore bounds $\dt U_\varepsilon$ in
$L^2_tL^2_x$; the Leray projection has operator norm one on $L^2$.
\cite[Proposition~12.1]{CKNLean}, applied to this family with datum
$b$, supplies a subsequence converging strongly in $L^2$ on bounded
space-time cylinders and weakly in the energy gradients. Its compactness
input is \cite[Lemma~10.6]{CKNLean}, which gives the required
time-pairing modulus; the present $L^2_tL^2_x$ time-derivative bound is
an additional estimate, not a replacement for that input. The uniform
$H^2$ bound improves local convergence
to strong $L^2_tH^1_x$: for a cutoff $\chi$ and
$W=\chi(U_\varepsilon-U)$, integration by parts gives
$\|\nabla W\|_2^2\le\|W\|_2\|\Delta W\|_2$;
the first factor tends to zero and the second is bounded. Weak
compactness retains $U\in L^\infty_tH^1_x\cap L^2_tH^2_x$ and
$\dt U\in L^2_tL^2_x$. The convergence
$\Jeps U_\varepsilon-U\to0$ locally in $L^2_tH^1_x$ follows by
writing it as $\Jeps(U_\varepsilon-U)+(\Jeps U-U)$ and using the
uniform derivative bound and approximate-identity convergence. The
quadratic transport term therefore converges in distributions to
$(U\cdot\nabla)U$; pressure is recovered by the fixed double-Riesz
formula of \cite[Section~8]{CKNLean}, and lies in
$L^2(\R^3\times(0,\tau_b))$ because $U\otimes U\in L^\infty_tL^2_x$ and the
Riesz projection is bounded on $L^2$. The initial value follows from
the uniform time-pairing modulus and $\Jeps b\to b$ in $L^2$.

For the time continuity, apply Lemma~\ref{lem:lps-Hilbert-triple}
componentwise. Here $U\in L^2_tH^2_x$, so the weak derivatives $\partial_jU$
and $\partial_j\partial_jU$ lie in $L^2$, and $\dt U\in L^2_tL^2_x$. Part~(i)
of the lemma gives the $L^2$ continuity of $U$ together with
$\|U(t)\|_2^2-\|U(s)\|_2^2=2\int_s^t\!\int U\cdot\dt U$. Part~(ii), applied to
each derivative $\partial_jU_i$ and summed over $i,j$, gives the $L^2$
continuity of $\nabla U$ together with
$\|\nabla U(t)\|_2^2-\|\nabla U(s)\|_2^2
=-2\int_s^t\!\int\Delta U\cdot\dt U$. Thus $U\in C_tH^1_x$. Choose this
continuous representative for the strong solution;
it agrees almost everywhere with the limit field of
\cite[Proposition~12.1]{CKNLean}. The resulting energy identity
gives the $L^2$ trace and, after pairing
the projected equation with $-\Delta U$, the $H^1$ trace. The
nonlinearity is in $L^2_tL^2_x$, so the equation and spatial cutoffs
justify the energy equality for every subinterval. These properties,
weak continuity, and the weak equation verify (LH1)--(LH5) with datum
$b$; in particular no positive-time inequality is inserted as an
extra hypothesis.
\end{proof}

\begin{lemma}[Serrin control of the $H^1$ norm]\label{lem:lps-H1-estimate}
Let $U$ be a strong solution on $[t_0,t_1]$, that is, a solenoidal
field in the class \eqref{eq:lps-strong-class} with an $L^2$ pressure as in
Proposition~\ref{prop:lps-local-strong}, with
$t_0<t_1\le T$, and suppose it has either mixed norm in
Theorem~\ref{thm:lps} there. For finite $s$, with
$\ell=2s/(s-3)$, there is a constant $C_s$ depending only on $s$ and
the whole-space Sobolev constant such that
\begin{equation}\label{eq:lps-H1-gronwall}
\|\nabla U(t)\|_2^2+
\int_{t_0}^{t}\|\Delta U(\tau)\|_2^2\,d\tau
\le\|\nabla U(t_0)\|_2^2
\exp\left(C_s\int_{t_0}^{t}\|U(\tau)\|_s^\ell\,d\tau\right)
\end{equation}
for $t\in[t_0,t_1]$. At $s=\infty$ the same estimate holds with
$\ell=2$ and an absolute constant.
\end{lemma}

\begin{proof}
Pair the momentum equation with $-\Delta U$. This is legitimate for
$U\in L^2H^2$ and $\dt U\in L^2L^2$ by part~(ii) of
Lemma~\ref{lem:lps-Hilbert-triple}, which gives
$\|\nabla U(t)\|_2^2-\|\nabla U(s)\|_2^2=-2\int_s^t\langle\dt U,\Delta U\rangle$.
At almost every time $\dt U=\Delta U-(U\cdot\nabla)U-\nabla p$ in $L^2$, and
$\langle\nabla p,\Delta U\rangle=0$ because $\Delta U$ is solenoidal and
$\nabla p\in L^2$ is a gradient of an $L^2$ function; hence
$\frac{d}{dt}\|\nabla U\|_2^2=-2\|\Delta U\|_2^2
+2\langle(U\cdot\nabla)U,\Delta U\rangle$ for almost every $t$. If
$3<s<\infty$, let $q=2s/(s-2)$. H\"older and homogeneous Sobolev
interpolation applied to $\nabla U$, whose $L^6$ norm is controlled by the
full Hessian norm $\|\nabla^2U\|_2=\|\Delta U\|_2$ (Plancherel), yield
\begin{equation}\label{eq:lps-H1-nonlinear}
\left|\int(U\cdot\nabla)U\cdot\Delta U\right|
\le\|U\|_s\|\nabla U\|_q\|\Delta U\|_2
\le C_s\|U\|_s\|\nabla U\|_2^{1-3/s}
             \|\Delta U\|_2^{1+3/s}\,.
\end{equation}
Young's inequality with powers $2/(1+3/s)$ and
$2/(1-3/s)=\ell$ bounds the last expression by
$\|\Delta U\|_2^2/2+C_s\|U\|_s^\ell\|\nabla U\|_2^2$.
At $s=\infty$, use
$\|U\|_\infty\|\nabla U\|_2\|\Delta U\|_2$
and Young with powers two. The coefficients are integrable on the
interval. Indeed,
$Y(t)=\|\nabla U(t)\|_2^2+
\int_{t_0}^t\|\Delta U(\tau)\|_2^2\,d\tau$
satisfies $Y'(t)\le C_s\|U(t)\|_s^\ell Y(t)$ for almost every $t$
(with $\ell=2$ at $s=\infty$). Gronwall gives the displayed
estimate with this same $C_s$. This is the whole-space calculation of
\cite[Lemma~8.16, printed pp.~171--172]{RobinsonRodrigoSadowski2016},
including its separately displayed $L^2_tL^\infty_x$ case.
\end{proof}

\begin{lemma}[Continuation beyond the comparison interval]\label{lem:lps-continuation}
Let $(u,Du)$ satisfy the hypotheses of Theorem~\ref{thm:lps}, and
choose $t_0<T$ in the intersection of the good-time set of
Lemma~\ref{lem:lps-good-times} and the equality set of
Lemma~\ref{lem:lps-energy-equality}. Let $S$ be the set of
$t'>t_0$ for which there is a strong solution $U$ on $[t_0,t']$
with $U(t_0)=u(t_0)$ in $L^2$ and $U=u$ almost everywhere on
$\R^3\times(t_0,\min\{t',T\})$. Then $S$ contains a time larger
than $T$.
\end{lemma}

\begin{proof}
\emph{Step 1: comparison after a restart by concatenation.} Let $t'>t_0$ and
let $W$ be a strong solution on $[t_0,t']$ with $W(t_0)=u(t_0)$ in $L^2$. We
claim that $W=u$ almost everywhere on $\R^3\times(t_0,\sigma)$,
$\sigma=\min\{t',T\}$. Let $V$ be the field equal to $u$ at times $t<t_0$ and
to $W$ at times $t\ge t_0$, and let $DV$ be defined in the same way from $Du$
and $DW$. We check that $(V,DV)$ is a Leray--Hopf solution with datum $a$ on
$[0,\sigma]$.

(LH1): $V$ and $DV$ are measurable and square integrable on
$\R^3\times(0,\sigma)$, being so for $u$ on $(0,t_0)$ and for $W$ on
$(t_0,\sigma)$; $\|V(t)\|_2$ is essentially bounded, because $u$ is and
$W\in C([t_0,\sigma];L^2)$; the weak-gradient and weak-divergence clauses hold
at almost every time from those of $u$ and the strong class. (LH2): the
pairings with a fixed $w\in L^2$ are continuous for $t\le t_0$ by (LH2) for
$u$, and continuous on $[t_0,\sigma]$ by the strong $L^2$ continuity of $W$;
they match at $t_0$ because $W(t_0)=u(t_0)$, a measurable slice since
$t_0$ is a good time. (LH3): let $\varphi\in\Cc(\R^3\times(0,\sigma);\R^3)$ be
divergence free. Testing the equation of $u$ with $\varphi$ times a time
cutoff tending to the indicator of $(0,t_0)$, and using (LH2), shows that the
weak-equation integral of $u$ over $(0,t_0)$ equals $-\langle
u(t_0),\varphi(t_0)\rangle$. The strong equation of $W$ (pressure terms drop
against divergence-free tests) and $W\in C([t_0,\sigma];L^2)$ show that the
integral of $W$ over $(t_0,\sigma)$ equals $+\langle
W(t_0),\varphi(t_0)\rangle$. The two boundary pairings agree and cancel, which
is the weak equation for $V$ on $(0,\sigma)$ with no mass at $t_0$. (LH4): for
$t\le t_0$ this is (LH4) for $u$. For $t\ge t_0$, the energy identity of the
strong solution (valid for every strong solution: apply part~(i) of Lemma~\ref{lem:lps-Hilbert-triple} with $r=\dt W=\Delta W-(W\cdot\nabla)W-\nabla p$, and use $\langle\Delta W,W\rangle=-\|\nabla W\|_2^2$, $\langle(W\cdot\nabla)W,W\rangle=0$, $\langle\nabla p,W\rangle=0$), $\frac12\|W(t)\|_2^2+\int_{t_0}^t\|DW\|_2^2
=\frac12\|W(t_0)\|_2^2$, together with
\eqref{eq:lps-energy-equality} at $t_0\in E$ and $W(t_0)=u(t_0)$, gives
$\frac12\|V(t)\|_2^2+\int_0^t\|DV\|_2^2=\frac12\|a\|_2^2$. (LH5): the strong
trace at zero is that of $u$.

The restriction of $(u,Du)$ to $[0,\sigma]$ is a Leray--Hopf solution with the
same datum, and it satisfies the mixed-norm hypothesis on $(0,\sigma)$ since
$\sigma\le T$. Lemma~\ref{lem:lps-comparison} on $[0,\sigma]$, with $u$ in the
Serrin role and $V$ as the rival solution, gives $V=u$ almost everywhere on
$\R^3\times(0,\sigma)$, hence $W=u$ almost everywhere on
$\R^3\times(t_0,\sigma)$. The pressures entering the comparison are those of
Theorem~\ref{thm:assoc-pressure} for $u$ and for $V$; the strong solution
enters only through the Leray--Hopf class of $V$.

\emph{Step 2: the set $S$ is nonempty.}
Proposition~\ref{prop:lps-local-strong} translated to $t_0$ gives a strong
solution $W$ on $[t_0,t_0+\tau_b]$ with $W(t_0)=u(t_0)\in H^1\cap\J$. By
Step~1, $W=u$ almost everywhere on the common interval, so
$t_0+\tau_b\in S$.

\emph{Step 3: extension.}
Suppose every $t'\in S$ satisfies $t'\le T$. On any such interval,
$U=u$ almost everywhere, so Lemma~\ref{lem:lps-H1-estimate} gives
\[
\sup_{t_0\le t\le t'}\|U(t)\|_{H^1}^2
 +\int_{t_0}^{t'}\|\Delta U(t)\|_2^2\,dt
\le M,
\]
where $M$ depends only on $\|u(t_0)\|_{H^1}$,
$\|u(t_0)\|_2$, and the Serrin integral on $(t_0,T)$, not on $t'$.
The $L^2$ part is bounded by the strong energy identity.
Proposition~\ref{prop:lps-local-strong}, translated to any endpoint
$t'\in S$, therefore provides a lifespan $\tau_*>0$ bounded below
in terms of $M$ alone. Restart from $U(t')\in H^1\cap\J$ and glue
the two strong fields. Continuity in $H^1$ at the junction eliminates
a distributional time jump; the two $L^2$ time derivatives remain
integrable on the union, and the pressures are glued in the same way. The
glued field is a strong solution on $[t_0,t'+\tau_*]$ with the same value
$u(t_0)$ at $t_0$, so Step~1 shows that it agrees with $u$ almost everywhere
on $\R^3\times(t_0,\min\{t'+\tau_*,T\})$. Thus $t'+\tau_*\in S$.
If $\sup S\le T$, choose $t'\in S$ with
$t'>\sup S-\tau_*/2$; the preceding extension exceeds $\sup S$,
a contradiction. Hence some member of $S$ exceeds $T$.
This uses the continuation remark after
\cite[Corollary~6.9, printed p.~137]{RobinsonRodrigoSadowski2016};
on $\R^3$ the local construction is the one underlying
\cite[Theorem~6.15, printed pp.~144--145]{RobinsonRodrigoSadowski2016}.
\end{proof}

\section{Smooth representatives}\label{sec:lps-smooth}

The strong solution identifies the weak field on every interval after a
good time. On a positive-time slab it then gains spatial regularity by a
direct linear ladder: a bound for $\dt U$ in $L^2_x$ gives the $H^2$
regularity of every slice, and the whole-space heat gain, applied to the
projected convection term as a source, gives one further order at a time.
The projected equation supplies all time derivatives. The endpoint $t_1$ is
included by working on closed time intervals throughout. On $\R^3$ there is no
boundary compatibility condition.

\begin{lemma}[Joint smoothness from Sobolev curves]\label{lem:lps-Bochner-joint-smooth}
Let $t_0<t_1$ and let $V$ be a space-time $L^2_{\rm loc}$ field on
$\R^3\times(t_0,t_1)$. Suppose that for every $\delta\in(0,t_1-t_0)$ there
are curves $Z^j\in C([t_0+\delta,t_1];H^m(\R^3))$, for all integers $j,m\ge0$
(each $Z^j(t)$ being one element of $\bigcap_mH^m$), such that
$Z^0(t)=V(\cdot,t)$ for almost every $t$, and for every
$\psi\in C_c^\infty(\R^3)$ and $j\ge0$ the function
$t\mapsto\langle Z^j(t),\psi\rangle$ is differentiable on $[t_0+\delta,t_1]$,
one-sidedly at the endpoints, with derivative
$\langle Z^{j+1}(t),\psi\rangle$. Then $V$ has a representative that is
jointly $C^\infty$ on $\R^3\times(t_0,t_1]$, with one-sided time derivatives
at $t_1$.
\end{lemma}

\begin{proof}
Fix $\delta$. The continuous Sobolev embedding $H^m\hookrightarrow C_b^k$,
$m>k+3/2$, turns each curve $Z^j$ into a continuous curve in $C_b^k$ for every
$k$. Since $Z^{j+1}$ is continuous with values in $H^m$, the
differentiability of the pairings with all test functions gives
$Z^j(t)-Z^j(s)=\int_s^tZ^{j+1}(\tau)\,d\tau$ as a Bochner integral in $H^m$,
hence in $C_b^k$, so $Z^j$ is a $C^1$ curve into $C_b^k$ with derivative
$Z^{j+1}$ (one-sided at the endpoints). Evaluation of a $C_b^k$ function at
$x$ is jointly $C^k$ in the function and the $x$ variable. Composing the
evaluation with the curves shows that
$(x,t)\mapsto Z^0(t)(x)$ has continuous mixed partial derivatives
$\partial_t^j\partial_x^\alpha=\partial_x^\alpha Z^j(t)(x)$ of every
order $j+|\alpha|$, jointly in $(x,t)$ on $\R^3\times[t_0+\delta,t_1]$; it is
therefore $C^\infty$ there, with one-sided derivatives at the ends of the
interval. It equals $V$ almost everywhere. For $\delta'<\delta$ the two
smooth functions agree almost everywhere on the common set, hence everywhere
by continuity, so they patch as $\delta\downarrow0$ to a representative on
$\R^3\times(t_0,t_1]$.
For formalization, the embedding into $C_b^k$ can be constructed by
Fourier inversion and differentiation under the integral, as in the
$H^4\hookrightarrow C_b^2$ argument of the CKN paper~\cite{CKNLean}, iterated to higher
orders.
\end{proof}

\begin{proposition}[Smoothing of strong solutions]\label{prop:lps-smoothing}
If $U$ is a strong solution on $[t_0,t_1]$ with $t_0<t_1$, that is, a solenoidal
field in the class \eqref{eq:lps-strong-class} with an $L^2$ pressure as in
Proposition~\ref{prop:lps-local-strong}, then for every $\delta>0$ with
$t_0+\delta<t_1$ and every integer $m\ge0$,
$U\in C([t_0+\delta,t_1];H^m(\R^3))$ and
$U\in L^2(t_0+\delta,t_1;H^{m+1}(\R^3))$.
It has a representative smooth jointly in space and time on
$\R^3\times(t_0,t_1]$, with one-sided time derivatives at $t_1$.
\end{proposition}

\begin{proof}
Write $Q=\R^3\times(t_0,t_1)$. The strong class gives
$U,\nabla U,\nabla^2U,\dt U\in L^2(Q)$ and $U\in C([t_0,t_1];H^1)$. At almost
every time $\dt U=\Delta U-(U\cdot\nabla)U-\nabla p$ with all terms in $L^2$.
The fields $\dt U$ and $\Delta U$ are solenoidal and $\nabla p$ is a gradient,
so the Helmholtz decomposition gives the projected equation
\begin{equation}\label{eq:lps-projected}
\dt U-\Delta U=-\Leray B,\qquad B=(U\cdot\nabla)U,
\end{equation}
in which the pressure no longer appears.

\emph{Step 1: $\dt U$ is bounded in $L^2_x$ after positive time.}
Let $M(t)=\|U(t)\|_{H^2}^2$, so $M\in L^1(t_0,t_1)$, and let $K$ be the
constant of the embedding $H^2\hookrightarrow L^\infty$. For small $h>0$ put
$v_h(t)=h^{-1}(U(t+h)-U(t))$ on $[t_0,t_1-h]$. Since
$U(t+h)\otimes U(t+h)-U(t)\otimes U(t)$ equals
$U(t+h)\otimes(U(t+h)-U(t))+(U(t+h)-U(t))\otimes U(t)$, the difference of
the equations at times $t+h$ and $t$, divided by $h$, states that $v_h$ is
solenoidal and solves $\dt v_h-\Delta v_h+\nabla\cdot F_h+\nabla q_h=0$, with
flux $|F_h|\le(|U(t+h)|+|U(t)|)|v_h|$. We have $v_h\in L^2H^2$ and
$\dt v_h\in L^2L^2$, so part~(i) of Lemma~\ref{lem:lps-Hilbert-triple}
gives the energy identity for $E_h=\|v_h\|_2^2$; the pressure term vanishes
against the solenoidal $v_h$, and
\[
E_h'\le-2\|\nabla v_h\|_2^2
 +(\|U(t+h)\|_\infty+\|U(t)\|_\infty)\,2\|v_h\|_2\|\nabla v_h\|_2
\le-\|\nabla v_h\|_2^2+\kappa_hE_h,
\]
where $\kappa_h(t)=2K^2\bigl(M(t)+M(t+h)\bigr)$ has integral at most
$4K^2\|M\|_{L^1}$, uniformly in $h$. Moreover
$\|v_h(s)\|_2^2\le h^{-1}\int_0^h\|\dt U(s+\theta)\|_2^2\,d\theta$, so
$\int_{t_0}^{t_0+\delta}E_h\le\|\dt U\|_{L^2(Q)}^2$. Averaging the
Gronwall inequality $E_h(t)\le E_h(s)\exp\int\kappa_h$ over
$s\in(t_0,t_0+\delta)$ gives, for $t\in[t_0+\delta,t_1-h]$,
\[
E_h(t)\le B_\delta:=\delta^{-1}e^{4K^2\|M\|_{L^1}}\|\dt U\|_{L^2(Q)}^2 ,
\]
independently of $h$. Since $v_h\to\dt U$ in $L^2(Q)$, weak lower
semicontinuity gives $\|\dt U(t)\|_2^2\le B_\delta$ for almost every
$t\in(t_0+\delta,t_1)$.

\emph{Step 2: $H^2$ bound of the slices and the base of the ladder.}
At almost every time $U(t)\in H^2$ solves the stationary system
$\Delta U=\dt U+(U\cdot\nabla)U+\nabla p$. Test it with $\Delta U$: the
pressure drops because $\Delta U$ is solenoidal, and
$|\langle(U\cdot\nabla)U,\Delta U\rangle|\le\|U\|_6\|\nabla U\|_3\|\Delta U\|_2
\le C\|\nabla U\|_2^{3/2}\|\nabla^2U\|_2^{3/2}$ by Sobolev interpolation, with
$\|\nabla^2U\|_2=\|\Delta U\|_2$. Young's inequality gives
$\|\nabla^2U(t)\|_2\le2\|\dt U(t)\|_2+C\|\nabla U(t)\|_2^3$. By Step~1 and
$U\in C([t_0,t_1];H^1)$, $\sup\|U(t)\|_{H^2}\le C_\delta$ for almost every
$t\in(t_0+\delta,t_1)$. Consequently $B=(U\cdot\nabla)U$ is bounded in $H^1$
there: $|\nabla B|\le|\nabla U|^2+|U||\nabla^2U|$, with
$\|\nabla U\|_4^2\le C\|U\|_{H^2}^2$ and
$\|U\|_\infty\|\nabla^2U\|_2\le C\|U\|_{H^2}^2$.

We use the whole-space heat gain. \emph{Let $m\ge1$, $a<a+\sigma<a+\beta$,
let $z\in L^2((a,a+\beta);H^m(\R^3))$ and $G\in L^2((a,a+\beta);H^{m-1}(\R^3))$
satisfy $\dt z-\Delta z=G$ in distributions. Then
$z\in C([a+\sigma,a+\beta];H^m)\cap L^2((a+\sigma,a+\beta);H^{m+1})$ with norm
at most $C(\|z\|_{L^2H^m}+\|G\|_{L^2H^{m-1}})$, where $C$ depends only on
$m$, $\sigma$ and $\beta$.} This is the whole-space form of
Lemma~\ref{lem:local-heat-gain}: multiply by a time cutoff vanishing for
$t\le a+\sigma/2$ and equal to one for $t\ge a+\sigma$; each derivative
$\partial^\gamma$, $|\gamma|\le m$, of the product solves a heat equation with
square-integrable source or divergence-form source and zero initial trace, and
the energy estimate of Lemma~\ref{lem:localized-vorticity-energy} with $v=0$
applies with no spatial cutoff. The estimate includes the value at the
terminal time $a+\beta$ in its supremum, so the closed terminal interval is
covered.

Apply it, using the bounds above with $\delta/2$ in place of $\delta$, with
$m=2$ on $(a,a+\beta)=(t_0+\delta/2,t_1)$ and $\sigma=\delta/2$
to $z=U$ and $G=-\Leray B$; here $\|\Leray B\|_{H^1}\le C\|B\|_{H^1}$ because
$\Leray$ is a bounded Fourier multiplier. This gives $U\in C([t_0+\delta,t_1];H^2)\cap
L^2(t_0+\delta,t_1;H^3)$ for every $\delta$.

\emph{Step 3: the ladder.} We show by induction on $m$ that for every
$\delta>0$ with $t_0+\delta<t_1$, $U\in C([t_0+\delta,t_1];H^m)\cap
L^2(t_0+\delta,t_1;H^{m+1})$. The cases $m\le1$ are contained in the strong
class and $m=2$ is Step~2. Let $m\ge2$ and assume the statement for $m$.
Apply it with $\delta/2$ in place of $\delta$. Then $U$ is bounded in
$H^m$ and $\nabla U\in L^2H^m$ on $(t_0+\delta/2,t_1)$, so by the algebra
inequality $\|fg\|_{H^m}\le C_m\|f\|_{H^m}\|g\|_{H^m}$ (valid since
$m>3/2$), $B=\sum_jU_j\partial_jU\in L^2H^m$ there, and $\Leray B\in L^2H^m$.
Now $U\in L^2H^{m+1}$ solves \eqref{eq:lps-projected} with source
$-\Leray B\in L^2H^m$. The heat gain of order $m+1$ with
$(a,a+\beta)=(t_0+\delta/2,t_1)$ and $\sigma=\delta/2$ gives
$U\in C([t_0+\delta,t_1];H^{m+1})\cap L^2(t_0+\delta,t_1;H^{m+2})$.
This proves all the spatial assertions on the closed interval $[t_0+\delta,t_1]$.

\emph{Step 4: the time chain.} Fix $\delta$ and put $I=[t_0+\delta,t_1]$. Define
$Z^0=U$ and, recursively,
\[
Z^{j+1}=\Delta Z^j-\Leray\sum_{i=0}^{j}\binom ji
\,(Z^i\cdot\nabla)Z^{j-i}.
\]
If $Z^0,\dots,Z^j\in C(I;H^m)$ for every $m$, then $Z^{j+1}\in C(I;H^m)$ for
every $m$, by the algebra property and the boundedness of $\Leray$ on $H^m$;
so by Step~3 and induction all $Z^j$ belong to $C(I;H^m)$ for every $m$.
By \eqref{eq:lps-projected}, $\dt Z^0=Z^1$ in the sense of distributions on
$(t_0+\delta,t_1)$. If $\dt Z^i=Z^{i+1}$ for $i\le j$, then the product rule
for products of curves that are continuous into $H^m$ and have continuous
distributional time derivatives (by time mollification), together with
Pascal's rule for the binomial coefficients, gives $\dt Z^{j+1}=Z^{j+2}$. For
each $\psi\in C_c^\infty(\R^3)$ the pairing $\langle Z^j(t),\psi\rangle$ is
thus a $C^1$ function on the open interval with the continuous derivative
$\langle Z^{j+1}(t),\psi\rangle$, and therefore is differentiable on the closed
interval $I$, one-sidedly at the ends, with this derivative.

\emph{Step 5: joint smoothness.} Step 4 gives exactly the hypotheses of
Lemma~\ref{lem:lps-Bochner-joint-smooth}, for every $\delta$, with
$V=U$; the identification $Z^0(t)=U(t)$ for almost every $t$ was
established in Step~3. The lemma produces a representative smooth jointly on
$\R^3\times(t_0,t_1]$ with one-sided derivatives at $t_1$. No continuation
beyond $t_1$ is used for the endpoint: all estimates hold on the closed
intervals $[t_0+\delta,t_1]$.
This is the whole-space version of
\cite[Theorems~7.1, 7.3 and 7.5, printed pp.~149--154]{RobinsonRodrigoSadowski2016};
the argument above uses a direct heat-gain ladder and the Fourier
multipliers of the CKN paper~\cite{CKNLean} in place of the periodic Galerkin basis used for the
detailed source proof.
\end{proof}

\begin{remark}[An a posteriori energy inequality]\label{rem:lps-Hm-energy}
Once Proposition~\ref{prop:lps-smoothing} is proved, $U\in C^1(I;H^m)$, $I=[t_0+\delta,t_1]$, for
every $m$ and $\delta$ by Step~4, and for $m\ge2$ with
$\|U\|_{H^m}^2=\sum_{|\alpha|\le m}\|\partial^\alpha U\|_2^2$ pairing the
projected equation with $\partial^\alpha U$ and integrating one derivative by
parts in the divergence form of the convection term gives
\begin{equation}\label{eq:lps-Hm-energy}
\frac{d}{dt}\|U\|_{H^m}^2+\|\nabla U\|_{H^m}^2\le C_m\|U\|_{H^m}^4 .
\end{equation}
Indeed the heat pairing gives $-\|\nabla U\|_{H^m}^2$, the projection is
orthogonal and drops out against the solenoidal $\partial^\alpha U$, and
the tensor-gradient pairing is bounded by
$C_m\|U\|_{H^m}^2\|\nabla U\|_{H^m}$ by the algebra property, after which
Young's inequality is used. This inequality is not used in the proof of the
proposition, which does not pass through regularized solutions at higher
order.
\end{remark}

\begin{proof}[Proof of Theorem~\ref{thm:lps}]
Lemmas~\ref{lem:lps-mixed-interpolation} and
\ref{lem:lps-energy-equality} give energy equality for $u$ at almost
every measurable time. At $t_0=0$, the strong initial trace (LH5)
gives the common $L^2$ trace $a$ through those times for $u$ and for
any rival Leray--Hopf solution. Theorem~\ref{thm:assoc-pressure-forced}
and the $L^4$ estimate give the $L^2$ canonical pressure for $u$;
Theorem~\ref{thm:assoc-pressure} gives an $L^{5/3}$ pressure for the
rival. Lemma~\ref{lem:lps-comparison} with $t_0=0$ identifies them
almost everywhere. The rival's (LH4) supplies its energy inequality.

Choose $t_0$ in the intersection of the full-measure good-time and
energy-equality sets of Lemmas~\ref{lem:lps-good-times} and
\ref{lem:lps-energy-equality}. Subtracting two instances of
\eqref{eq:lps-energy-equality} gives the equality from $t_0$ for
almost every later time. The trace assertion of the energy lemma
identifies the strong $L^2$ trace there with $u(t_0)$; no positive-time
energy clause is added to the stated Leray--Hopf definition.
Proposition~\ref{prop:lps-local-strong} constructs a strong solution
from this datum. It belongs to $L^2_tL^\infty_x$ by
$H^2\hookrightarrow L^\infty$ and its $L^2_tH^2$ bound.
Lemma~\ref{lem:lps-comparison}, with that strong field in the Serrin
role and the original weak field in the energy role, identifies them
on their common interval. Lemma~\ref{lem:lps-continuation} repeats
this identification at each restart and extends the strong field
beyond $T$.

For any $\delta\in(0,T)$ choose such a $t_0<\delta/2$.
Proposition~\ref{prop:lps-smoothing} gives a smooth representative on
$\R^3\times[\delta,T]$. Representatives from two choices of good
time agree almost everywhere on the overlap. Since both are
continuous, they agree everywhere there: a nonempty open set on
which they differ would have positive measure. The representatives
thus patch as $\delta\downarrow0$. The strong continuation past $T$
supplies the one-sided space-time derivatives there.
\end{proof}

\begin{proof}[Proof of Corollary~\ref{cor:ess-smooth}]
Theorem~\ref{thm:ess-l5-unique} gives $u\in L^5(Q_T)$ and its
uniqueness conclusion. Fubini identifies this spacetime norm with
$L^5(0,T;L^5(\R^3))$, and $3/5+2/5=1$. Apply
Theorem~\ref{thm:lps} with $(s,\ell)=(5,5)$ to obtain the smooth
representative through the top time.
\end{proof}


\appendix
\section{Corrections to the sources}\label{app:errata}

\subsection*{Escauriaza--Seregin--Šverák (2003)}

\begin{enumerate}
\item \textbf{Cylinder-sign typo.} Theorem~1.4, printed p.~214, writes the ambient cylinder as \(B_1\times(0,1)\), while its Sobolev hypothesis and target backward cylinder use negative times. The coherent cylinder is \(B_1\times(-1,0)\). The same inconsistency occurs in the IMA preprint. The manuscript states the negative-time domain in Theorem~\ref{thm:ess-local} and proves the shifted-cylinder argument in Lemma~\ref{lem:regular-point-shift}.

\item \textbf{Test-space mismatch.} Definition~2.1, printed p.~215, names spatial tests \(\phi\in C^\infty_0(\mathbb R^3)\), but its local energy formula uses \(\phi(x,t)\) and \(\partial_t\phi\). The IMA preprint repeats the mismatch. The precise space-time test class is not fixed there. The manuscript fixes its suitable-solution test clauses in Definition~\ref{def:sws}.

\item \textbf{Higher-order regularity delegated.} Lemma~2.2, printed p.~216, asserts bounds for every spatial derivative order but refers \(k>1\) to an external Stokes regularity result instead of proving the induction. Its proof ends on printed p.~220; §3 begins on p.~221. Constants are indexed by \(k\), but their dependence is not developed. Theorem~\ref{thm:vorticity-regularity} supplies the finite weak-to-regular bootstrap needed here, but does not reproduce the source's all-orders theorem or supply the all-orders exterior estimates in (3.31), (3.38), and (3.42)--(3.46).

\item \textbf{Compressed Campanato compactness.} Propositions~2.4--2.6 and the proof of Lemma~2.2, printed pp.~216--220, appeal to “well-known compactness” and leave the relevant function-space operators and pressure-limit details compressed. No result of the manuscript relies on this argument; the manuscript uses the CKN regularity criterion~\cite{CKNLean} and proves the compactness it needs in \cite[Lemma~11.1]{CKNLean}.

\item \textbf{Time representative and all-time \(L^3\) bound.} The assertion following (3.1), printed p.~221, extends weak continuity from the Leray--Hopf representative and the essential \(L^\infty_tL^3_x\) bound to every time, but does not spell out the density and lower-semicontinuity argument. The manuscript writes that representative and modulus in Lemma~\ref{lem:weak-cont-L3}.

\item \textbf{Strong regularity estimates are used without a detailed proof.} The local estimate (3.10), printed p.~222, is cited through mixed-norm coercive estimates and supplies much more than the weak continuity and \(L^3\)-slice bound needed there. The manuscript replaces those conclusions by the momentum-pairing modulus in Lemma~\ref{lem:weak-cont-L3}; this does not claim the full derivative package in (3.10).

\item \textbf{Singular-point lower bound is externally imported.} The lower bound (3.13), printed p.~223, is attributed to Seregin--Šverák [39], a result not used here. The contrapositive of CKN Theorem~A does not replace it directly: Theorem~A gives full-neighborhood regularity below the top face but only one-sided Hölder control at the top, so its contrapositive does not yield ESS's lower bound at a full-neighborhood singular point. The manuscript avoids the singular-point formulation. Lemma~\ref{lem:thmA-top} proves top-boundary \(\varepsilon\)-regularity with threshold \(\varepsilon_0/8\), and Definition~\ref{def:good-point} defines good points by that smallness. Lemma~\ref{lem:good-open-glue} shows that the solution is Hölder on a compact set of good points, and that a point which is not good satisfies the lower bound \eqref{eq:bad-point-lower-bound} at every admissible scale. The blow-up argument of Proposition~\ref{prop:blowup-limit} uses only this lower bound, so no input from [39] is needed.

\item \textbf{Blow-up compactness conclusion is stronger than the replacement lemma.} The estimates (3.21)--(3.23), printed pp.~224--225, invoke nonstationary Stokes coercivity to obtain derivative bounds, strong \(L^3\) convergence, and \(C([a,b];L^2_{\rm loc})\) convergence. The cited nonstationary Stokes coercivity is not proved in detail. The manuscript compactness lemma \cite[Lemma~11.1]{CKNLean} gives strong local space-time \(L^2\), weak gradients, and weak slice convergence at every time; it does not claim the stronger \(C_tL^2_{\rm loc}\) topology. The blow-up proof here instead handles its endpoint and strong-\(L^3\) passages using pairing continuity, weak slice convergence, lower semicontinuity, and time-strip estimates.

\item \textbf{Pressure-splitting normalization matters.} Equations (3.11)--(3.19) begin on printed p.~223 and end on p.~224. In these equations, \(p_1\) is defined by a local Dirichlet problem and \(p_2\) is that split's harmonic remainder. The decay in (3.19) belongs to this fixed remainder; it cannot be transferred to a different harmonic remainder chosen after rescaling. The manuscript's Lemma~\ref{lem:pressure-split} instead fixes a whole-space Riesz pressure of the zero-extended tensor and its harmonic difference, then proves the interior bound and rescaled vanishing directly. It does not transfer (3.19) to that changed split or redefine the remainder after rescaling.

\item \textbf{Time reversal in the vorticity equation.} The §3 vorticity equation and (3.32), printed pp.~226--227, use \(\partial_t-\Delta\); §5 is written for \(\partial_t+\Delta\) in its forward variable. Reversing time reconciles them. This is a convention to state when applying Theorem~\ref{thm:bu}, not a sign correction to either displayed equation.

\item \textbf{Exterior continuation proof is abbreviated.} Equations (3.38)--(3.46) are printed on pp.~227--228, and the closing text continues on p.~229. They use all-order velocity, time, and pressure bounds, a divergence correction via stationary Stokes, local strong solvability, and linear regularity; the closing harmonic \(L^3\) Liouville step is only described. The manuscript does not follow this route. Theorem~\ref{thm:vorticity-regularity} and Lemma~\ref{lem:vorticity-top-extension} give a bounded, continuous vorticity of class \(W^{2,1}_2\) with zero trace, which is the class of Theorem~\ref{thm:uc} and Theorem~\ref{thm:bu}. The infinite-order vanishing hypothesis of Theorem~\ref{thm:uc} follows from Lemma~\ref{lem:flat-from-vanishing}, and the Liouville step is Lemma~\ref{lem:liouville-L3}. No higher-order regularity is used.

\item \textbf{Initial-time \(L^5\) estimate.} Lemma~7.1, printed pp.~243--244, explicitly cites an external theorem for its general mixed-norm estimate (7.4); Theorem~7.4, printed pp.~246--248, abbreviates the contraction and limit details. These are inputs to the \(L^5\) and uniqueness conclusions of the source's Theorem~1.3, not to its local Theorem~1.4. The manuscript proves those conclusions in Theorem~\ref{thm:ess-l5-unique} without the general estimate (7.4): Lemma~\ref{lem:pv-small-heat} proves the special heat estimate that is needed, and Proposition~\ref{prop:pv-local-solution} carries out the fixed-point construction and its limit (see Remark~\ref{rem:pv-departure}).

\item \textbf{Pointwise boundary trace and dimension conventions.} Theorem~4.1, equations (4.1)--(4.3), printed p.~230, evaluates \(u(x,0)\) although the cylinder is open and does not specify a trace or representative. It states a \(W^{2,1}_2\) hypothesis and a differential inequality but concludes pointwise vanishing. The source uses a spatial \(\mathbb R^3\) cylinder, writes an \(\mathbb R^n\)-valued field, and uses a \(\rho^{n-1}\) factor in the proof; it also concludes on the whole ball although Lemma~4.3 estimates only an inner ball before saying the extension is immediate. The notation \(A_0(R,T)\) in §4 and the small absolute \(A_0\) in §5 are unrelated. Theorem~\ref{thm:uc} fixes the continuous trace convention and dimension three; its proof uses the Gaussian-box estimate in Lemma~\ref{lem:uc-gaussian} and gives the explicit radius iteration that enlarges the inner-ball estimate to the whole ball.

\item \textbf{Gaussian propagation omits closure and radius steps.} Lemma~4.3, printed pp.~230--233, applies Proposition~6.1 to a cutoff field in \(W^{2,1}_2\), while that proposition is stated for smooth compactly supported fields; the density passage and the singular \(\varepsilon\downarrow0\) bound are not supplied. It also moves from an inner-ball estimate to the whole target ball without the radius iteration and cites the local \(L^2\)-to-pointwise estimate without proof. The manuscript supplies the smooth-to-\(W^{2,1}_2\) passage in Lemma~\ref{lem:carleman-sobolev}, proves the Gaussian-box estimate in Lemma~\ref{lem:uc-gaussian}, and carries out the explicit radius iteration in Theorem~\ref{thm:uc}.

\item \textbf{Half-space trace and weak-gradient hypotheses.} Theorem~5.1, equations (5.2)--(5.4), printed p.~234, does not define the zero initial trace needed for extension across \(t=0\); without it, the extension can create a boundary distribution. Condition (5.4) also omits the weak spatial gradient assumed explicitly in Seregin's corresponding class. The source sets the spatial cylinder in \(\mathbb R^n_+\) and the field target in \(\mathbb R^n\), then states that the result also holds for every finite target dimension \(m\). The manuscript and the formalization fix \(n=3\). The source's repeated symbol \(A_0\) denotes unrelated quantities in §§4 and 5. Theorem~\ref{thm:bu} makes the continuous trace and weak-gradient clauses explicit.

\item \textbf{Omitted quantitative choice and noncompact test field.} Lemma~5.2, printed pp.~234--236, existentially chooses \(A_0<1/32\); it does not display a numerical choice small enough for the later bounds, although the stated existential choice permits one. This is an omitted explicit quantitative choice, not a threshold contradiction. Lemma~5.3, printed pp.~236--238, applies Proposition~6.2 to a noncompact Sobolev cutoff without a cutoff/tail limit or density argument and compresses the propagation step from Theorem~4.1. Its general-dimensional notation writes \(\mathbb R^3_+\) where \(\mathbb R^n_+\) is required. The manuscript computes the absorption constants in Proposition~\ref{prop:carleman-halfspace}, extends the estimate by density in Lemma~\ref{lem:carleman-sobolev}, and carries out the needed cutoff and iteration in Theorem~\ref{thm:bu}.

\item \textbf{Final large-growth iteration is compressed.} The proof of Lemma~5.4 and Theorem~5.1, printed p.~238, invokes the small-growth lemma with a parameter \(M>A_0\), outside that lemma's stated hypothesis; the trace and rescaling steps are also skipped. A valid repetition requires translating and rescaling each slab and justifying its trace. The manuscript gives the corresponding explicit iteration in Theorem~\ref{thm:bu}; the source's compressed step is not a counterexample to Theorem~5.1.

\item \textbf{Smooth Carleman estimates applied beyond their stated class.} Propositions~6.1--6.2, printed pp.~238--243, are stated for smooth compactly supported fields, but the preceding applications require Sobolev cutoff fields. Proposition~6.1 asserts an absolute positive constant \(c_0\) but gives no numerical value; Proposition~6.2 compresses the final Cauchy--Schwarz absorption without recording the uniform constants. The manuscript computes explicit constants in Proposition~\ref{prop:carleman-gauss} and Proposition~\ref{prop:carleman-halfspace}, and supplies the density extension in Lemma~\ref{lem:carleman-sobolev}.
\end{enumerate}

\subsection*{Sources of the existence and pressure results}

The corrections to Ożański--Pooley, Tsai, Robinson--Rodrigo--Sadowski and Lemarié--Rieusset, and the corrections to Theorem~1.1 and equations (3.2)--(3.4) of Escauriaza--Seregin--\v{S}ver\'ak, concern the existence and associated-pressure results of Section~\ref{sec:existence-in-ckn}. They are recorded in the appendix of the CKN paper~\cite{CKNLean}.

\subsection*{Seregin lecture notes}

\begin{enumerate}
\item Appendix~A.1, Proposition~1.1 and Proposition~1.2, printed pp.~229--234, give the matching Carleman calculations but also compress the final \(L^2\) absorption/Cauchy--Schwarz step after (A.1.26); in particular the uniform constant is not derived there. The manuscript computes the constants in Proposition~\ref{prop:carleman-gauss} and Proposition~\ref{prop:carleman-halfspace}.

\item Appendix~A.2, Theorem~2.1 and Lemma~2.2, repeats the pointwise Gaussian route and compressed inner-to-whole-ball continuation; Appendix~A.3 repeats the half-space proof's final iteration without fully stating its trace/rescaling details. The manuscript instead supplies an integral-box estimate and explicit radius/large-growth iterations in Lemma~\ref{lem:uc-gaussian}, Theorem~\ref{thm:uc}, and Theorem~\ref{thm:bu}.

\item Appendix~A.3, condition (A.3.4), explicitly includes weak \(\nabla u\in L^2_{\rm loc}\), which ESS's (5.4) omits. The appendix's initial trace convention is still unstated. Its \(Q_1^+\) notation uses \(\mathbb R^3_++e_n\) in an \(n\)-dimensional argument; it should be \(\mathbb R^n_+\). The manuscript fixes spatial and target dimension three and makes the weak-gradient and continuous trace clauses explicit in Theorem~\ref{thm:bu}. The notes also describe the first Carleman cutoff with support outside a larger parabolic cylinder; after the separate cutoff at the initial time, the estimates agree with the published proof. Appendix~A.3 has the ordinary editorial typos “result result” and “fist integral,” which do not change the argument.
\end{enumerate}

\bibliographystyle{amsplain}
\bibliography{refs}
\end{document}
